% ============================================================
%  Harald Høffding, "Spinozas Liv og Lære. Et Bidrag til Tænkningens
%  Historie i det 17de Aarhundrede"
%  "Spinoza's Life and Doctrine. A Contribution to the History of Thought
%  in the Seventeenth Century"
%  København: Rudolph Klein, 1877.  VIII, 166 pp.
%
%  TRANSLATION (English), from transcription.tex (Danish) in this directory,
%  the text of record (collated 2026-10-03; see its header).
%
%  DRAFT COMPLETE (2026-10-03): the whole book, pp. III--166, translated
%  from the skeleton in reading order (TRANSLATION-PLAYBOOK.md, the batch
%  loop); to be read by an independent reviewer.
%
%  BOOK-SPECIFIC CONVENTIONS (beyond TRANSLATION-PLAYBOOK.md §2):
%   * Italics in the Danish (\textit: Latin titles and terms) -> \textit;
%     letterspacing (\emph: emphasis, names in the notes) -> \emph.
%   * The print's misprints and errata: translate the sense, with a comment
%     at the site.  Errata (p. 166): p. 8 "for" l. "fra" (from); p. 104
%     "hallucition" l. "hallucination"; p. 132 "gneerositas" l.
%     "generositas".  Misprints: p. 14 "førte"; p. 81 "Mulighednn"; p. 102
%     "Tingene reale, lovmæsssige"; p. 105 "Ideassosiationerne"; p. 110 "før"
%     (sense "for"); p. 125 "Existens", "Kavsalitetsprincip"; p. 141
%     "Hejbergs"; p. 154 "parodoxe"; p. 160 "hinanden-"; p. 165 "mennelige"
%     (menneskelige).
%   * The sections of the last part are headed by numerals only (I.-V.);
%     their titles come from the print's contents page, which numbers them
%     1.-5.  Headed here as in the text, titled in the Contents as in the print.
%   * Quotations in Latin, French, German, Dutch stay in the original
%     (Høffding himself translates the French verse of p. 56 on p. 166).
%     ɔ: -> "i.e.".
%   * Page markers \opage{N} are placed by hand, inline, at the English
%     word where printed page N begins (checked: 1--166 in order, once each).
%   * Høffding's Danish names for Spinoza's terms (p. 73) are kept in
%     English: Grundvæsen = fundamental being (substance), Grundform =
%     fundamental form (attribute), Enkeltvæsen = individual being (mode);
%     Forestilling = representation (imagination, where it renders
%     Spinoza's imaginatio).  The List of Terms (p. 165) keeps the print's
%     order and gives the Danish equivalents in italics.
% ============================================================
\documentclass[12pt]{book}
\usepackage[utf8]{inputenc}
\usepackage[T1]{fontenc}
\usepackage[english]{babel}
\usepackage[protrusion=true,expansion=true]{microtype}
\usepackage[margin=1.3in]{geometry}
\usepackage{setspace}
\usepackage{amsmath}
\usepackage{enumitem}
\usepackage{libertinus}
\usepackage{libertinust1math}
\usepackage{textalpha}
\usepackage{fancyhdr}
\usepackage[colorlinks=true,linkcolor=black,urlcolor=black,
  pdftitle={Spinoza's Life and Doctrine},
  pdfauthor={Harald Høffding}]{hyperref}
\setlength{\headheight}{15pt}
\pagestyle{fancy}
\fancyhf{}
\fancyhead[LE,RO]{\thepage}
\fancyhead[RE]{\textit{Spinoza's Life and Doctrine}}
\fancyhead[LO]{\textit{\leftmark}}
\setlength{\parindent}{1.5em}
\setstretch{1.3}
\newcommand{\tocline}[3]{\noindent\makebox[2.2em][r]{#1}\hspace{0.6em}#2\dotfill\ #3\par}
\usepackage{marginnote}
\newcommand{\opage}[1]{\marginnote{\footnotesize\textit{[#1]}}}

\title{Spinoza's Life and Doctrine\\[0.4em]
  \large A Contribution to the History of Thought\\ in the Seventeenth Century}
\author{Dr.\ H.\ Høffding}
\date{Translated from the Danish\\[0.6em]
  \small \textit{Spinozas Liv og Lære. Et Bidrag til Tænkningens Historie\\
  i det 17de Aarhundrede.} Copenhagen: Rudolph Klein, 1877}

\begin{document}
\frontmatter
\maketitle

% ===== PREFACE (pp. III--V) =====
\chapter*{Preface}
\markboth{Preface}{Preface}
\opage{III}The writing of this little book is owed to the esteemed publisher's
invitation to include Spinoza in his series of ``cultural-historical
personalities.'' I have not been blind to the difficulties of such a task.
Spinoza's significance for the history of culture is closely bound up with his
place in the history of abstract speculation, and it cannot rightly be understood
without a closer insight into his philosophy. I have tried to give as easily
comprehensible a survey of it as possible, and I have laid out the book so that in
the first three parts the reader makes a preliminary acquaintance with Spinoza's
most important ideas before a connected account of his doctrine is given in the
last part.

We have in our literature a thorough and penetrating monograph on Spinoza by the
late Professor Brøchner. Through the clear light in which that work sets the
various elements of Spinoza's philosophy, it holds a high place in the great
literature on Spinoza. But it is not easily accessible to one who wishes to make
Spinoza's acquaintance without rather advanced philosophical preparation. I have
conceived my little book as a kind of \opage{IV}introduction to Brøchner's, in that I
have renounced completeness and thoroughness in detail, as well as learned
documentation. I have, however, been able to use what was not yet available in
complete form when Brøchner wrote his treatise: the interesting work ``On God and
Man,'' which shows us the first form of Spinoza's philosophy and points clearly to
its historical presuppositions, as well as the letters from and to Spinoza that
have been found since and that shed light on several points in his biography. Van
Vloten has collected everything pertaining to this in his book ``Ad Benedicti de
Spinoza opera qvæ supersunt omnia supplementum'' (Amstelodami 1862).

The standpoint from which I regard and judge Spinoza's philosophy differs from
Brøchner's. My late teacher, of whom I have so often been reminded while working
on this book, stood nearer to Spinoza as a speculative philosopher than I do. His
criticism is therefore directed, not against the very principle of Spinoza's
speculation, but against the inconsistencies in its execution. I had to go
further, not only by virtue of my own philosophical views, but also because I set
myself the task of indicating the significance of Spinoza's theories in the light
of the later history of the sciences, and especially of the research of our own
time. But in spite of this more far-reaching criticism, my admiration for the great
thinker has grown rather than diminished. Spinoza may be said to have celebrated a
double resurrection, after having lain for more than a century ``like a dead
dog'': the first time at the end of the eighteenth century, in the philosophical,
\opage{V}poetic and religious awakening; the second time in our own day, through the
great hypotheses of natural science and the reawakening of philosophy after the
slackness that the surfeit of speculation had brought about. It is characteristic
of his position in the history of thought that two such different epochs could
meet in adherence and sympathy toward him. In many ways his method and his way of
thinking are even directly opposed to those that prevail in the research of our
time; but here we have an example of how---to use an image I have employed
somewhere in \emph{this} book---workers who dig a tunnel, each from his own side,
meet in the middle, provided only that each, from his own starting point, goes
forward along the right path.

\medskip
\noindent\hspace*{3em}June 1877.\par
\hfill\emph{Harald Høffding.}\par

% ===== CONTENTS (p. VII), as printed =====
\chapter*{Contents}
\markboth{Contents}{Contents}
\hfill Page\par
\tocline{}{Introduction}{1}
\tocline{}{Spinoza's Life}{19}
\tocline{}{Spinoza and His Correspondents}{44}
\tocline{}{Spinoza's Doctrine}{72}
\tocline{}{\quad 1. Fundamental Concepts}{72}
\tocline{}{\quad 2. Mind and Matter}{89}
\tocline{}{\quad 3. Theory of Knowledge}{102}
\tocline{}{\quad 4. Psychology and Ethics}{120}
\tocline{}{\quad 5. Political Theory. Philosophy of Religion}{147}
% (as printed: section 3 begins on p. 103, not 102; the part is headed
%  "Spinoza's Philosophy" on p. 72 -- see the transcription.)

\mainmatter
% ===== INTRODUCTION (pp. 1--18) =====
\chapter*{Introduction}
\addcontentsline{toc}{chapter}{Introduction}
\markboth{Introduction}{Introduction}
\opage{1}The great importance of the history of science is being more and more
recognized in our day. Not only does a strong historical interest go hand in hand
with the interest in the scientific problems themselves---a wish to know the
earlier state of these problems and the way in which they were formerly
treated---but it holds of many questions that they are first rightly understood
when one sees how they have historically arisen and been conceived. Often the path
by which a scientific truth has gradually been won is the same one that, on a
shortened scale, must be traversed by everyone who wants to make it his own. But
there is surely no science for which history has such great importance as for
philosophy. This must be ascribed, first, to the fact that philosophy is richer in
endeavors than in results. Given a sum of results, one can to a certain degree
forget when, how and by whom they were won. The Pythagorean theorem can be fully
appropriated even if one has never heard anything of Pythagoras. But where the
chief weight must be laid on the ever renewed striving,---where it is precisely a
matter of \opage{2}sharpening the sting of the problems in order to set thought in
motion anew, there historical orientation has a peculiar importance. The earlier
attempts become partly models, partly warnings for the later ones. It lies,
further, in the nature of philosophical ideas that they often come into the world
too early. They are often ideals, set up and won by a bold speculation that goes
beyond the actual results of real knowledge and anticipates a goal that in reality
can be reached only laboriously and after long ages. Since the final task of
philosophy is to develop a coherent and comprehensive world-view, philosophical
ideas are directed toward the unity and the whole of things; but they presuppose,
besides, that the nature and the interrelations of particular things are known.
The same thing, moreover, happens within the real sciences too, with the views and
hypotheses that express the connection of phenomena and give their explanation.
Huyghens put forward the so-called wave theory in the seventeenth century to
explain the phenomena of light, and only in our century has it won general
recognition. Lamarck's theory of the origin of species became an object of wider
attention and assent only through Darwin's investigations. The significance of
scientific ideas, in philosophy as in the other sciences, depends on how far they
have really been drawn from nature and carried through with consistency and
clarity. If they have roots in reality and inner consistency, their victory will
come, early or late.

Spinoza's philosophy is of great interest in both respects. It is no accident that
the call \opage{3}to commemorate the two-hundredth anniversary of his death by raising
his statue has found support around Europe, not only among philosophers by
profession, but also among eminent representatives of the real sciences. For his
system does not merely have a classical significance through its purity and clear
coherence, so that we can learn from it, as from a model, how a scientific
world-view must be built. That would not be enough to awaken a general interest.
It is precisely the systematic form of a philosophical doctrine that most readily
falls prey to transience. We can move about in Spinozism, as in the other
philosophical systems, much as we do in the buildings of antiquity, which we admire
and from which we can learn much, although we would not ourselves build in the same
way now. But in our day one cannot study philosophy and general natural science
without one's thought being led back to Spinoza in many ways. Many of his
fundamental thoughts may be said to have found their confirmation in the results of
recent science. Out of the perishable form of the system the kernel of the ideas
develops, and it is then of great interest to see how, by what method and through
what motives Spinoza came to the ideas at which we have now arrived, in part by
other paths.

But Spinoza's ideas were more to him than objects of theoretical speculation. They
fused with his personal view of life, and precisely the seriousness with which he
asserted the right and the capacity of scientific inquiry to determine the view of
life became decisive for his fate and brought him into conflict with the
ecclesiastical \opage{4}authority, which here recognizes only the right of tradition.
His personal life and his character thereby acquire a significance in philosophical
respects as well. Philosophy is, to be sure, first of all a theoretical concern; but
it cannot and must not give up the demand that the truth it has recognized become
determinative for life; it is to act as a salt against the insipidity of dogmatism.
It is therefore its pride that it can point in its history to figures like
Spinoza's, and we shall see in what follows how his character and biography have
played no small part in the struggle for the legitimacy of free inquiry.

As already indicated, a philosophical system must be judged from a twofold side.
One must examine its inner coherence and consistency, and its agreement with
reality as this is accessible to us through scientific experience. Neither of
these two points of view can be dispensed with. What thought seeks in philosophy
is a scientific conception of life and the world, in which all the particular given
facts are joined into a whole and brought into inner coherence with one another.
Unity and coherence must therefore be present; not only logical contradictions, but
also what is detached, isolated and merely empirical is unphilosophical. But this
unity and coherence must not be bought by setting aside essential facts, nor by
distorting and arbitrarily reinterpreting them. The principle that lies at the basis
of a philosophical system must therefore be chosen so that, in its further
unfolding, it can contain and illuminate everything that stands firm through
scientific experience. At different \opage{5}times philosophy has laid different
weight on these two sides. The dogmatic and idealistic tendency in philosophy lays
the chief weight on the carrying through of the principle; the critical and
realistic tendency, by contrast, lays the chief weight on agreement with experience
and on the grounding of the principle itself. The last-named tendency seems more
and more to be gaining the upper hand, and it is in any case the prevailing one in
our day.

The history of philosophy presents us with a series of systems, each of which lays
its own principle at the basis. Which principle is seized upon in the particular
system is naturally not arbitrary. On a deeper view\footnote{See H.~Brøchner:
Bidrag til Opfattelsen af Filosofiens historiske Udvikling. Kbhvn.\ 1869.} it will
appear that the leading idea of a philosophical system emerges from the whole
development of the mind, just as it in turn acts back upon that development. The
individual thinker finds his principle, so to speak, in the spiritual atmosphere in
which he has developed, and separates it out of this atmosphere, just as a chemist
separates a substance from its compound with several other substances in order to
consider it by itself. Manifold principles that are essential elements in the
content of the general consciousness have in this way received their closer
treatment and elucidation in different philosophical systems. But philosophy must
necessarily come to the insight that it is of the greatest importance to examine
how far the principles themselves can be grounded, and how it is possible to reach a
true and universally valid principle. English philosophy struck out on this path
from the very beginning, and \opage{6}Kant, the greatest philosopher of modern times,
has with great profundity shown it to be the only sure path to philosophical
knowledge. A philosophy that starts from a principle which is not itself grounded
will always come to stand close to theology; there will be a spiritual kinship,
however much the philosopher in question asserts the independence of rational
knowledge, and however energetically he has emancipated himself from faith in
authority. It is above all in method that this kinship between dogmatic or
speculative philosophy and theology shows itself. Both proceed deductively and seek
to derive their doctrine from one highest principle, which for theology is given by
authority and tradition, while in speculative philosophy it is an idea that presents
itself to reason with immediate clarity and seems to be self-evident. Critical
philosophy is, so to speak, philosophy to the second power. It will not merely
philosophize about a given principle, but wants to understand its origin and its
justification. It demands, moreover, an account not only of the principle of
authority but also of the ``self-evident'' idea. It sets to work inductively and
analytically, insists on getting firm ground under its feet, and recognizes only
those principles that follow with necessity from the analysis of the facts of
experience.

Spinoza belongs decidedly to the dogmatic-speculative tendency in philosophy. His
method is the deductive one. His fundamental concept, the idea of substance, arose
for him partly under the influence of the teaching of the popular religion in which
he had grown up, partly through the influence of the preceding philosophy. It stood
before him with an inner truth, raised above the need for \opage{7}actual proof.
Substance (Dutch: \textit{zelfstandigheid}) is that which is in and of itself, which
presupposes nothing whatever else. The idea of substance lies at the basis of the
higher popular religions' conception of the Deity as the absolute being, the origin
of all things, eternal, imperishable and without limits. But the religious
conception, governed by imagination, cannot carry this idea through with full
consistency; it always involuntarily places finite existence outside the absolute
being, and even ascribes to it the power to hinder, or at least to restrain, its
working. In a dim sense of the contradiction in which it thereby entangles itself,
the religious conception then seizes upon the dogma of creation as the solution of
the riddle: there is indeed something outside the absolute being, with independent
reality, but it has come into being through the absolute being's own will and has
reality only through it. Spinoza takes up the idea of substance in such a way that
he carries it beyond theological dualism (the doctrine of twoness). But the
connection with the religious circle of ideas is seen not only in the fact that in
his chief work (the Ethics) he directly designates substance as God, but especially
in the whole tone and manner of expression of an earlier work of his (``On God, Man
and His Well-Being''), which was found in Amsterdam fifteen years ago. He uses
religious conceptions throughout, but gives them a rationalistic application and
development. How near he still stands here to the circle of ideas of the Old
Testament is seen from an utterance such as this: ``True love always springs from
the knowledge that its object is glorious and good. Can it then burn more fiercely
toward \opage{8}anyone than toward the Lord our God? For he alone is glorious and is
the perfect good.'' Spinoza is a religious thinker. The idea of God stands before
him with inner clarity and validity, and his whole thought moves within it. But
through the further development he gives the idea of God he steps beyond the limits
of positive religion, and thereby draws upon himself the charge of irreligion and
atheism. He saw clearly that the being which is in and of itself must be the only
being and must contain everything within itself, so that there can be no true being
outside it. He thus abolished the dualistic opposition between God and the world and
rejected the dogma of creation. He is thereby the real founder of monism, or the
doctrine of unity. Philosophy, which seeks unity and coherence in everything, must
necessarily be monism, and in so far Hegel was right when he said: ``Spinoza is so
cardinal a point in modern philosophy that one may truly say: either you have
Spinozism, or you have no philosophy.''

Spinoza is naturally not the first to have put forward monism. Within Greek
philosophy he has predecessors in the Eleatics, Plato and Aristotle; but their ideas
have hardly had any influence on him. It seems to be otherwise, judging by what can
be traced in the treatise just mentioned, with the Jewish and Arabic philosophers
of religion of the Middle Ages, % the print: „for Middelalderen“, l. „fra“ (Rettelser): of, from
in whom a strong pantheistic tendency appears. In the eleventh century there lived
in Saragossa a Jewish poet named Ibn-Gebirol, whose religious hymns won great fame.
He was the real regenerator \opage{9}of Hebrew poetry, and he holds the first place
among the Jewish poets of the Middle Ages. But the learned connoisseur of the Jewish
and Arabic literature of the Middle Ages, the Frenchman S.~Munk, has also shown
(see his ``Mélanges de philosophie juive et arabe'') that Ibn-Gebirol is the same as
the author against whom the Christian scholastics polemicized so eagerly under the
distorted name Avicebron, and who may be called the classical pantheist of the
Middle Ages, as Spinoza is that of modern times. In his most important work, ``The
Source of Life'' \textit{(fons vitæ)}, he developed---relying on Aristotle and the
Neoplatonists---the idea of one form and one matter as revealing themselves in
everything, in the spiritual as well as in the corporeal. In the infinite series of
beings there is no leap. There is one substance, which subsists by itself and bears
all differences, which takes on all forms and gives everything its being and its
name. Alongside these pantheistic ideas there runs, to be sure, an attempt to hold
fast to the dogma of creation; but ``The Source of Life'' is nonetheless a purely
philosophical work, the only Jewish work from the Middle Ages in which no passages
of Scripture or traditions are cited. There is, according to S.~Munk, an inner
connection between Ibn-Gebirol's poetry and his philosophy: ``it is the same circle
of thoughts that reveals itself under two different forms, and while one always sees
in his poems a mind sunk in deep musings and a stranger to what interests the
crowd, his philosophical works often bear traces of the poet's lively imagination
and enthusiasm.'' An Arabic philosopher of the following century, Ibn Roschd (the
Averroës of the scholastics, % the print: „Avèrroës“
born in Cordova, died \opage{10}1198), expresses similar ideas, apparently without
knowledge of Ibn-Gebirol's writings. He polemicizes directly against creation out of
nothing. Eternal matter contains all the forms within itself. While Muhammadan,
Jewish and Christian theologians alike abolish the concept of nature, philosophy
asserts a constant transition from possibility to actuality according to the law of
eternal necessity. Man thinks by participation in the absolute reason, which is
eternal, while every individual thinking being is perishable. True knowledge, which
is reached by becoming a participant in the divine reason, is the highest
perfection; once it is reached, there is no longer any question of reward and
punishment.---A somewhat younger contemporary of Ibn Roschd was the Jewish
theologian and philosopher Moses ben Maimon (Maimonides), born in Cordova in 1138,
died in Cairo in 1204. He sought to unite faith and knowledge, interpreting the
Bible in a rationalistic way, conceiving the religious ceremonies as aids to
exercise in true knowledge, and maintaining that scientific insight into God's being
is the necessary presupposition of the love of him. His great and remarkable work
``The Guide of the Perplexed'' (Moreh Nebuchim), which has been translated into
French by S.~Munk, has had great influence on many eminent minds within Judaism; it
was the reading of it that led men like Spinoza, Mendelsohn and Maimon from the
synagogue over to philosophy. It is probably through it that Spinoza came to know
Ibn Roschd's ideas.---According to what Grätz states in his ``Geschichte der Juden''
(vol.~10), there is said to be a certain agreement between some of Spinoza's ideas
and \opage{11}expressions and the so-called Kabbalah, a fantastic religio-philosophical
mysticism that arose within Judaism. Grätz believes in particular that he can
demonstrate a similarity between a kabbalistic work by Abraham de Herrera, published
in 1656, and several propositions in the first book of Spinoza's ``Ethics.''

Since we know from Spinoza's biography that he spent his first years of youth
studying Jewish theology and philosophy, it is of interest to see what stimulating
influences in the direction of his later standpoint he may have received here,
directly or indirectly. But these influences were surely only preliminary; he
himself declared that it was Descartes's writings that had taught him philosophy.
There is a great difference between those medieval speculations and Spinoza's clear
and sober thought, which took mathematics as its model.

From the earlier treatise cited above one also sees that Spinoza was influenced by
the greatest philosopher of the Renaissance, Giordano Bruno. This restless and
fermenting spirit himself points back to Ibn Gebirol and Ibn Roschd, as well as to
the Neoplatonists and Aristotle. But apart from these older ideas, which he sought
to revive, Giordano Bruno's whole world-view is essentially determined and motivated
by Copernicus's discovery of the earth's motion around the sun. He was the
enthusiastic prophet of this discovery throughout Europe, and he became its martyr.
Through this discovery the old opposition between heaven and earth was effaced, and
an infinite horizon opened up, in which Bruno's imagination loved to lose itself.
His enthusiasm for nature \opage{12}can still be traced in some of Spinoza's
utterances. What separates Spinoza from Bruno is not only his strict, systematic
thinking, but above all the great advance made in natural science and philosophy in
the interval between their appearances. Bruno's conception of nature is poetic and
idealistic, and it does not rest on a knowledge of the definite laws of nature and
of the real connection of phenomena. But at about the same time that he suffered a
martyr's death at the stake of the Inquisition, Galileo made his great discoveries
and laid the foundation of modern exact natural science by applying mechanics to
astronomy. The mechanical laws of motion, which appear with particular clarity in the
motions of the heavenly bodies, now became the foundation of the conception and the
explanation of nature. Galileo is the founder of mathematical physics, which rests
on the principle that all physical phenomena can be explained as motions that can be
measured either directly or indirectly. Often it is only with the help of experiment
that one becomes able to measure; but as soon as this has succeeded, the starting
point for mathematical calculation has been won. The goal toward which the
mathematical natural science founded by Galileo strives is to derive the facts given
in experience with mathematical necessity from as few and as simple presuppositions
as possible.

But the path of experimental inquiry is long; for long periods it attains only
scattered results and opens only distant prospects of total views and comprehensive
standpoints. For many philosophers there has therefore been something unsatisfying
about this method. Only \opage{13}English philosophy at once attached itself---with
Bacon, Hobbes and Locke---to the principle and method of experimental inquiry. But
it could do so only by provisionally leaving the great speculative problems in
abeyance. Where there was a lively need to discuss the questions that lie on the
border between religion and philosophy, this method could not yet suffice. English
empiricism (the philosophy of experience) was long able to flourish very well with
the usual theological conceptions as its background. It was otherwise on the
Continent, where philosophy, because of its speculative tendency, at once collided
with theology and had to declare itself for or against it. Experimental inquiry
acquired only a subordinate significance there; the essential interest is directed
toward the development and the carrying through of general ideas. This appears
clearly in Spinoza's nearest predecessor, the Frenchman René Descartes.

Descartes was dissatisfied with the scholastic philosophy he had learned at the
Jesuit college, and sought a new one with a firmer foundation. He then resolved to
disregard all his earlier presuppositions, all the content of knowledge he had
previously taken in, and to doubt everything---not in order to become a sceptic, but
in order to come to know the nature of human knowledge and to see how far reason can
reach by its own powers. Now since doubt is itself a thinking, thinking itself is the
only thing of whose reality there can be no doubt, and everything that I see as
clearly and distinctly as that I think must be true. Descartes therefore sets up the
rule of truth that made so \opage{14}strong an impression on Spinoza: that clear and
distinct conception is the condition of our being able to take anything for true.
The first use he makes of this rule is to show that the very concept of doubt leads
us to the concept of an absolutely perfect being. For, he says, how could I know that
I doubt and that I desire a true knowledge, i.e.\ that I lack something and am not
wholly perfect, if I did not have the idea of a being more perfect than my own, by
comparison with which I can see the defects of my nature? God, or the absolute being,
who contains all perfection within himself, therefore becomes the highest principle
of philosophy, and true knowledge consists in deriving everything from God's being
by clear and distinct inferences. ``Since God,'' says Descartes, ``is alone the true
cause of everything that is or can be, it is evident that we shall follow the best
way of philosophizing if we seek to derive the explanation of created things from
God's being.'' Thus he derives the highest laws of nature from God's immutability,
which is one of his perfections, and from the highest laws of nature, or the laws of
motion, the special laws. Galileo's procedure is too empirical for Descartes; he
reproaches Galileo for not going back to the ultimate grounds. The first of the laws
of nature that he derives from God's immutability is precisely the law of inertia,
which Galileo had demonstrated by experiment: everything remains, so far as it
depends on itself, in the same state, and can never be changed except by external
causes.

Descartes will not, like Galileo, go from effects to causes, but conversely, from
causes to \opage{15}effects, and the highest cause is given to him immediately in the
very nature of thought. His method is therefore thoroughly deductive. The facts of
experience, by contrast, receive only the significance of guiding thought in a
certain definite direction. The model that influences Descartes here is mathematics.
He was already thinking of it when he set up clear and distinct knowledge as the mark
of truth. Mathematical knowledge stands as a great example of transparent clarity,
since in it everything is reduced to relations of magnitude that find their exact
determination. Descartes agrees with Galileo that natural science is to be developed
mathematically; but he believes that he can explain all special phenomena from the
general mathematical laws. He thus makes no distinction between physics and
mathematics. ``My whole physics,'' he says in one of his letters, ``is nothing but
mathematics.'' Everything that did not have mathematical clarity was for him obscure
and confused. He saw the essence of matter in extension. All the other properties of
matter can, in his opinion, be thought away, but not extension in length, breadth and
height; from it, as the essential, all other properties must therefore be derived.
The fundamental properties of matter are divisibility and mobility, and the various
material phenomena all arise through the motion of the parts.

It lay all the nearer for Descartes to make the mathematical the main thing in his
conception of the world, since he was himself a mathematician of the first rank; he
is the first who applied algebra to geometry, and thus the founder of analytic
\opage{16}geometry and mechanics. And in general it must have lain near for the thinker
who wished to set up a new system and a new method in place of medieval scholasticism,
with its subtleties of reasoning and its verbal explanations, to seize upon the
mathematical method, which is the ideal and the model of all inquiry. There lay in
this, moreover, more than a form; there lay a whole program, a prophetic annunciation.
The world no longer stands as a play of obscure forces, but as an infinite series of
phenomena standing in definite relations to one another, for which clear laws hold.
Descartes thus sets up the demand that everything be explained according to necessary
laws, a demand whose progressive fulfilment is the natural science of modern times.
What Galileo showed in a single domain, Descartes made the principle of the whole
conception of the world.

That Descartes makes the transition, with such astonishing speed and ease, from his
starting point, the doubt of everything, to the idea of the absolute being, and makes
this idea the principle of his doctrine, can be explained only by the influence of
religious faith, and we have here a classic example of the peculiarity of the
principles of speculative philosophy that was emphasized above. But what must be
emphasized here in particular, because it is precisely at this point that Spinoza
enters into opposition to Descartes, is that the latter takes up the idea of the
absolute being in the same form in which it appears in positive religion. He was
himself a believing Catholic, and however energetically he holds fast to the validity
and the unrestricted dominion of the laws of nature, he nevertheless traces them back
in the last instance to God's almighty will and \opage{17}explains them by God's
immutability. \emph{If} God willed, he could abolish the laws of nature, and even the
mathematical laws; 2 and 3 make 5 only because God wills it. And at another point
divine intervention becomes positively necessary for Descartes's system. If the
essence of matter consists solely in extension, where then is the cause of material
motion to be sought? Since it does not lie in the essence of matter, which is purely
geometrical, it must lie in God's will. God, says Descartes, set the parts of matter
in a certain motion at the creation.

Since Descartes still holds fast to the dogma of creation, he assumes other
substances besides the absolute one. God has brought forth mind, the thinking
substance, and matter, the extended substance. But here we stand precisely before the
contradiction that philosophizing thought discovers in the religious conception: that
on the one hand we have an absolute substance, and yet on the other hand finite
substances, brought forth by it. Descartes does remark, to be sure, that he does not
take the word substance in the same sense of God and of finite beings; but he does
not show how---once he has used it of God---he can have any right at all to use it of
finite beings. In this inconsistency lies hidden the problem that was the deepest
motive of Spinoza's thought. When Descartes designates both mind and matter as
substances, he thereby comes to set them in opposition not only to the absolute
substance but also to each other. He gets not only a dualism between the infinite and
the finite \opage{18}(God and the world), but also one between mind and matter. With
respect to both questions he is influenced by theological conceptions. He had clearly
seen the difference between the mental (the phenomena of consciousness) and the
material, and the impossibility of deriving the one from the other. But this
difference became for him a complete opposition, so that he and his school in the end
appealed to a miracle to explain the connection of the soul with the body. Here too
Spinoza steps in with clear consistency and carries the problem a mighty step further.

Having now briefly brought out the historical conditions under which Spinoza's
thought developed, we pass on to an account of his life.

% ===== SPINOZA'S LIFE (pp. 19--43) =====
\chapter*{Spinoza's Life}
\addcontentsline{toc}{chapter}{Spinoza's Life}
\markboth{Spinoza's Life}{Spinoza's Life}
\opage{19}It will already have been seen from the foregoing that Spinoza's doctrine stood
in sharp opposition to the religious conceptions that alone prevailed in his time,
and when one now recalls how mighty the political and ecclesiastical dominion of
authority was in the seventeenth century, it seems a marvel that he was allowed to
live as quietly and undisturbed as he did. There was surely no other country than
Holland where this would have been possible. Holland is the cradle of religious
freedom. In the year 1609 the Dutch had ended their heroic struggle for freedom
against mighty Spain,---as a historian has said, the only modern counterpart to the
Greeks' struggles for freedom. As the Greeks saved Western culture, so the Dutch saved
religious and political freedom. In the manifesto in which they declared themselves
free of Spanish rule in 1581, they appealed to ``the law of nature,''---a portent of
the North Americans' Declaration of Independence and of the Declaration of the Rights
of Man in the French Revolution. Freedom in one domain must in the long run bring
freedom in all the others with it.

\opage{20}Protestantism in its first period was not tolerant, and it also had
difficulty in being so, since it was a matter of a struggle of life and death with
the Roman Church. Strict discipline had to be kept within the ranks of the
combatants. William I of Orange seems to have conceived the religious struggle from
a more universal point of view, as a struggle for general religious freedom. But his
successors let themselves be guided above all by their aspiration to found a
monarchy, and they therefore joined the orthodox clergy, whose influence on the
common people thereby also benefited them. Over against the court, the army, the
clergy and the common people stood the republican party, which relied chiefly on the
wealthy burgher class. In religious respects it joined the so-called Arminians, who
contested orthodox Calvinism, asserted freedom of conscience, and allowed everyone to
interpret the Bible as he himself wished. While all the other church communities of
the time hedged themselves about with a substantial apparatus of ``symbolic books,''
anyone could be a member of the Arminians' congregation who was free of idolatry,
offensive conduct and intolerance and acknowledged Scripture as the rule of his life,
however he might interpret it. It must be said to their honor that they upheld
toleration before they were themselves persecuted. For persecution did not fail to
come. The Reformed synod at Dordrecht condemned their doctrine and expelled them from
the church community. Maurice of Orange supported the orthodox, who held to him
because the republicans and the Arminians wished to restrict the influence of the
clergy. Only some time after Maurice's \opage{21}death were the Arminians allowed to
hold divine service. From them again sprang several smaller sects, which followed
freer forms in faith and life.

It was no wonder that those who were persecuted for their faith in other countries
cast their eyes on Holland and chose it as their refuge. While the Dutch were
fighting for freedom, and indeed throughout the whole seventeenth century, the fires
of the Inquisition burned merrily in Spain and Portugal. Persecution fell
particularly on the Jews who had let themselves be baptized under compulsion, the
so-called Marannos, many of whom secretly held fast to their old faith and their old
customs. Already during the Dutch war of independence many Marannos had fled to
Holland. In 1598 the first synagogue was built in Amsterdam, and two more soon
followed. At first they were indeed suspected of being Spanish spies; but once it
had been established that this was groundless, and since they also brought large
capital with them and gave Amsterdam an impetus forward by their commercial spirit,
the authorities showed them great accommodation. The Jewish community grew strongly.
In the year 1639 a college for Jewish language and theology was founded, where
Spinoza pursued his first studies.

Baruch Spinoza was born on the 24th of November, 1632. The chief source for his
biography, the pastor Colerus, who is, however, no reliable witness as regards
Spinoza's childhood and youth, relates that he was born in Amsterdam. But Grätz, the
learned author of the ``Geschichte der Juden,'' has concluded from a passage in one
of Spinoza's letters that he must have been born in Spain and spent his early
childhood there. Spinoza \opage{22}says, namely, that he knew a Jew who was called
Judas the Faithful and who suffered a martyr's death at the stake with great heroism.
Grätz now informs us that this Judas the Faithful was a Spanish nobleman who was led
by zealous study of the Old Testament to the conviction of the truth of Judaism, and
who was therefore burned by the Inquisition in the year 1644. How could Spinoza have
known him, if he had not lived in Spain as a child?---Spinoza's parents are called
Portuguese Jews; but this was the usual name for all Jews from the Pyrenean
peninsula. Either before or, according to Grätz's hypothesis, after the son's birth
they had, like so many other Marannos, taken refuge in Holland. The father seems to
have been a well-to-do merchant. Because of the boy's good gifts he was destined for
the learned profession and attended the Jewish college, where he immersed himself in
the study of the Old Testament and the Talmud. He seems also to have made the
acquaintance of the Jewish philosophers of religion of the Middle Ages and of
kabbalistic mysticism; of one of his teachers, Manasseh ben Israel, it is known that
he embraced the kabbalistic doctrines. His keen understanding soon found
contradictions in the theological doctrine. He often embarrassed his teachers with
his objections and questions. One of his biographers, the physician Lucas, relates
that when the rumor of his doubts spread, two of his fellow pupils came to him and
pressed him with questions concerning various points of the orthodox doctrine, which
Spinoza tried in vain to evade by saying that they had, after all, Moses and the
prophets. Later \opage{23}the same two persons appeared as witnesses against him and
accused him of having mocked Moses and the Jewish people. This caused all the
greater stir in the Jewish community because people had thought to see in the gifted
young man a future pillar of the synagogue and the college. He had been the favorite
of his Talmud teacher Morteira, and only reluctantly did the latter resolve on severe
steps against him, especially since Spinoza was very modest and reserved and led a
quiet and peaceful life. When Spinoza noticed that he was suspected, he withdrew
still further, and his visits to the synagogue became rarer. He now began to
associate with Christians and felt the need of a more universal education than the
biblical-theological one. He learned Latin, which was then, of course, the common
language of learning, and later had as his teacher in the classical languages a
physician named van Ende, a well-known humanist, who kept a kind of school for young
people in his house. Colerus relates that Spinoza fell in love with his daughter, who
helped her father with the teaching, but that she preferred one of his fellow pupils,
who better understood the art of gallantry. But this story, on which Auerbach has
built a novel, can hardly be true, since it has been established that van Ende's
daughter was only twelve years old at the time when Spinoza was excommunicated. Van
Ende was very free in his religious views and taught his pupils, as Colerus says,
other things than Latin.\footnote{Van Ende later settled in Paris, and was hanged
there as a Dutch spy and an accomplice in the Chevalier de Rohan's conspiracy.}
Spinoza's \opage{24}doubts and misgivings found their nourishment here. Soon he gave
up theology altogether and threw himself into scientific and philosophical studies.
``Descartes's writings fell into his hands; he read them with avidity and often
declared afterwards that he had drawn his philosophical knowledge from them. He was
delighted with Descartes's principle that one must accept nothing as true that has not
been proved by good and solid reasons'' (Colerus). In Descartes's doctrine he found a
world-view that was built throughout on thought and reason. Acquaintance with it
completed his emancipation from Jewish theology. His fellow believers harbored the
suspicion that he intended to go over to Christianity, and feared that he would
become a dangerous adversary of the Jewish religion. It must, as Grätz remarks, have
made the rabbis little inclined to tolerance that a multitude of Jews were at the
same time suffering martyrdom in Spain. Spinoza, moreover, was not the first example
of apostasy from the faith of the fathers in the Amsterdam congregation. Some twenty
years before, the synagogue had repeatedly excommunicated Uriel da Costa, a
descendant of Marannos, who was led back to Judaism by reading the Old Testament and
fled with his family to Amsterdam, where he had himself circumcised, but who,
disappointed in his faith in the ideality of Judaism, came out vehemently against the
``Pharisees'' and was led by continued brooding to deistic views. Da Costa's restless
and confused mind led him to no firm standpoint. Despair at being shut out, by the
excommunication, from all intercourse with his kinsmen made him \opage{25}submit to a
hard and humiliating penance in order to be received into the synagogue again. But
soon afterwards he mocked Judaism anew, was excommunicated again, and in the end
killed himself in despair.

With Spinoza, however, they would gladly have avoided an open breach. He was offered
a not inconsiderable yearly pension if he would promise not to come out publicly
against Judaism and to visit the synagogue now and then. When this offer was not
accepted, a fanatic made an attempt to assassinate Spinoza; but he escaped with a few
tears in his cloak. At last (1656) they proceeded to the solemn excommunication,
which shut him out from intercourse with all his kinsmen. Lucas remarks: ``The
excommunication has such weight among the Jews that the excommunicated man's best
friend does not dare to do him any service or even to speak to him, lest he come to
suffer the same punishment. Therefore those who fear solitude would rather suffer any
other punishment whatever than the curse.'' Spinoza had early learned to value
solitude, and through his studies and his whole spiritual development he had,
moreover, been initiated into a more universal world and had gained a larger
spiritual fatherland. When he learned that they meant to expel him, he is said to
have said: ``They are not forcing me here to anything I would not have done of my own
accord. Since they will now have it so, I take the path that opens before me, and I
can console myself that I go away more innocent than the ancient Hebrews went out of
Egypt, although my livelihood is no more assured than theirs.'' But the rabbis were
not content with the excommunication. They \opage{26}informed the authorities that
Spinoza had blasphemed God and Moses, and demanded that he be removed from Amsterdam.
Perhaps they feared that Spinoza would draw others after him; the formula of
excommunication itself mentions that he has led some astray, and there are also, as
we shall soon hear, other traces of his having had disciples already at this time.
When the Reformed clergy joined their Jewish colleagues, Spinoza was banished from the
city for a time ``for the maintenance of order and subordination.'' A Spanish apology
that he is said to have written on this occasion has not been preserved; presumably
he set out in it the views on religious freedom that he later developed at length in
his ``Theologico-Political Treatise.'' He took up residence with a friend in the
country outside Amsterdam. It now stood him in good stead that he had in his time,
following the precept the Talmud gives to all the learned, learned a craft, namely
the grinding of optical lenses. By this work he earned what he needed for his living.

The oldest work we have from Spinoza's hand must have been written shortly after his
expulsion from the synagogue. It is the treatise ``On God, Man and His Well-Being,''
originally written in Latin, which was found fifteen years ago in a Dutch translation.
It is of great interest in showing us Spinoza's view in the making. He has corrected
Cartesianism with the recognition that, when the concept of substance is conceived
consistently, there can be only one substance, whose forms of manifestation all things
are. His peculiar deterministic conception likewise appears clearly and definitely;
everything, the human \opage{27}will included, follows the law of necessity. But he
still holds the Cartesian conception of the relation between soul and body, which he
lets act upon each other. Finally, he uses religious conceptions and expressions
throughout, and assumes final causes in nature. This last is of particular interest,
since it points decidedly to other predecessors than Descartes,---to the thinkers of
the Middle Ages and the Renaissance.\footnote{See above, pp.~8--11.} Neither Descartes
nor Spinoza's later doctrine ascribes validity to final causes in the conception of
nature; both maintain the strict mechanical mode of explanation at every point.

In the preface to the Dutch translation of this treatise it is said that it had been
written down in its time by Spinoza for the use of his disciples, who wished to pursue
the study of ethics and philosophy.---The treatise itself ends with these words: ``I
ask the friends for whom I have written this not to wonder at what is new in it, since
you well know that a thing does not cease to be true because it is not accepted by
many. And since the character of the age in which we live is not unknown to you, I beg
you earnestly to be very careful about acquainting others with these things. I will
not say that you are to keep them entirely to yourselves, but only that, if you ever
begin to communicate them to others, no other aim may move you than the good of your
fellow men.''

In the year 1661 Spinoza (who after his expulsion from the synagogue had exchanged his
Hebrew name \opage{28}Baruch for the Latin Benedictus, which means the same) moved to
Rhynsburg, a small town near Leyden. Here, during the Arminian persecution, a small
free congregation had formed, which would have nothing to do with priests or creeds,
but required only, quite in general, a faith in God and Christ; they were called
Collegiants or Rhynsburgers, and they still existed at the beginning of this century.
Spinoza seems to have been on a friendly footing with them, as well as with the
so-called Mennonites (Anabaptists). It was in the Collegiants' former orphanage that
several hitherto unknown letters from and to Spinoza were found some years ago.
Spinoza never went over to Christianity, but he seems to have sympathized above all
with the freer, half rationalistic forms in which he found it among the Arminians and
Collegiants.

In his rural abode he continued his studies with zeal, and one sees from his letters
that his system was finished in all its characteristic fundamental features at the
beginning of his stay in Rhynsburg. He remained in touch with the circle of younger
men that had already formed around him in Amsterdam, and that made up a kind of
philosophical congregation, striving to live their way into the new view of life.
Strangely enough, the first editor of Spinoza's letters struck out of them every
allusion to this circle of adherents; no doubt because it was not advisable then,
immediately after Spinoza's death, to draw attention to it. From the manuscripts found
in the Collegiants' orphanage, on the other hand, we receive information about it. One
of Spinoza's younger friends, Simon \opage{29}de Vries, who on his advice was studying
chemistry and medicine in Amsterdam, writes to him to get enlightenment on some
difficult points in a manuscript Spinoza had entrusted to him, which was probably the
draft of the ``Ethics.'' But while in the original edition of the letters he makes
this inquiry in his own name, in the original manuscript that was found he speaks on
behalf of several. ``I have,'' he says, ``long wished to visit you, but my time and the
hard winter have not allowed it. I lament my fate that we are now separated by so
great a distance. How happy is your housemate, who can talk with you of the highest
things at meals and on walks. But although we are physically far from each other, you
are nevertheless always present to my mind, especially when I am busy with your
writings. Now since not everything in them is clear enough to my companions---for
which reason we have begun our college again,---and also so that you will not think I
have forgotten you, I write you this letter. Our college is arranged in the following
way. One of us reads aloud (each in his turn), explains everything as he has
understood it, and proves it following the order of your propositions. Now when we
cannot agree on the understanding of one point or another, we note it down and write
to you about it, so that you may explain it to us, and so that under your guidance we
may defend the truth against the superstitiously pious and Christian, and stand firm
against the whole world.''

The first work of Spinoza's to be printed was a systematic exposition of Descartes's
philosophy, which \opage{30}he had given to a young man who wanted to be taught by
him, and whom he did not yet wish to initiate into his own doctrine. It was published
in 1663 by his friend, the physician Ludvig Meyer, under the title \textit{Renati des
Cartes Principia philosophiæ more geometrico demonstrata}. Meyer remarks in the
preface that the work must not be regarded as an expression of the author's own views,
and names, as points in which these differ from Descartes's, the questions of the
freedom of the will, of the human soul as a substance of its own, and of the limits of
knowledge. As to the last point, Descartes had declared, in favor of theology, that
there was much that surpassed the human power of comprehension. Spinoza, on the
contrary, believed that the problems at which Descartes had stopped could indeed be
solved if one only applied the right procedure. In an appendix, \textit{Cogitata
metaphysica}, he indicates somewhat more closely the transition from Descartes's views
to his own.

Besides the circle of younger men in Amsterdam with whom Spinoza remained constantly
in touch, he had many correspondents and received many visits. His time was heavily
occupied, since besides the lens grinding, which was his source of livelihood, he
continued his studies and busied himself with chemical and physical experiments. It
was perhaps in order to be less disturbed that in the year 1664 he left
Rhynsburg\footnote{The little house in which he lived in Rhynsburg still stands, and
the street is sometimes called Spinoza Street. On the house there is the following
inscription from the year 1667: Ach, waren alle menschen wijs, En wilden darby wel: De
aard waar haar een Paradijs, Nu is ze meest een Hel. The first three lines could agree
with Spinoza's thoughts; the fourth is wholly at odds with his optimistic conception.}
and settled in \opage{31}Voorburg, a mile from The Hague. But since he had gained many
friends in The Hague, he moved to that city in the year 1669 at their urging, and lived
there for the rest of his life. First he lodged with a widow, later with the painter
van der Spijk. Most of his time he spent in his room at his work. He diverted himself
by catching flies and putting them into spiders' webs, or by setting spiders to fight,
a sight at which he could sometimes burst into peals of laughter. He also often
examined insects, or particular parts of them, under a microscope. He was a skilled
draftsman, and a whole collection of portraits he had made was found after his death;
among them one of himself in the dress of Masaniello.\footnote{The portrait from which
the vignette of this volume is taken was painted by Spinoza's landlord van der Spijk.
There is great dissimilarity between the various portraits of Spinoza that exist.} He
often went down to his landlord's family and talked with them about everything under
the sun, especially when there was sickness or sorrow in the house. His way of life was
frugal in the highest degree, and his wants were few. The pastor Colerus cites several
details in this respect in his biography, from the bills that were found after his
death. Neither pleasures nor money played any part for him. He showed great
unselfishness toward his sisters, although they had tried to exclude him from his
paternal inheritance. After he had had his right of inheritance recognized by a
lawsuit, he left them \opage{32}almost the whole inheritance. His friend Simon de Vries wanted to give him a sum
of money so that he could live more comfortably; but he declined it, just as he also
opposed de Vries's making him his universal heir; a pension that de Vries had settled
on him at his death he accepted only after having reduced it considerably.

In the year 1670 he published his \textit{Tractatus theologo-politicus}, in which he
examined the relation of religion to philosophy and showed the necessity of religious
freedom. This work is the first proclamation of unconditional freedom of belief.
During the struggles between Catholics and Protestants in the sixteenth and
seventeenth centuries the vanquished had no doubt often pleaded the cause of
toleration, but always with great restrictions. Spinoza asserts the right to express
one's conviction as a human right of which the state cannot deprive the individual,
so long as he does not encroach upon the rights of others. The spiritual nature of
men is so various that they cannot all feel satisfied with the same opinions. Spinoza
shows, further, that there is a religiousness independent of dogmas and ceremonies;
he applies a historical criticism to the biblical books, by which he became the
forerunner of the whole of modern biblical criticism; and he seeks to give a
psychological and historical explanation of the spiritual state of the prophets.
Since he showed, moreover, how dangerous it is when the clergy have too great an
influence in the state, it is no wonder that a great stir arose over the anonymous
book. By bold utterances that made God's power coincide with nature's, he also gave
offense to \opage{33}those who had no love for the hierarchy. Several of his
correspondents grew cooler from this time. The book was banned, but soon appeared
again under other titles.

Spinoza's position was not without danger, although he sought to live as
inconspicuously as possible. He had in a way challenged by his book the mighty party
of Orange and of the clergy, which had the common people on its side. He was,
moreover, a friend of the famous statesman Jan de Witt, who stood at the head of the
republican party, and against whom the preachers stirred up the people in their
sermons. The disasters of the French war brought de Witt down, and the mob, which
regarded him as a traitor, tore him and his brother to pieces in the most horrible
manner. At about the same time the great Condé, who was in command at Utrecht, wished
to speak with the famous philosopher and sent him a safe-conduct. At the wish of his
friends Spinoza accepted the invitation; but when he arrived, Condé had suddenly been
called away. The Marshal of Luxembourg, however, received him with great attention. It
was suggested to him that he dedicate a work to Louis XIV, who would then give him a
yearly pension; but Spinoza would not hear of it. When he came back to The Hague,
people crowded together outside his lodging in a threatening manner, because the
rumor went about that he was in treasonable communication with the enemy. His
landlord was frightened at this; but Spinoza calmed him and declared that ``as soon
as the crowd made a move to attack the house, he would go out to them, even if he
should suffer \opage{34}the same fate as the poor brothers de Witt; he knew in his own
heart that he was a good republican, who had never had anything in view but the honor
and the advantage of the state.''

Not only in France, but also in Germany, attention had been drawn to Spinoza. In the
year 1673 he received from the Elector Palatine, Karl Ludwig, known for his interest
in learning, through Professor Fabritius, an invitation to take up a professorship of
philosophy at Heidelberg. He was promised the greatest freedom to philosophize, so
long as he did not offend the religion established in the state. But Spinoza did not
accept the invitation. ``In the first place,'' he wrote in his reply, ``I should not
be able to carry philosophy further if my time were taken up with teaching. In the
second place, I do not know where the limits of the freedom to philosophize that you
mention lie, so that I might not appear to wish to offend the established religion.
Religious quarrels arise not so much from burning religious zeal as from the various
passions of men and their love of contradiction, which leads them to disparage and
condemn everything, even what is right. And since I have already experienced this in
the private and solitary life I now lead, I must fear it all the more if I accept such
a dignity.'' Spinoza had good reason to be wary of appearing as a public teacher of
philosophy; he would hardly have been able to put forward one or two of his doctrines
before the state religion would have been regarded as offended. At that time, in spite
of all fine promises, there could be no question of \opage{35}real freedom of
teaching.\footnote{Even in our own day it has been declared among us that a university
teacher may not attack the doctrines of the state church. One of our newspapers held
that the late Professor Brøchner had misunderstood his position in this respect. But a
public teacher of a science surely cannot reasonably be bound by other laws than those
that lie in the very nature of the science; they will surely contain every needful
guarantee. It must be precisely his highest duty to separate out of the kingdom of
truth everything that has unwarrantably crept into it. The life of the state and of
the people would in the end languish if certain doctrines were given the privilege of
being exceptions to this rule.}---It seems as though the Elector knew Spinoza only as
the author of the work on Descartes's philosophy, not as the author of the
Theologico-Political Treatise. Even so, it was a sign of liberality to wish to appoint
an excommunicated Jew who had not gone over to Christianity.

Spinoza himself did not regard his relation to the popular religion as a hostile one.
True religiousness did not, in his conviction, depend on any confession of faith
whatever. Those who could not reach the knowledge of the divine by way of thought must
find their rest in the doctrines of a positive religion. He liked to talk with his
landlord and his family when they had been to church, and to ask them what profit they
had had from it. In Voorburg he joined in signing a petition concerning the filling of
a pastorate, and in The Hague he several times heard his landlord's pastor, whom he
praised highly. When his landlady asked him whether he thought she could be saved in
the religion to which she belonged, he answered: ``Your religion is good, and you
should not seek any \opage{36}other or doubt that it will lead you to salvation, if
only you apply yourself to piety and lead a peaceful and quiet life.'' He saw far too
clearly how much it takes to form a view of life for oneself, once one has left the
inherited one, to lead away from it anyone who could not yet stand on his own feet. But
he naturally also saw how much the inherited doctrines can hinder progress in the
knowledge of truth, since human sloth is glad to rest content with what has been
instilled in it as a child. The greatest misfortune they cause, however, is by the
strife and contention they occasion; it is the various opinions received through
tradition that divide men, whereas by following nature they would reach agreement, as
children of one and the same mother.

Just as Colerus emphasizes his respect for the faith of others, another theologian
too, Sebastian Kortholt, who some twenty years after Spinoza's death collected
information about him, and who otherwise condemns him in very fanatical terms, says
that ``no oath and no impudent speech about God was ever heard from Spinoza's mouth.''
On the other hand, the Reformed theologian Limborch relates that he---quite
unexpectedly---met Spinoza at a dinner, and that during the grace at table the latter
showed his ungodly disposition, ``in that by his expression he seemed to mock the
stupidity of those praying.'' % the print: „syntes at at spotte“
This account is wholly at odds with what is otherwise known of Spinoza's character and
conduct, and it is surely only a testimony to how disposed Limborch was to find
``signs of ungodliness'' \textit{(signa animi irreligiosi)} in the free\opage{37}thinker
in whose company he had---very much against his will---found himself. The good
theologian, moreover, as has rightly been remarked, cannot have been very absorbed in
his prayer if he sat watching Spinoza's face.

Spinoza's conversation is said to have been witty and lively, and it held a great
attraction for many. His great acuteness and fine irony made what he said both
instructive and entertaining. His bearing always showed the stamp of mildness and
self-control; he was never seen either much saddened or overjoyed, and when feeling
became too strong he would rather withdraw than give it vent before others. In the
last part of his life his health was greatly weakened; he suffered from a wasting
disease of the chest, which his strenuous studies had nourished and perhaps brought
on. In this sickliness one may perhaps seek the reason for his small productivity
toward the end of his life. The ``Ethics'' seems to have been finished already in the
sixties, and in those years he was also working on the ``Treatise on the Emendation of
the Intellect,'' which was found unfinished after his death. In the year 1675 he gave
in to the requests of his friends and traveled to Amsterdam to have the Ethics printed.
But the rumor at once spread that he meant to publish a book that would prove there
was no God. The theologians complained of him at court and to the authorities, and the
Cartesians, who were regarded as his kindred spirits, joined in, in order to clear
themselves of all complicity. Spinoza therefore decided to postpone the publication to
better times; but he did not live to see them. Only after \opage{38}his death did
Ludvig Meyer publish the Ethics, an unfinished political treatise, the unfinished
Treatise on the Emendation of the Intellect, an unfinished Hebrew grammar, and a
collection of letters.

The ``Treatise on the Emendation of the Intellect'' is of interest, among other
things, because we see from it how Spinoza's thought springs in its last ground from
an ethical motive. He begins this work in the following way: ``After experience had
taught me that everything we meet with in ordinary life is vain and futile, and since
I saw that there is nothing good or bad in the things we most fear, except in so far
as the mind is moved by them, I resolved to inquire whether there is anything that is
a true good, that can communicate its being to us, and that can affect the soul by
itself, even if all other goods are rejected; in other words, whether there is
anything by whose acquisition and possession I could enjoy a lasting, eternal and
supreme happiness.'' In contrast to the false and perishable goods, the knowledge of
the eternal and of our unity with it, or the knowledge of our mind's connection with
the whole of nature, is then set up as the true and highest good. To reach this goal
we must gain a true knowledge of the nature of things and strive that as many as
possible share it with us. A purification and reformation of knowledge must therefore
be undertaken, and a community formed in which as many as possible can be brought up
and developed toward the goal indicated.

In his chief work, the ``Ethics'' \textit{(Ethica ordine geometrico demonstrata)},
Spinoza carried through the mathematical method he had already applied in his
exposition \opage{39}of Descartes's doctrine, a method that can be used only in a very
forced way in treating metaphysical, psychological and ethical questions. The way in
which Spinoza applies this method, however, not only testifies to his great acuteness,
but also brings clarity and transparency into the course of thought, something in
which other speculative methods are otherwise sorely lacking. Spinoza gives,
moreover, in the prefaces and appendices to the individual books and in the remarks
(scholia) to the individual propositions, connected developments that in part touch
on what is most interesting in the book, and to which one should chiefly keep at a
first reading.\footnote{Besides in Latin, one can now also read Spinoza in two German
translations (by Auerbach and by Kirchmann) and in an excellent French translation (by
Saisset).} The ``Ethics'' is divided into 5 books. The first (de deo) treats of God as
the absolute substance, of his attributes and of the relation of things to him, and
thus gives the fundamental features of Spinoza's general world-view. The second book
\textit{(de natura et origine mentis)} treats of the nature of man, his relation to
God and to other finite beings, in particular of the relation between mind and matter
and of human knowledge. The third book \textit{(de origine et natura affectuum)}
treats of the nature and origin of the passions and gives a masterly psychological
portrayal of them. In the fourth book \textit{(de servitute humana s.\ de affectuum
viribus)} it is shown what dominion the passions can gain over man. It is the
portrayal of man in his state of bondage, which gradually passes over into freedom,
\opage{40}as his true nature develops. Finally, the fifth book \textit{(de potentia
intellectus s.\ de libertate humana)} portrays man in the eternal reality of his ideal
being. Thus this magnificent work gives us, in Herbart's apt comparison, the picture of
a drama in five acts. It has exposition, complication and the untying of the knot.
First come God and the human mind; then follows what weighs the mind down, the
affects, and the portrayal of human bondage; finally spiritual freedom.

To the ``Ethics'' is joined the ``Political Treatise,'' on which Spinoza was working
when he died. Two things, after all, were presupposed for the realization of the
highest good: true knowledge, which the ``Ethics'' shows us in its application to
man's relation to the passions, and the right community, whose nature is set forth in
the ``Political Treatise,'' which examines natural right, the power of the state, and
the origin and end of society.

But it is not only by his profound and liberal writings that Spinoza has worked for
the cause of philosophy and of free inquiry. His whole life, whose purity and ideal
direction even his opponents had to acknowledge, has been an important argument in
the dispute between philosophy and theology. It lay in the spirit of those times to
regard it as an impossibility that a man who placed himself outside positive religion
could be a good man and lead a moral life. Spinoza's life therefore became a
theological problem. The first who, after Spinoza himself, declared with energy that
the good and the right have their sanction in themselves and do not depend on the
commands of authority was Pierre Bayle. He pointed in particular precisely to
\opage{41}Spinoza (whose true significance he otherwise failed to see) as proof of the
proposition that morality does not depend on positive religion. In our literature
Holberg and Eilskov were influenced by Bayle's ideas; they emphasize, however, above
all another proposition of his, namely that it is natural disposition and temperament
rather than articles of faith that determine men's conduct. Eilskov holds that when
freethinkers do not live ungodly lives, it must be explained by the fact ``that they
have followed the nature of their disposition more than the laws of virtue; they have
lived so more because they found in themselves no natural impulse to live otherwise
than from fear that they would otherwise transgress their duties. That is without
doubt the reason for Spinoza's quiet and honorable way of life.'' By this conception,
as Holberg says, ``the whole disputation became futile'' (Ep.~210), that is, the
question was disposed of by the recognition that character and conduct are relatively
independent of theories and views, in believers and unbelievers alike. It would also
be wrong to think of Spinoza as a saint; in him, as in everyone else, the moral
character was surely determined by his original temperament, and not without reason
has his disease of the chest also been pointed to as an explanation of his reserved,
resigned way of life. Perhaps one may also reproach him with a certain passivity; his
enthusiasm was in any case of the quiet kind, resting in itself, and it did not lead
him---like a Giordano Bruno or a J.~G. Fichte---to offer open defiance to his age in
the struggle for his ideas. But what is great and beautiful in his character is, on
the other \opage{42}hand, precisely the harmony between his personality and his
doctrine; for in the faith in the eternal ground of nature that stirs in everything and
in all lies also the assurance that the true is eternally victorious. And his life
always keeps the significance that by its testimony it has helped to set irrelevant
moral considerations aside in the discussion of philosophical and scientific questions.

If in those times it was a question whether a freethinker could live a moral life, it
was subject to still greater doubt whether he could die in peace. Many different rumors
were therefore in circulation concerning Spinoza's death. Some claimed to know that he
had stupefied himself with mandrake juice when he felt death approaching, others that
he had uttered cries of despair and forbidden his landlady to let any pastor near him,
so that no one should witness his weakness. Fortunately two theologians undertook,
shortly after Spinoza's death, to inquire into the truth of these rumors, and both of
these ``unwilling'' witnesses declare them wholly groundless. Spinoza's landlord and
landlady % the print: „Vært og og Værtinde“
were still living and could relate everything that was known of his death. His
weakness of the chest had steadily increased; but neither he himself nor those around
him expected that he would die so soon. The day before, he had been downstairs with
his landlord and smoked a pipe while talking with him about the day's sermon, and even
on the morning of the day he died (the 21st of February, 1677) he was likewise
downstairs with his landlord and his family. All the more astonished were they, when
they came home from church in the afternoon, to hear \opage{43}that he had died at
three o'clock. No one had been with him except his friend, the physician Ludvig Meyer
of Amsterdam, who left again the same evening. One of our two authorities, Seb.\
Kortholt, says expressly that he ``breathed his last sigh in peace,'' although he adds
that the learned have disagreed as to whether such a death befitted an atheist. And
Colerus refutes all the rumors, although shortly afterwards he is scandalized that
Spinoza is called ``the \emph{late blessed} Mr.\ Spinoza'' in bills that came in to the
estate. But although Colerus regarded Spinoza as damned in the other world, the
history of philosophy must nevertheless be very grateful to the worthy pastor for
having striven so eagerly in his biography, which in all its naïveté is an interesting
little book, to clear his memory in this present world. It is also a reconciling
feature that Spinoza, the advocate of resignation and toleration, found a
truth-loving and impartial biographer in the camp of his opponents.

% ===== SPINOZA AND HIS CORRESPONDENTS (pp. 44--71) =====
\chapter*{Spinoza and His Correspondents}
\addcontentsline{toc}{chapter}{Spinoza and His Correspondents}
\markboth{Spinoza and His Correspondents}{Spinoza and His Correspondents}
\opage{44}Before we pass on to a connected account of Spinoza's philosophical doctrine, it
will be of interest to dwell on his correspondence. His ideas appear in it in part in
a somewhat freer form than in his writings, and thereby become more easily
accessible. It also serves to characterize his historical position to see how his
nearest contemporaries received his ideas. An account of the relation between Spinoza
and his correspondents thus forms a fitting transition between the biography and the
system.

In the Introduction we brought out the opposition between empirical inquiry into
detail and speculative, deductive thought. This opposition comes out in an
interesting way in the letters exchanged between him and his friend Henry Oldenburg.

Oldenburg was born in Bremen and was, under Cromwell and Charles II, the envoy of the
Lower Saxon Circle in London. He took a great interest in physical and chemical
experiments, was the first secretary of the Royal Society, founded in 1662, and a good
friend of the famous Boyle, the founder of analytical chemistry.

\opage{45}Oldenburg was no great mind, but he was enthusiastic about his acquaintance
with Spinoza, whom he visited in Rhynsburg, and about the vigorous upswing that
experimental inquiry was taking in England at that time. He took pleasure in being the
go-between for his two brilliant friends, but he is not free of adorning himself with
his relation to them. Through him a discussion was carried on between them, prompted
by a work of Boyle's that had been sent to Spinoza. Spinoza objected to Boyle's wishing
to conclude from his experiments that the difference between the properties of bodies
could be traced back to differences in motion, figure and other mechanical relations.
Spinoza holds, like Descartes, that everything in nature must be explained according
to the laws of mechanism, but denies that the justification of this mode of
explanation can be established by experiment. The proof that can be in question here
he regards as given by Descartes and Bacon. If all phenomena in nature presuppose
motion, and motion in turn extension in space, then the essence of matter is one with
extension, and the laws of extension, the propositions of geometry, thus become the
real laws of nature, to which everything must be traced back. Spinoza lacked an eye
for the significance of inductive inquiry. What he sought was the essence of things,
expressed in clear concepts. But experience does not teach us, as he says in another
of his letters, the necessary essence of things, but only what actually exists. What
truly is, for Spinoza, was not the particular things and phenomena, but the general
laws and the infinite being that lies at the basis of everything. Experience can have
only the significance of guiding our thoughts in \opage{46}certain definite directions,
but it cannot give us a true knowledge. Boyle, on the other hand, aimed precisely at
the investigation of the particular and special, and left the inner essence of things
alone. He collected observations and analyzed the particular substances, but he did not
aspire to set up any theory. It was also precisely careful observation that chemistry
needed in order to get beyond alchemy, and it owes much to Boyle's sceptical-empirical
spirit. Even a layman can still read with interest and profit his ``\textit{Sceptical
Chymist}'' and the other treatises in which he combats peripatetic and alchemical
chemistry. Oldenburg urges his friends each to use his own method with zeal. ``You,''
he writes to Spinoza, ``I ask in particular to go on laying the principles of things
with your mathematical acuteness, just as I constantly urge my noble friend Boyle to
confirm and illuminate science by repeated and exact experiments and observations.''
This was in reality the best solution of such a dispute. Neither of the two methods
can be dispensed with; what matters is that each find its proper domain and its true
limits.

Oldenburg was not so fortunate in another discussion with Spinoza. It appeared here,
as was touched on in the Introduction, that the purely empirical urge to inquire can be
very limited; in the particular domain one has chosen, one seeks clear grounds and
proofs, but one can for all that follow, in one's view of life, the first tradition
that comes along and take reason captive. Oldenburg wrote his first letter to Spinoza
shortly after his visit to Rhynsburg. He had felt \opage{47}so drawn to Spinoza that he
at once felt the need to enter into communication with him, especially since much of
what they had talked about was not yet clear to him. What he particularly balked at was
Spinoza's conception of God as the infinite substance, whose forms of expression all
things are. Since Oldenburg always starts from the usual theological conceptions, it is
no wonder that he does not understand Spinoza. It seems as though the latter sees this
and is glad to seize on the chemical exchange he entered into with Boyle through
Oldenburg, in order to end the fruitless discussion. Yet at this stage Oldenburg
appears as the eager and enthusiastic friend of free inquiry. He urges Spinoza to
publish his writings and to speak his mind in spite of the theologians' yapping
\textit{(qvicqvid theologastri ogganire poterint)}. ``Your republic,'' he writes,
``enjoys, after all, a very great freedom. One should therefore also philosophize in it
with freedom. Put away all fear of stirring up the petty men of our time! Long enough
have ignorance and quackery had free play: let us now hoist the flag of true science
and search out the secrets of nature still more deeply than before!'' This zeal,
however, was somewhat cooled by closer acquaintance with Spinoza's doctrine, and above
all by the ``Theologico-Political Treatise'' and the great offense it gave. Both
Oldenburg's and Boyle's views were orthodox. Another of Spinoza's correspondents, the
physicist and philosopher Tschirnhausen, who was staying in London at this time, found
that Boyle and Oldenburg had formed ``strange notions'' of Spinoza's person; he
succeeded, however, in ridding them of these prejudices, so that they again came to
think favorably of him. After an interruption of some ten \opage{48}years the
correspondence is resumed, and Oldenburg begins by declaring that he had held an
unfavorable opinion of Spinoza because he had applied to his views the measure ``that
the usual theologians and the inherited confessions of faith supply''; now, on the
contrary, he has seen that it was a hasty judgment. He wishes to be initiated more
closely into Spinoza's doctrine, and offers his help in gaining it entry among
sensible and good men and in defending it against prejudice. But when he hears that
Spinoza means to publish his Ethics, he earnestly admonishes him not to say anything
that could in any way seem to shake piety and virtue---``especially since this
degenerate and wicked age seeks nothing more eagerly than doctrines that can defend the
prevailing vices.'' Spinoza asks for a closer explanation of what doctrines these are
that could offend piety and virtue; for what agrees with reason he also holds to be
most serviceable to virtue. Oldenburg then names three points. Spinoza seems to confuse
God and nature. He abolishes the truth and validity of miracles, and it is on them,
after all, that the certainty of revelation rests. Finally, his opinion of the person
of Christ is unclear. If only he can satisfy all true Christians on these points, his
cause will be out of danger. Spinoza answers roughly as follows. 1) I conceive God as
an indwelling (immanent) cause; he does not act beyond the limits of his being, for
his being is infinite and has nothing outside itself. The apostle Paul, after all, also
says that in God we live and move and have our being. And when I am accused of
confusing God and \opage{49}nature, nature is taken to mean the outward mass, corporeal
matter, and I am wholly misunderstood. 2) To found revelation on miracles is to found
it on ignorance. No one knows what the powers of nature can accomplish, and because we
cannot find the natural causes in a particular case, we are not entitled to deny their
existence. Religion ought to be judged by the truth of its teaching; otherwise it does
not differ from superstition, which builds on ignorance. 3) It is not necessary to know
Christ after the flesh; what matters is to partake of the eternal Son of God, that is,
the eternal, divine wisdom that has revealed itself in all things, most in the human
mind, and most of all in Christ. But the teaching of some churches that God has taken
on human nature I do not understand; to me it is as absurd as saying that the circle
has taken on the nature of the square.---On this last point Spinoza is a forerunner of
the modern speculative philosophy of religion, and he has in particular expressed a
thought that Strauss developed further in the concluding treatise of his famous
``Leben Jesu.''---Oldenburg now at last comes out with the real chief objection: that
Spinoza teaches an eternal necessity in all things, from which it must follow that no
one is responsible for his actions. Spinoza answers that the necessary connection in all
things is precisely the foundation of his whole doctrine. But, he says, I do not teach
fatalism; the necessity that is the expression of a being's own nature is one with true
freedom; it is not an outward but an inward necessity, and it does not abolish the
difference between good and evil.

\opage{50}It was Spinoza's distinction and his greatness that he did not want thought
only up to a certain point, that he did not stop halfway. His striving for a coherent
scientific world-view was fully justified, even if he cherished illusions about the
power of speculative thought. A philosophical problem, once it has been posed, cannot
be pushed back out of regard for any other presuppositions whatever. But many speak with
apparent recognition of reason and thought, and when it then becomes serious, when the
honesty of the recognition is to show itself in one's bowing to the result, they jump
off, like soldiers who get cannon fever. It can be a long and laborious way for the
individual before he brings harmony between his feelings and interests and the clear
truth that knowledge shows him. But the way must and shall be traveled, sooner or later.
Along this way the race goes forward, just as the individual goes forward along it in
quiet.---Oldenburg was one of those who turn back when the bullets begin to whistle.
And in him this stands out all the more, since he began by being so bold about it and
came forward as the enthusiastic man of progress.

One feels more sympathy in this respect with another of Spinoza's correspondents,
although it went with him in a similar way. A merchant from Dordrecht, Willem van
Bleyenbergh, occupied himself in his leisure hours with philosophical studies and had
read Spinoza's exposition of Descartes's philosophy several times. He now writes to
Spinoza and presents himself as a man who is guided by a need for pure and genuine
truth and seeks \opage{51}to satisfy this need through science. But he, like Oldenburg,
is repelled by Spinoza's determinism. If God acts in everything, he says, then it is
God, after all, who acts in Adam's sin too, and so either God himself must be evil, or
sin must not be anything real and positive.---Spinoza meets him with great kindness.
``Nothing,'' he says, ``do I value more than to enter into friendship with men who
sincerely love truth. Such a friendship is founded on the love that each of them has
for the knowledge of truth, and therefore one can no more dissolve it than one can
refrain from embracing the truth one has once recognized. Nothing but truth can bring
about a deep bond between different dispositions and characters.'' With respect to the
question Bleyenbergh had raised, Spinoza declares that concepts such as good and evil
arise and are applied when we compare actions and events with one another by a certain
standard. But our relative concepts cannot be applied to the absolute being. God can no
more be characterized by concepts taken from human nature than a man by concepts taken
from the nature of the elephant. (In another letter Spinoza says that if the triangle
could speak, it would say that God was ``triangular in an eminent sense''---\textit{eminenter
triangularis,}---and the circle, that he was ``circular in an eminent sense''---\textit{eminenter
circularis.)} We therefore may not ascribe hatred or anger to God either, as the
theologians do, and as Scripture seems to do, since it accommodates itself to the
conception of the common people, for whom God stands as lawgiver and judge \opage{52}and
not as the eternal and infinite being that dwells in everything and from which
everything springs. Because everything happens with necessity, the difference between
the good and the wicked is not effaced, any more than the difference between joy and
sorrow. The good man lives in the clear consciousness of his harmony with the idea of
God, which determines all his actions, while the consciousness of the evil man is filled
with finite and earthly things.---Bleyenbergh now explains his standpoint more closely.
He has two rules by which he philosophizes: the one is the clear and evident conception
of his understanding, the other God's revealed word and will. Even if something stands
before him as clear and evident, he nevertheless holds that he is mistaken if it
conflicts with Scripture. Spinoza replies to this that the Bible can teach us nothing
whatever in philosophical respects. What it teaches is, when one applies a philosophical
measure, often wholly incomprehensible, whereas one can never doubt the truth of what
is proved clearly and distinctly. But what is the use of discussing when one party
expressly declares that a proof makes an impression on him only so far as it agrees
with what Scripture---according to his own or some theologians' interpretation---teaches,
and that rational knowledge is therefore, properly speaking, a matter of indifference to
him?---A visit Bleyenbergh paid to Spinoza in Voorburg brought no closer understanding
about, and Spinoza at last broke off the correspondence in a friendly way and referred
him to oral discussion. Bleyenbergh later published polemical writings against
Spinoza's works.

We cannot here inquire more closely with what right the charge was constantly brought
against Spinoza \opage{53}that his doctrine of necessity (determinism) abolishes all
freedom and morality. In the last chapter of my work ``Om Grundlaget for den humane
Etik'' I have treated the questions that belong here at greater length. The difficulty
there may be here does not in any case stem from determinism in itself, but from the
speculative conception that lies at its basis. And this difficulty is, as Spinoza's
critics always overlook, common to pantheism and theism. Whether one assumes an eternal
substance that stirs in everything or an almighty will that creates everything, it
becomes equally difficult to understand how a finite being can possess an absolutely
free will. Consistent theologians like Augustine (and after him Descartes and Leibnitz)
therefore answer the question in a similar way to Spinoza. The difference is only that
theology places the difficulty a stretch further back, as philosophical difficulties
in general tend to recur in theology at higher powers. The difficulty has in this case
arisen from the deductive method, which is, after all, common to theology and
speculative philosophy. If we apply the inductive method to psychological and ethical
questions, on the other hand, everyone admits that we understand an action only when we
can trace it back to the individual's character, previous history and circumstances of
life. There can in principle be nothing in human actions that could not be explained in
this way if all the circumstances were known to us. Determinism is therefore
unconditionally right and is involuntarily applied in practice by each and every one.
But it does not follow that we are reduced to letting ourselves be carried along by the
stream. We \opage{54}stand toward our own mental and bodily nature as toward outward
nature: if we can make ourselves masters of the conditions, we can also determine the
course of things, divert or restrain the stream; but whether we can do so, only
experience can decide in each particular case. Miracles we can work neither in inner
nor in outer nature.---This way of regarding the matter is, moreover, as we shall see
later, not foreign to Spinoza, although the deductive method is the one that governs his
views.

Spinoza stood in friendly relations with Orobio de Castro, a Spanish Jew, who after
frightful tortures in the prisons of the Inquisition had fled to Holland and there gone
over to the religion of his fathers. He did not share Spinoza's views and later wrote
against him. He sent Spinoza a letter in which Lambert van Welthuysen of Utrecht, an
adherent of the Cartesian philosophy, gives his opinion of Spinoza's
``Theologico-Political Treatise.'' It is remarkable, but not unique, that Welthuysen,
who in his writings shows himself a rationalist and was accused by the theologians of
Socinianism, speaks of Spinoza's doctrine with far greater vehemence than the orthodox
Oldenburg and Bleyenbergh. The author of the work in question, says Welthuysen, has put
off all religion; at the least he goes further than the deists, with whom these ungodly
times swarm, and teaches atheism covertly (\textit{clam}), in that he lets the universe
spring with necessity from God, which is the same as making God one with the world,
and in that he teaches that the worship of God is only to serve to promote virtue and
\opage{55}piety, but has no absolute value in itself. Then all forms of religion become
equally good, the Koran is put on a level with the word of God, and one cannot prove
that Muhammad was a false prophet!---The draft of Spinoza's reply has been found, and
one can see from it that he was for a moment brought out of his contemplative
equilibrium and used strong expressions about the charge that was here directed against
him. Even in the letter he sent Orobio de Castro in reply, he says that he does not know
whether the writer of the letter spoke from malice or from ignorance. Can a man be
accused of having put off all religion when he acknowledges God as the highest good,
who ought to be loved by a free mind, when he finds in this love our highest happiness
and our true freedom, and places the reward of virtue in virtue itself? Is there then
nothing in virtue and knowledge themselves that can give joy, but must one look only to
whether reward and punishment are in store, so that one lets oneself be compelled like
a slave and would rather follow one's desires, were it not for the threatening
punishment? As for Muhammad, he was certainly a false prophet, for he denied the freedom
that the true religion allows. But the spirit of Christ must nevertheless be said to be
present among Turks and heathens alike, when their worship of God consists in justice
and love of mankind.---In a later letter to Welthuysen himself, Spinoza acknowledges his
striving after truth and his sincerity.

In Welthuysen's letter we see Spinoza called an atheist for the first time, though
covertly (\textit{clam}). That was the view of him that was the general one for over a
hundred years. Bayle says of him that he is the first \opage{56}who brought atheism into
a system. Voltaire has given expression to the same view in a satirical poem \textit{(Les
systèmes)}:

\begin{quote}
Alors un petit juif, au long nez, au teint blème,\\
Pauvre, mais satisfait, pensif et retiré,\\
Esprit subtil et creux, moins lu que célébré,\\
Caché sous le manteau de Descartes son maître,\\
Marchant à pas comptés, s'approcha du Grand-Être:\\
``Pardonnez-moi,'' dit il, en lui parlant tout bas,\\
``Je crois, entre nous, que vous n'existez pas.''
\end{quote}
\noindent (See the end of the book.) It was---and to some extent still is---usual to
call all those atheists who depart from the prevailing religious conceptions. To know
what is meant by this word, one must therefore first know what is meant by God. As a
rule one puts into the word atheism something more as well than the theoretical denial
of the Deity; since God stands as the highest that can be thought, he who rejects his
existence must necessarily---so it is inferred---be a low and contemptible man, one who
denies what gives life its worth and higher significance. With respect to the extension
of this concept there is therefore great disagreement. Abbé Mersenne, Descartes's
friend, says that the number of atheists is so great that one must wonder that God lets
them live; in Paris alone there are 50,000---one can find 12 together in a single house!
But among the atheists he then also counts Charron, Machiavelli, Cardanus, Campanella,
Vanini and Fludd.---Eilskov remarks in his ``Filosofiske Breve'' that most people call
Spinoza's doctrine atheism, but he does not think they are right in this. ``Spinoza did
indeed confuse God with the world, in that \opage{57}he made both one substance; but he
nevertheless granted to this one substance God's foremost mark, that it existed by the
necessity of its nature, consequently that everything that happened in the world was a
consequence of the unalterable connection of things.'' Eilskov would therefore rather
call Spinoza a deist than an atheist, rather say of him that he denied providence than
that he denied God.---What separated Spinoza from the so-called deists was that the
latter did not, to be sure, assume a providential government of particulars or a
supernatural revelation, because everything since the creation went on according to the
laws of nature established once and for all, but nevertheless assumed a personal God.
What gave offense in Spinoza was precisely that he denied God as a single, determinate
being, an object of representation and imagination. God did not stand for him outside
nature,\footnote{Compare Heiberg's ``Protestantismen i Naturen'': The God you seek is
your own God; what more would you have? He does not step forth as an object among many
others. In seeking him, you have already found him; in asking, the answer is already
won.} but was what works and is in all nature. But for such a view a word was lacking,
although it was not the first time it had been expressed. The word pantheism seems to
have been introduced first by the English philosopher John Toland at the beginning of
the eighteenth century. The new word was like a portent of the change in conception
that was to take place. Toward the end of the century the greatest \opage{58}minds vied, after Lessing's
example, in praising Spinoza; the pious Novalis called him ``a God-intoxicated man,''
and Schleiermacher in his ``Speeches on Religion'' pointed to Spinoza as a hero of
religiousness. Only now did the positive side of Spinoza become the object of more
general recognition; until then people had had eyes only for the negative. As long as
the dualistic conception of nature continued to prevail, it could not be otherwise
either. Spinoza was quite right when he said that his opponents conceived of nature in
an entirely different way from him. For him the essential thing in nature was the
eternal power and the eternal laws that reveal themselves in everything. But such a
value his opponents, given their dualistic views, could not ascribe to nature; it stood
for them as a power that drags man down, as finite and external, as a barrier beyond
which they yearned.

It is illuminating with regard to Spinoza's concept of God, what he states in one of
his letters, that he has just as clear a concept of God as of a triangle, but not as
clear an image; we can think God, but not imagine him. Imagination gives us a single
image, encloses the object in a given frame, but thought grasps---not the single object
by itself---but the relation and connection between objects. That God is thought, not
imagined, thus means that he is conceived as the law or the power in everything, not as
a single being subsisting by itself that has something outside itself. Precisely
because the usual conception wants to grasp God in an image, it comes \opage{59}to set
``the world'' outside him. Yet there is hardly so great a difference between
imagination and thought as Spinoza assumes. Although thought overcomes the oppositions
and the glaring externality in which imagination easily entangles itself, it cannot
itself avoid positing new oppositions; it establishes relations, lawfulness and
connection, but thereby always presupposes a plurality of terms between which the
relations hold and the laws apply. An absolute unity is the absolutely
incomprehensible, precisely because it would not stand in relation to anything
whatever. Our concepts do not reach beyond the world of phenomena and relations. That
is not to say that our world is the only or the highest one; a consistent theory of
knowledge leads to assuming an infinity of possibilities. What we can maintain is only
that the truth and necessity we know must have their place and their validity within a
possible larger connection. What this place is, we lack every means of deciding. But in
working for what is true, beautiful and good, we are sure that we are working in the
service of true reality. For such a monistic world-view Spinoza stands as the great
prophet, although he went by speculative paths that we cannot tread with the same
confidence as he.

---We have mentioned in an earlier place a letter from Spinoza's faithful friend Simon
de Vries, in which the latter expresses his sorrow at being separated from him and
calls a housemate of Spinoza's happy, because he always has access to him. In his
reply Spinoza says that de Vries \opage{60}should not envy his housemate: ``There is no
one who is more of a burden to me, and with whom I am more on my guard. Therefore I ask
you and all of you not to communicate my views to him until he has reached a more
mature age. He is still too childish and inconstant, and is more eager for what is new
than for what is true. But I hope that after the passage of some years he will correct
these childish faults, and as far as I can judge from his character, I am sure of it,
which is why I also love him for his good nature.'' It is no doubt the same young man
of whom Spinoza tells in another letter that he would not initiate him into his own
views and had therefore prepared an exposition of Descartes's philosophy for him. And
finally we presumably meet the same person again in Spinoza's correspondence under the
name of Albert de Burgh. He writes to Spinoza from Florence and begins by telling that
he ``by God's infinite mercy has become a member of the Catholic Church.'' He mentions
his former admiration for Spinoza's understanding and acuteness and for his ardent
search for truth; but all the more does he now grieve that Spinoza is led about and
enticed by ``the wretched and insolent prince of the evil spirits.'' What is your whole
philosophy but illusions and fantasies? And yet you build on it not only your spiritual
peace in this life, but also the eternal salvation of your soul? You believe that you
have found the best philosophy; but there have been and are so many philosophers who
disagree with you! In your ``Theologico-Political Treatise'' you have tried with
devilish cunning to separate \opage{61}theology from philosophy, although in you
yourself they coincide. You claim that what you say there about the teaching of
Scripture you have learned from Scripture itself. But how dare you rely on your own
interpretation? Only by the testimony of the whole Church and the apostolic traditions
can Scripture be interpreted. You cannot appeal to the Protestants, for they are just
as wretched as you and walk in just as great darkness. How dare you, a paltry man, a
miserable worm that shall itself become food for worms, rely on your reason against the
testimonies of the many holy and wise men the Church can show? Dare you consider
yourself wiser than patriarchs, prophets, apostles, indeed Christ himself? Do you think
you can explain everything? Is it not devilish pride that you would judge the terrible
mysteries of Christ's life and death, which the Church declares incomprehensible? For
so many years the Church has stood firm and unalterable, confirmed by miracles and by
the heroism of the martyrs, and instead of bowing before it you would rather trust in
the offspring of your own brain and make yourself more wretched than the beasts, however
foolishly puffed up you are over the excellence of your mind and your teaching!---This
letter I have written to you with a Christian intention, to show my love for you,
although you are a heathen, and to beg you not to go on corrupting others.

Spinoza at first did not want to answer this letter, because he thought that reasons
would make no impression here. But some common friends, who like himself had cherished
good hopes for Burgh, persuaded him to make an attempt. He rebuts the young man's
fanaticism in a dignified manner. \opage{62}---I do not imagine, he says, that I have
found the best philosophy, but I believe that I have insight into the true philosophy.
There is no conviction so clear and certain as that which rests on clear reasons. The
true makes itself evident to thought and thereby at the same time teaches us where
falsehood lies \textit{(verum est index sui et falsi)}.\footnote{In one of his letters
to Bleyenbergh Spinoza says: When I have found a solid proof, I can never fall into such
thoughts that I doubt what I have proved, and therefore I rest content with what the
understanding teaches me, without suspicion that I am deceived. And even if I come to
see that I have been mistaken, I still regard it as good fortune, for I thereby take a
step forward, and that can only increase the joy and tranquility with which I strive to
lead my life.} But you, who claim that you have found the best religion, what reasons
can you give for it, since you do not trust reason? If you appeal to the inner
testimony of the Spirit, then all others can do the same. If you point to the
excellent and holy men who have lived in the Roman Church, then all the sects among the
Protestants can each bring a similar proof. All have had their martyrs. There were, as
you know, some of your own family who under the Duke of Alba's regime suffered every
possible torment for the sake of their faith. I myself have known a Jew by the name of
Judas who, half dead, sang a psalm on the pyre of the Inquisition. If you point to the
Roman Church's long past and continuous tradition, it is far surpassed by that of the
Jews, who trace their lineage and their traditions right back to Adam, and whose
religious community stands unshaken to \opage{63}this day, in spite of all persecutions
past and present.---You, who once worshipped an infinite being, now worship a God whom
Chastillon let his horses eat in the town of Tienen.\footnote{The Huguenot Marshal
Chastillon threw consecrated hosts to the horses of his cavalrymen.} Acknowledge, then,
that true Christianity consists in justice and love, and stop calling absurd errors
mysteries!---

While the correspondents named so far\footnote{Yet another name, and a very famous one,
seems as if it must be mentioned here. Mr.~A.~D. Jørgensen, namely, has in the course
of occupying himself with Steno's biography and writings come to the conviction that
Steno, who studied in Leyden while Spinoza lived in nearby Rhynsburg, and who himself
occupied himself with Cartesian philosophy, stood in close relation to Spinoza. His
``\textit{Epistola ad novæ philosophiæ reformatorem},'' which was written after the
famous anatomist and geologist's conversion to Catholicism, where he found peace for
his brooding mind, seems to be directed against Spinoza (particularly against the
``Theol.-Pol. Treatise''). He speaks of his opponent with great respect and alludes to
their former intimate acquaintance. Unfortunately the work is not to be found in our
libraries.} essentially represent the struggle of the old world-view against the new
doctrine, which is incomprehensible to it, we still have two groups of Spinoza's
correspondents left, both of which stand in a sympathetic relation to him; the one
consists of his disciples proper, who adopted his doctrine, the other of those who see
through it in such a way that they are able to direct a real philosophical critique
against it and point beyond it to new paths for philosophical inquiry.

We have spoken above of Simon de Vries and his relation to Spinoza. In a similar
beautiful relation of piety \opage{64}stood several young men, most of them physicians
and of Jewish descent, belonging to the little circle that studied his doctrine from
his manuscripts. A couple of beautiful letters have been found to one of these,
Dr.~med. Bresser. In the one, Spinoza complains that Bresser has left without taking
leave of him. He urges him to correspond and encourages him to lay down a fixed plan
of study and to devote the best part of his life to the cultivation of his mind. ``So
that you may take courage to write to me with greater freedom, you must know that I
suspect that you have far too great a distrust of your abilities and therefore fear to
come forward with questions or statements that might betray ignorance.''---The other
letter deals with the right method one should apply in one's studies. Spinoza stresses
that---apart from the purely logical side of the matter---what counts is steady work
of thought and firmness of mind and will, which cannot be attained unless one settles
on a definite way of life and prescribes for oneself a definite goal.---In the advice
he here gives his friend there lies an indication of the road he himself had traveled.

To the last group belong Leibnitz and Tschirnhausen. Of these Leibnitz has a famous
name in the history of science; Tschirnhausen first really acquired interest for the
history of philosophy when, through the previously mentioned discovery of Spinoza's
manuscripts, it was found that he was the hitherto unknown author of the most acute of
the letters addressed to Spinoza. However much these thinkers have in common with
Spinoza, they nevertheless go beyond his standpoint. They at the same time represent
the reaction against him, inasmuch as \opage{65}Leibnitz in particular again lets
philosophy enter into a compromise with theology that was by no means to the advantage
of the clarity and thoroughness of his thought. On the whole Leibnitz was the complete
opposite of Spinoza. In place of the latter's concentrated calm and absorption in one
thought there appears in the former an astonishing versatility: in mathematics and
physics, in philosophy, history and jurisprudence he is at home, and everywhere he
intervenes in an original way and makes new discoveries. While Spinoza keeps away from
public life and spends his time grinding his lenses and pursuing his studies in his
frugal room, Leibnitz throws himself into a whirl of practical affairs, frequents the
courts, philosophizes with queens and princesses and theologizes with high
ecclesiastical dignitaries. This many-sidedness and outward orientation were not to the
advantage of his character as thinker and man. He readily came to terms and showed
considerations he ought not to have shown. His conduct toward Spinoza shows him in a
particularly unattractive light. In Spinoza's correspondence there is only one letter
from Leibnitz to Spinoza and one from the latter to the former, and both deal with
optical, not with philosophical, questions. But from the manuscripts found in their
time we obtain new information here as well. Tschirnhausen, who met Leibnitz in Paris,
describes him as ``free of the usual theological prejudices,'' and tells that he set
great store by Spinoza's ``Theologico-Political Treatise''; Leibnitz had even,
according to his statement to Tschirnhausen, written an approving letter to Spinoza
about this book; but it has unfortunately been lost. Later he has hinted that for a
time he inclined toward Spinozism. Tschirnhausen \opage{66}asks Spinoza for permission
to acquaint Leibnitz with the manuscript of the Ethics. But Spinoza will not yet give
his consent to this. He can see from Leibnitz's letters, to be sure, that he is a man
of free outlook and versed in all the sciences; but he will nevertheless wait to
communicate his writings to him until he obtains closer information about his
character and his activity in Paris.\footnote{Leibnitz was at that time in the service
of the Elector of Mainz and had traveled to Paris with great political plans.} The
distrust that Spinoza thus intimates proved here---just as with Albert de Burgh---not
to be unfounded. When Leibnitz, after his visit to Paris and London, returned to
Germany by way of Holland, he visited Spinoza several times, as he himself relates, and
talked with him at great length. But officially he would later not properly acknowledge
this acquaintance. In his \textit{Théodicée} he relates in passing that he visited
Spinoza, but gives the impression that their conversation turned on ``anecdotes of the
time.'' To judge from the relation in which he then stood to Spinoza, and from the
great problems with which both men were occupied, they surely talked about other things
than anecdotes. Later Leibnitz in a letter described Spinoza as ``\textit{novateur trop
connu}.''---The acquaintance would have been compromising, especially for one who
wished to have the appearance of being orthodox and wished to reconcile faith and
knowledge. But what did Leibnitz gain by all his flirting with church doctrine? People
did not believe him; the people of Hanover changed his name to ``Løwenichts,'' i.e.\
Glaubenichts, and when he died no clergyman followed his coffin.

\opage{67}It must not be forgotten, however, that Leibnitz really worked out a
philosophical view whose strength lay precisely where Spinoza's weak side had been.
Spinoza traced everything back to the one substance; but Leibnitz started from the
idea that one and the same universe is apprehended in infinitely different ways at the
different stages of development at which existing beings find themselves; there are
thus really infinitely many worlds or substances, which he called monads. This way of
looking at things is not entirely foreign to Spinoza's system, but plays only a
subordinate role there, and if definitely brought forward would lead to conflict with
his fundamental view. The intermediary between Spinoza and Leibnitz, Tschirnhausen, has
pointed precisely to this point in his letters to Spinoza.

Walther Ehrenfried von Tschirnhausen was a Saxon nobleman who in his youth traveled to
Holland to study. He came into contact with the circle of Spinoza's young friends and
through them with Spinoza himself. He is known as a mathematician and physicist (he is,
among other things, the inventor of the burning glass) and as a philosopher stood close
to Spinoza and Leibnitz. His work \textit{Medicina mentis}, a kind of logic, has so
much in common with Spinoza's treatise ``On the Emendation of the Intellect'' that he
has even been accused of plagiarism. In any case his most important doctrines are a
further development of what Spinoza had already taught, and he often uses even the
same words and turns of phrase. In a few places he alludes to Spinoza, but never names
him (only as \textit{qvidam} or in a similar way), naturally for the same reason as
\opage{68}the one that led Leibnitz to such great caution. Tschirnhausen had,
moreover, dedicated his work to Louis XIV in order to obtain a pension.---Here he
interests us as the author of those letters in Spinoza's correspondence that betray the
sharpest eye for the system's fundamental concepts and thereby also for its weak sides.
It is significant that these letters were written at the time when Tschirnhausen was
associating with Leibnitz in Paris; the points they touch on are precisely the ones
that came into their own in Leibnitz's philosophy.

The exchange between Tschirnhausen and Spinoza concerns two main points. How can one
derive from the general concept of extension the manifold actually existing bodies with
their determinate figure and motion? With this question the whole deductive and
abstractly mathematical method that Descartes had introduced into philosophy, and that
Spinoza had brought to its completion, is struck. In extension in itself there lies
neither figure nor motion. Descartes, as Tschirnhausen remarks, was even better placed
on this point; he saw that motion could not be derived from mere extension, and assumed
that God at the creation had imparted to matter the motion it had. But Spinoza, who has
of course given up this theological notion, cannot explain actual
motion.\footnote{This objection to Spinoza's doctrine was later set out at greater
length in an acute essay by John Toland in his Letters to Serena (1704).} Spinoza
replies to this that Descartes's concept of matter as mass at rest was unsatisfactory,
but that he himself had not yet properly \opage{69}thought the matter through. He
expresses himself in this debate with a lack of assurance that must be explained by the
fact that here he stood at the historical limit of his thought. A new concept was
necessary in the conception of nature, a concept that Leibnitz in particular brought
forward both in physics and in philosophy: the concept of force. Spinoza has, to be
sure, as we shall see, put this concept in the closest connection with his fundamental
concept, substance, but it does not yet pervade his conception of nature, and on closer
consideration it proves to stand in disharmony with the concept of substance as Spinoza
conceives it.---But what offended Tschirnhausen in Spinoza's conception of matter was
not merely its one-sided determination as extension in space, whereby the properly
physical was set aside; it was also the demand that material phenomena should be
derived from one single concept. In order to draw conclusions, he says, we always need
several concepts; from a single one we cannot derive a plurality of determinations.
When, for example, I consider the circumference of the circle alone, I can conclude
nothing from it except that it is everywhere equal to itself and uniform, whereby it
differs from other curves; but other properties I cannot derive. When, on the other
hand, I also take other things into account, namely the radii drawn from the center,
two or more lines intersecting one another, etc., then I can derive further properties
from this. Spinoza replies that this holds in mathematics and logic, but not as regards
real concepts, and therefore not for the concept of substance either; but he does not
indicate where this difference should \opage{70}come from. Tschirnhausen's objection is
of great interest because it points to the law of relativity, which is more and more
recognized in philosophy. Every sensation, every representation and every thought
stands in connection with other sensations, representations and thoughts, and is
possible only in and through such a connection. Before something absolute and
self-contained in and for itself, thought therefore stands powerless; it cannot think
it, because there is nothing with which it can compare and measure it. Tschirnhausen's
objection, which is directed only at the concept of matter, strikes the very
fundamental concept in Spinoza; for his ``substance'' is that which exists in itself
and is to be understood through itself; but everything we understand, we always
understand by means of something else. The necessity of striking out on another path
and using another method than Spinoza's becomes evident here.

Tschirnhausen's second question concerns Spinoza's doctrine that substance reveals
itself under infinitely many attributes or fundamental forms, but that we know only two
of them (mind and matter). Spinoza's line of thought was this: the more perfect a being
is, the more fundamental forms it must also have under which to reveal itself, and
since God is the infinite being, his essence must express itself in an infinity of
fundamental forms. But, says Tschirnhausen, since every thing belongs to the essence of
God or substance, every thing must be expressed under infinitely many fundamental
forms, and there are thus really infinitely many worlds. And even if it is one
substance that exists under the infinitely many attributes, one universe \opage{71}that
is apprehended in infinitely many ways, still we know only two of these ways, and the
other fundamental forms thus become like foreign worlds to us.---Here we are brought
close to the change that Leibnitz, as just indicated, introduced into the concept of
substance. Leibnitz's monads correspond to the infinitely many attributes of the
Spinozistic substance, for the world is represented in each monad in a peculiar way.
By putting the monad in place of substance Leibnitz paved the way for the critical
philosophy, which starts from the idea that what we know about the world is conditioned
by the nature and conditions of our consciousness and therefore holds only from our
standpoint. Spinoza boldly started from the idea of the absolute being and maintained
an absolute truth; he still saw only imperfectly that ``what is seen depends on the
eyes that see,'' and that we know light only as refracted into a multiplicity of rays.
It is a truth that English philosophy, Leibnitz, Kant, the physiology of the senses
and the hypothesis of evolution have each in its own way expressed.

Thus we see Spinoza in his correspondence placed between past and future. We see his
thought in its justified struggle against prejudices, as well as in its limitation,
where it meets the seeds of the future.

% ===== SPINOZA'S PHILOSOPHY (pp. 72--164); the contents page: "Spinozas Lære" =====
\chapter*{Spinoza's Philosophy}
\addcontentsline{toc}{chapter}{Spinoza's Doctrine}
\markboth{Spinoza's Philosophy}{Spinoza's Philosophy}

\section*{I.}
\addcontentsline{toc}{section}{1. Fundamental Concepts}
\opage{72}Spinoza begins his main work, the Ethics, with a series of definitions that
already give us his whole system in brief. ``By cause of itself''---so runs the first
definition---``I understand that whose essence involves existence, or, in other words,
that which cannot be conceived except as existing.'' To this the third definition
attaches itself: ``By substance I understand that which is in itself and is conceived
through itself, that is, that whose concept does not need the concept of any other
thing from which it must be formed.'' That which exists in and through itself must also
be the cause of itself. For what springs from something else, or has its cause outside
itself, cannot be understood in and through itself alone, and so cannot be substance.
As cause of itself, substance is eternal. For by eternity I understand, says Spinoza,
the very existence of things, insofar as it follows from the concept of things.
Substance is, further, infinite. For finitude is limitation and negation, \opage{73}and
if substance were finite, its existence would depend on something other than its own
essence.

The essence and content of substance find their expression in the attributes. By
attribute Spinoza understands what thought conceives of substance as making up
(constituting) its essence. The attributes are definitions of substance, are each by
itself substance itself from one particular side. The attribute must therefore, like
substance, be conceived in and through itself: each attribute by itself is infinite and
eternal, and the one attribute cannot be explained by means of the other, the one has
not arisen from the other.

While substance is that which exists in itself and is explained through itself,
Spinoza understands by mode \textit{(modus)} that which has its existence in something
else and is understood only through this. All modes are therefore more particular
determinations \textit{(affectiones)} of substance. Everything that is thus falls into
two classes: either it has its existence in itself (substance) or in another (modi),---
either it has an unconditioned or a conditioned existence.

These are the three fundamental concepts in Spinoza. If we look for Danish names for
them, we might perhaps call substance \textit{Grundvæsen} (fundamental being),
attribute \textit{Grundform} (fundamental form), mode \textit{Enkeltvæsen} (individual
being). The individual beings thus have their existence in the fundamental being, and
this manifests itself in the fundamental forms. The higher a being stands, the more
fundamental forms it has. The one and infinite fundamental being, which exists under
infinitely many fundamental forms, we call God. God exists with necessity; his
existence is grounded in his essence. Or \opage{74}suppose that God does not exist:
there would then have to be something that hindered and excluded his existence; but
where should this something be, since God is the infinite fundamental being outside of
which there is nothing? Besides God no fundamental being can be given or conceived.
Everything that exists thus has existence only as a particular form of the one
fundamental being.

As that which is conceived in and through itself, the fundamental being is the starting
point of thought; only through it is everything else understood. The right way to
philosophize is, says Spinoza, to start from the highest cause. But since all existence
is included in the fundamental being, Spinoza's philosophy not only has its starting
point in this concept; it cannot, if it wants to be consistent, get beyond it. If
Spinoza assumes an existence that lies outside the fundamental being, he contradicts
himself: there would then be something outside the one that is.

Of the infinitely many fundamental forms of the infinite being we know only two:
thought and extension. By extension Spinoza understands, as we know, what we call
matter, since he regards the geometrical properties as the essential ones. Extension is
a form of God's essence. One is right, to be sure, that God is not corporeal; for a body
is a determinate and limited figure. But extended or corporeal nature is not outside
God or created by God. The fundamental form (the attribute) is, after all, the
fundamental being (substance) itself, eternal and infinite as this is. People regard it
as unworthy to make the corporeal an expression of God's essence; but that is only
because they always think \opage{75}of finite bodies, which can be divided and
dissolved. The fundamental being as such is indivisible and eternal. Thus water as
water can be divided, and its parts separated from one another, but not insofar as it
is substance.---Spinoza stands here as a forerunner of the doctrine of the conservation
of matter, which modern chemistry has confirmed experimentally. The various particular
substances can be dissolved and converted into one another, but no matter perishes, no
atom is added and none is lost.---While Spinoza conducts a special defense of his
ascribing to God the fundamental form of extension, he does not do so as regards the
fundamental form of thought, since here he can rely on the common conception of God as
spirit.

According to Spinoza, then, we know God as the infinite power that works in the world
of thought as well as in that of matter. But we are naturally led here to ask why we
know only these two fundamental forms among the infinitely many. The reason does not lie
far off: we know ourselves as conscious and material beings, and fundamental forms other
than consciousness and matter are not accessible to us in the world of experience. This
leads us to see how Spinoza arrived at the concept of attribute or fundamental form at
all. He lays down this concept simply by a definition, without mentioning its origin.
But it is really won by an analysis of what is given in experience. Everything we know
can be referred to two great groups: phenomena of consciousness and material phenomena.
This is the deepest-lying difference we find in the whole world of experience. That
Spinoza's ``attribute'' is really only the characteristic mark \opage{76}of one of these
groups can be seen from the fact that he (right at the beginning of the ``Ethics'') uses
the expression ``kind'' in the same sense as ``attribute.'' In the same sense he
further also uses the expression ``nature,'' and in an earlier work ``property or
quality'' \textit{(proprietas sive qvalitas sive attributum. Princ. phil. Cartes. I.
def. 5)}. The attribute is thus really a quality, a peculiarity, that gives a certain
group or class of phenomena its fundamental stamp. Spinoza sets up as the 5th axiom in
the second book of the ``Ethics'': ``We feel and perceive no particular things other
than bodies and modes of thought''; and relying on this he then proves that extension
and thought are divine attributes. Here the real origin of the concept lies plainly in
view.

Spinoza conceives the attribute as the expression of the essence of substance. Where,
then, does he get the idea of substance? With him it is immediately given, as what is in
itself clearest and most evident in our knowledge. He therefore determines it straight
away by a definition. Its being (God's existence) is an eternal truth. But if we look
more closely, we shall be able to trace the real origin here too. In his exposition of
Descartes's philosophy Spinoza says (I. def. 5): ``Every thing in which, or through
which, something that we perceive, i.e., a property, a quality or an attribute, has its
existence, is called substance.'' ``Being as such, as substance in and for itself, does
not affect us, and therefore it must be explained by an attribute, from which it can
nevertheless be distinguished only logically'' \textit{(Cogit. metaph. I, 3, 1)}. Here
we see the traces of the method by which knowledge \opage{77}in reality climbs upward
to reach what is substantial in things: experience and the analysis of
experiences.\footnote{In \emph{Descartes} it is stated still more definitely that the
concept of substance is reached in this way. ``Substance cannot originally be known
merely from its being an existing thing, for that by itself alone does not affect us;
but we easily come to know it from one of its attributes. For from seeing that some
attribute or other is present, we \emph{conclude} that there must necessarily also be
present an existing thing or a substance to which this attribute can be ascribed''
(\textit{Princ. phil. 1. 52}).---In his polemic against Descartes, \emph{Gassendi}
rightly lays great weight on the fact that we reach the concept of substance by
inference. That is why, he says, this concept is so abstract, indistinct and
problematic, and it is not the knowledge of the properties (the accidents) that depends
on the knowledge of substance, but conversely the knowledge of substance that depends
on the knowledge of the properties. (\textit{Disqvisitio metaphysica. Amstelod. 1644
p. 121---123}).}

But experience can, says Spinoza, teach us only the existence of finite things, not the
fundamental being and the fundamental forms. The particular things we can learn from
experience we understand rightly only when we already know the inner essence of
things.\footnote{In a note to the unfinished treatise ``On the Emendation of the
Intellect'' Spinoza says that he intends to examine in detail the empirical method that
the modern philosophers employ. Unfortunately he did not carry out this plan.} He
therefore skips the path of experience and analysis as useless and begins by defining
the idea of the fundamental being, which he regards as immediately clear and evident.
It would be a self-contradiction for him that what includes all being in itself should
be known through something else. Experience is for him a ladder that reaches only
halfway, and which it occurs to him all the less to use, \opage{78}as he feels himself
on the heights from the very beginning. When a mathematician defines, his definition is
one with a construction, since it states how the thing defined comes into being. In the
definition of the circle as a line in which every point is at the same distance from a
certain given point lies the instruction for constructing a circle by letting a
straight line move about one of its end points. But for such a construction we have
nothing corresponding as regards the abstract metaphysical concepts. For Spinoza it
stood as a contradiction that what is in itself should not exist. ``To be able to exist
is a power, and the more reality belongs to the nature of a thing, the more powers it
has in itself to exist: hence the absolutely infinite being, or God, has absolute power
in itself to exist, and therefore exists absolutely. Many will perhaps not easily see
the evidence of this proof, because they are accustomed to consider only the things
that spring from external causes. But I am speaking here only of the fundamental being,
which cannot be produced by any external cause. For the things that arise through
external causes receive all the perfection or reality they have from the power of the
external cause, and their existence springs solely from the perfection of the external
cause, not from their own. On the other hand, the perfection that the fundamental being
possesses is owed to no external cause, and therefore its existence must follow from
its nature alone. There is therefore nothing of whose existence we can be more certain
than of the existence of the absolutely infinite and perfect being.'' We now understand
why Spinoza begins his system with the definition of the concept ``cause of itself.''
\opage{79}We understand, further, that this concept could not be grounded in
experience, which always shows us one thing conditioned and effected by another. It is
thus an idea of reason that Spinoza presupposes, and just as anyone who has the ability
to form and intuit figures in space can be compelled to acknowledge the reality of the
circle, so Spinoza holds that anyone who thinks through the concept ``cause of itself''
must assume the existence of such a cause.

Beyond and behind this concept we cannot get in Spinoza. Only by taking an entirely
different starting point and applying another method can one come to put the concept of
what is in itself up for debate. That happens in the empirical philosophy in England
and in the critical philosophy in Germany. What for Spinoza stood as the most certain
of all, the first and deepest starting point of thought, now becomes precisely the most
problematic, the most difficult concept of all. Kant's concept of the ``thing in
itself'' (das Ding an sich) is directly reminiscent of Spinoza's definition of
substance; but Kant designates by it precisely what constantly eludes our knowledge,
the limiting concept in the world of thought.

In his ``Short Treatise on God and Man'' Spinoza mentions (I. 7) the objection that God
cannot be defined, because every definition takes place by referring the thing defined
to a certain kind or class, whereas God, of course, cannot be referred to any class. He
grants that this is the doctrine of ordinary logic about definition, but does not
hesitate to reject it as regards the fundamental being. Substance, says \opage{80}he,
is after all a being existing through itself, which is known in and through itself. It
is only of the things that do not exist through themselves that it holds that they are
defined by being referred to a certain class. Here, then, Spinoza maintains that the
knowledge we have of substance is of an entirely different kind from the rest of our
knowledge. The question then is whether there really is such a double mode of
knowledge, or whether the kind and manner of our knowledge is not essentially the same
throughout. We shall return to this question in treating Spinoza's theory of knowledge.

As we have seen above, Tschirnhausen had already stopped at the difficulty that the one
fundamental being should reveal itself under infinitely many fundamental forms, of which
we nevertheless know only two. Two tendencies in Spinoza clash here, as Kuno Fischer
has clearly shown. According to his fundamental thought, that all is one, what holds
for substance must also hold for the individual beings; but more powerful than this
striving to conceive the absolute as immanent is his presupposition of its infinite
exaltation above every finite being. This exaltation shows itself in the fact that,
while we as finite beings have only two attributes, the infinite being has infinitely
many. Here we thus have a point where it can be seen with how much right I maintained
in the Introduction that Spinoza's speculative principle, the idea of the absolute
substance, is rightly understood only when one goes beyond philosophy, back to the
religious consciousness. In a note to the ``Short Treatise'' just mentioned
\opage{81}(I, 1)\footnote{It must be remarked, however, that it is not quite certain
whether the notes to this work derive from Spinoza himself.} it is said: ``After we
have considered nature, we have so far been unable to find in it more than two
fundamental forms that belong to the all-perfect being itself. These are not now enough
for us; they cannot be the only thing of which this most perfect being consists. On the
contrary, we find in ourselves something that plainly announces to us not merely more,
but infinitely many perfect fundamental forms, which must belong to that being for it
to be called perfect. But this something cannot spring from the two fundamental forms;
for two give only two and not infinitely many. From where, then? Not from myself; for I
cannot give what I do not have. From where else, then, but from the infinite
fundamental forms themselves, which tell us that they are, without at the same time
telling us \emph{what} they are; for only of two do we know what they are.'' This is the
only hint that occurs in Spinoza's writings in this respect, and it has an almost
mystical character. This side of the matter, too, we shall return to under Spinoza's
theory of knowledge. That a great truth nevertheless lies hidden in Spinoza's idea of
the infinitely many fundamental forms of the fundamental being, we have already seen
above. Precisely when we see the relativity of our knowledge, its dependence on our
nature and on our position in existence, and thus renounce an absolute knowledge, we
cannot but grant the possibility % the print: „Mulighednn“ (sense: Muligheden)
of infinitely many other standpoints on the same existence.

Although Spinoza, by calling the fundamental being God, \opage{82}indicates this
concept's kinship with the object of religious faith, he nevertheless gives more
particular determinations that remove the popular notions of God. Nothing, he teaches,
can be or be conceived without God. But from the necessity of the divine nature there
follows an infinity of effects. God acts according to the necessity of his nature;
there is nothing outside him that could determine or hinder his acting. This is true
freedom, which is one with inner necessity. God is the only free cause: finite beings
all have something outside themselves that hinders them from following their nature.
And since it is grounded in God's nature that infinite effects spring from it, these
effects must be eternal. God has from eternity produced everything that he produces. If
anything remained that he had not yet managed to produce, his omnipotence would not be
eternal. In God there is no difference between possibility and actuality. His essence
and acting are an eternal actuality. Time, number and measure find no application to
the fundamental being. In the eternal there is no past or future. When one ascribes
intellect and will to God, these words must be taken in an entirely different sense
from when they are used of men; the likeness becomes merely a likeness of name, like
that between the dog, the barking animal, and the constellation of the same name. It
would be especially wrong to think of God as acting according to purposes; that would be
to subordinate him to a higher power, to place something outside him by which he
directed himself in his acting; that is, it would be to abolish the very concept of
God. The word personality, which the theologians use to designate \opage{83}God's
essence, is far too indefinite and unclear to be used with profit.

The question presents itself here how the divine acting is to be understood when all
difference of time and all difference between possibility and actuality is to be
excluded. Light on this lies, in the first place, in the very concept ``cause of
itself.'' The acting of the fundamental being produces nothing but itself: its
producing of everything coincides with its producing of itself. ``The essence and
existence of things,'' says Spinoza, ``are a necessary consequence of the divine nature.
God must be called the cause of all things in the same sense in which he is called the
cause of himself. Finite beings are nothing but determinations \textit{(affectiones)} of
God's attributes, or forms \textit{(modi)}, in which God's attributes are expressed in a
certain determinate way.'' Just as by understanding we understand both the source of
thoughts and the sum of all thoughts, so God is the source of all things and the sum of
all things, the cause of everything and the totality of everything. God is, in other
words, an indwelling (immanent, reflexive) cause, not a transcending (transcendent)
cause; what he produces is not something outside him, but a form of his own essence.---
But in the second place it must be stressed that Spinoza often uses, as a paraphrase
for ``the acting of the Deity,'' the expression ``consequence of the nature of the
Deity,'' and even suggests that this latter designation is the most correct. The
comparison Spinoza so often uses points in the same direction: that everything follows
from God's essence with the same necessity with which it follows from the essence of
the triangle that the sum of its angles is equal \opage{84}to two right angles. The
relation between ground and consequence thus takes the place of that between cause and
effect. In the relation between ground and consequence time plays no role; the
consequence does not come to be, but is eternally one with the ground. Spinoza's
conception thus becomes this: that everything is as in an eternal now, because it has
its ground in the infinite substance, where no differences of time hold.

Just as the concept of the fundamental being, on consistent consideration, coincided
with the concept of God, so the concept of God in turn coincides with the concept of
nature. ``That infinite being, which we call God or nature, acts with the same necessity
with which it exists.'' But Spinoza now distinguishes between two forms of the essence
of the Deity or of nature: \emph{producing} nature \textit{(natura naturans)}, God as
free cause, and \emph{produced} nature \textit{(natura naturata)}, the totality of
individual beings, which can neither be nor be conceived without God. And within
produced nature a further distinction is made between what springs \emph{immediately}
from God's essence and what springs from it only \emph{mediately}. What follows
immediately from God's eternal and infinite nature must itself be eternal and infinite.
By this one is to think of the order of nature as a whole, the laws holding in the whole
universe, within the world of mind as well as within the world of matter. But what is
finite and has a limited existence cannot be produced immediately by God's absolute
nature. The finite can therefore be produced by God only insofar as he appears in
limited and finite form: the one finite thing is thus determined by another, this by a
third, and so on to infinity. Each individual being is a link \opage{85}in the chain of
nature and is determined by the other links in this chain.---Here an entirely different
causal relation is introduced from the one that holds for the fundamental being. The
individual beings stand independent, outside of and over against one another, and act as
mechanical causes, whereas the fundamental being acts only as an indwelling and
reflexive cause, never upon anything but itself, never meeting resistance.

Spinoza says in one of his letters that the concept of necessity is the chief
foundation of his doctrine. He denies that there is anything contingent. ``If men
clearly understood the whole order of nature, they would find everything just as
necessary as all the propositions of mathematics.'' But necessity thus appears in his
system in a double form: as inner necessity in the nature and acting of the fundamental
being, and as outer necessity in the interaction and mutual determinations of the
individual beings. Now since the fundamental being is Spinoza's first and last thought,
he ought consistently to have stopped at inner necessity; but then the world of
individual beings, concrete actuality, would have received no reality. Every particular
and finite thing is, according to Spinoza, a limitation and a determination of the
essence of substance; and, he adds, every determination is a negation. For the
determination does not necessarily belong to the nature of the thing: where the limit
is set, the thing, after all, ceases to be. Every determinate figure is thus a negation,
an abolition of infinite extension. To posit a particular finite thing as existing is to
abolish God's existence. Spinoza does teach expressly, to be sure, that although every
particular thing is determined by another thing \opage{86}to exist in a certain way, the
power by which it persists in its existence nevertheless follows from the eternal
necessity of God's nature. But this power must nevertheless always itself be a
determination and thus, on Spinoza's conception, a negation; it must be related to the
infinite divine power as figure is related to the whole of matter, and from the essence
of substance one cannot derive what abolishes substance. In Spinoza's system there can
really be no room for an outer necessity, an interaction between real individual beings.
If these are each of them negations, they cannot affect and determine one another; for
from nothing comes nothing. When Spinoza nevertheless teaches such an outer necessity
and presents nature as a world of elements that shape themselves into greater and
smaller wholes according to the laws of their interaction, this doctrine is not won by
the deductive path. It rests on a presupposition that is not contained in the idea of
the fundamental being, and he himself has stated in one of his letters that it is
experience he builds on here. To the duality within his system there thus corresponds a
duality in his method: the fundamental being and the individual beings are related to
each other as speculation and experience. Spinoza lays speculation at the foundation and
uses experience only where speculation does not suffice. Here he stands in clear
opposition to present-day philosophizing, which as a rule lays experience at the
foundation and lets speculation build its hypotheses on the basis of empirical
knowledge. This is all the more natural a procedure since only in this way does it
become possible to check the very \opage{87}principles from which speculation starts.
But it shows Spinoza's profundity that even a method so altered leads thought to a
world-view akin to his. Thus workmen who dig a tunnel through a mountain, each from his
own side, will meet in the middle, provided only that they go the right way, each from
his own starting point.

Jakobi regarded Spinoza as the only consistent philosopher and gave the choice between
his doctrine and a \textit{salto mortale} into faith. Later criticism discovered, to the
joy of many, the double current and to that extent the inconsistency in Spinoza. The
iron ring seemed broken. But it must be well noted that the double current does not
first appear at a certain point inside the system, but is present in the first design,
in the definitions and axioms with which it begins. When Spinoza begins not only with
the definition of ``cause of itself'' and of ``fundamental being,'' but in the same
breath defines the ``individual being,'' and when his first axiom is: everything that
is, is \emph{either} in itself \emph{or} in another,---then he in reality begins not
with one but with two presuppositions, draws from the very first on two different
sources. And the duality is in reality only on the surface; the two currents obey one
and the same fundamental law: it is the idea of lawful connection that governs
everything in Spinoza, and that attests a deeper unity behind the duality that
criticism can point out.

The proposition that every determination is a negation \textit{(omnis determinatio est
negatio)} has often, and rightly, been cited as characteristic of Spinoza's standpoint.
\opage{88}For him the positive was what exists in itself, which is raised above all
relations and all limits. But this conception is struck by Goethe's well-known words:
``In der Begrenzung zeigt sich erst der Meister.'' It is precisely our task to win
determinations, to grasp nuances, to determine relations. We find the positivity of
thought in the rich, manifold and articulated content it is able to win. Only through
determination does the content become positive. Every determination, we might say, is a
relation; each time we think something, we determine its relations and only thereby
come to know it \textit{(omnis determinatio est relatio)}. Our knowledge is relative,
and only thereby is it also positive. One will also easily see that the absolute and
infinite, what exists in itself, can be thought and designated only in a negative way,
by negation of the finite determinations. What exists in \emph{itself} we cannot explain
otherwise than as that which does \emph{not} need \emph{anything} \emph{else} in order
to exist. What Spinoza regarded as the true, the only positive, is thus precisely
negative, is a limiting concept that is thought only by means of the determinations
that are negated in it.

\section*{II.}
\addcontentsline{toc}{section}{2. Mind and Matter}
\opage{89}The distinctive and consistent carrying through of the concept of substance is
the first and most epoch-making feature of Spinoza's philosophy. But closely connected
with this is another main point. If one fundamental being reveals itself in everything,
the relation between mind and matter cannot be such as it is commonly imagined to be.
Spinoza holds, like Descartes, that the difference between the inner world of thoughts
and extension in space is for us insurmountable; thought and extension are for him two
fundamental forms in which the fundamental being reveals itself, wholly and fully in
each of them. From the definition of the fundamental form it then follows directly how
the relation between mind and matter must be. By fundamental form or attribute Spinoza
understands, after all, what thought asserts of substance as making up its essence.
But as an expression of substance or the fundamental being, every fundamental form must
be conceived in and through itself; the fundamental form is, after all, only the
fundamental being itself, seen from a certain side. The one fundamental form therefore
cannot have been produced by the other: mind not by matter, matter not by mind; all
fundamental forms are original and simultaneous. One and the same fundamental being
reveals itself at once under different forms, of which the one does not exclude the
\opage{90}other. Each particular mode of expression of a fundamental form must therefore
be understood and explained through other modes of expression of the same fundamental
form, and not through determinations that belong under another fundamental form:
thoughts can be explained only by thoughts, material motions only by material motions.
The cause must everywhere be of the same kind as the effect. But since it is one and
the same fundamental being that reveals itself in the two fundamental forms, there will
be a parallelism present: the order and connection of thoughts is the same as the order
and connection of things. God's power of thought and his power of action correspond to
each other, and everything that follows from God's infinite nature within the world of
thought follows in the same order and the same connection within the world of matter.
Whether, therefore, we conceive nature as matter or as mind, it is one and the same
causal series we have before us. And if other fundamental forms were accessible to us,
the case would be the same with them; each by itself would be an expression of the one
fundamental being, coherent in itself and revealed in an infinite series of phenomena.

It is characteristic of Spinoza's standpoint and procedure that he begins by
determining the relation between mind and matter in its generality, and only then
passes on to the relation between soul and body. This relation is a particular example
of that general relation.---Conceived under the fundamental form of mind, the
fundamental being manifests itself in a world of thoughts. The human mind is one of
these infinitely many thoughts. When we say, then, that the human mind perceives this or
that, it means that the fundamental being, not insofar as it is \opage{91}infinite, but
insofar as it finds its expression in the nature of the human mind, or makes up
(constitutes) the essence of the human mind, has this or that thought. To each
particular thought there corresponds, within the fundamental form of extension, a
material motion. The various consciousnesses and the various thoughts differ from one
another in the same degree as the corresponding material motions and bodies differ from
one another. The more impressions the body is capable of receiving, the more ideas the
consciousness is also fitted to perceive; and the more independent and active the body
is, the more clearly and distinctly the mind is able to think. Since the fundamental
forms are each by itself infinite, this must hold for everything actual: everything in
nature is thus animate, although in different degrees. Every existing being
has---as a member of the fundamental being---a double side from which it can be seen,
is both mental and material.---Since it here becomes of interest to know the different
nature of bodies, Spinoza inserts some propositions from natural philosophy.

Bodies differ from one another only by the different ratio between motion and rest.
Every body is determined to motion or rest by another body, this by a third, and so on.
The way in which a body is determined depends partly on its own nature, partly on the
nature of the body that affects it. One and the same body is moved differently
according to the difference of the bodies affecting it, and different bodies are moved
in different ways by one and the same body.

From the simplest, most elementary bodies there are formed \opage{92}composite bodies.
For when bodies of the same or different size are held together by other bodies in such
a way that they fit together with one another, or when they communicate their motions to
one another in a certain determinate ratio, they are said to be united with one another
and all together form one body or one individual, which is distinguished from other
individuals by the determinate union among the particular bodies of which it consists.
Such an individual remains the same even if some parts are separated off and others
taken in, and even if the size of the parts and the direction of the motions are
altered, provided only that the whole form of the composition remains the same.

Several such particular individuals can now be thought of as united into an individual
of the second order; several of this latter kind again into an individual of the third
order, and so on. Finally we then conceive the whole of nature as a single great
individual, whose particular parts vary infinitely without the whole individual
undergoing any change.

The way in which Spinoza here develops the concept of the individual is entirely in the
spirit of modern natural science. Instead of seeking inner, mystical unities, he seeks
the law that holds for the coming into being of the individual. Only thus can the
concept of individuality be maintained in the face of the difficulty that experience
shows us no absolute and sharply marked boundaries between particular beings. It is a
necessary consequence of modern natural inquiry that one must either give up the concept
of individuality altogether or else make it relative, \opage{93}use it in a wider
range, so that we, like Spinoza, get individuals of different degrees.

In teaching an ascending series of degrees of individuality in nature, Spinoza has at
the same time expressed the law of development in one of its most characteristic forms:
the transition from lesser to greater coherence. His account is a forerunner of Herbert
Spencer's \textit{First Principles.}\footnote{See my work ``Den engelske Filosofi i vor
Tid.'' Kbhvn. 1874. P.~136 ff.}

If we now recall that to every material form, every bodily individual, there
corresponds an ideal form, an individual thought, since everything is animate, although
in different degrees,\footnote{One would entirely misunderstand Spinoza's doctrine of the
universal animation of nature if one supposed (as, e.g., \emph{John Toland} does in the
essay mentioned earlier) that it is meant as an attempt to explain how motion arises in
matter. Spinoza denies, after all, in the most definite way that thought and matter can
act on each other. His line of thought is rather this: that matter is homogeneous in all
its forms and compositions; when, therefore, at a certain point we find phenomena of
consciousness connected with material motions, then all material motions, although in
very different degrees, must be accompanied by such. (\textit{Et. II, 13. Schol.}) This
theory arose in Spinoza through the resolute way in which he carries through the unity
and coherence of the material world. It still stands as a natural hypothesis, and has
in our time been put forward not only by philosophers but also by natural scientists
such as Zøllner, Haeckel and others.}---then Spinoza appears as a precursor of
Leibnitz's doctrine of monads. Just as the human body consists of many composite
individuals of different nature, so the thought that makes up the nature of the human
mind is also composed of many thoughts. \opage{94}The whole of nature consists of an
infinity of individual beings, each with its degree of representative power. And it is
precisely such beings that Leibnitz calls monads.\footnote{From another side---as
pointed out earlier---Spinoza points forward to Leibnitz by the doctrine of the
infinitely many fundamental forms of the one fundamental being.}

On this point too Spinoza's view is decidedly monism. Just as little as he could grant
a dualistic relation between the fundamental being and the individual beings, could he
grant one between mental and material phenomena. Everywhere he seeks the unity that
lies behind the differences. And this striving of his stands out especially strongly
when one compares him with his predecessor in philosophy, Descartes. In the Cartesian
philosophy soul and body were two entirely different substances, which could really be
connected with each other only by a miracle. Although Descartes, who had cultivated
anatomy with zeal, had a conception of the significance of the nervous system so clear
and apt for his time that modern neurophysiology has its predecessor in
him,\footnote{He was, according to \emph{Huxley}, to the nervous system what Harvey was
to the circulation of the blood.} it was nevertheless in a purely external way that he
thought of the soul as connected with the body. In the so-called ``pineal gland'' in the
brain he saw the place\footnote{That the ``pineal gland'' has no such significance
\emph{Steno} had already shown in a treatise of 1669.} where the soul acts on the
bodily and in turn is affected by it. Spinoza remarks in his critique of this doctrine
that the mental and the material have no common measure, cannot \opage{95}be compared
with each other and therefore cannot determine each other either. Neither can the body
determine the soul to think, nor the soul the body to motion or rest or anything else
whatever. What determines the soul to thinking must itself be a thought, and what sets
the body in motion or rest must itself be a body. Soul and body are one and the same
thing, which is considered now under the fundamental form of thought, now under that of
extension. We do not have---as in Descartes---first a bodily motion that sets the soul
in motion, and then a mental activity; but since the bodily motion in itself (in the
fundamental being) is one with the mental act, the difference in time falls away.

Spinoza feels that the deductive proof (from the concept of the fundamental form) will
not be enough here. ``People are,'' he says, ``firmly convinced that the body moves or
rests at the mere nod of the soul, and that it performs much that depends on the soul's
will and ingenuity. Therefore people will hardly consider the matter calmly unless I
draw proofs from experience.'' He then remarks first that no one can know how far the
power of the body reaches. No one has yet found the limit of what can be brought about
according to the laws of material nature. It is thus quite arbitrary to set definite
bounds here, beyond which nothing is supposed to be possible except through the soul's
intervention.---But what does one really mean when one speaks of the soul setting the
body in motion? No one can form a clear and distinct concept of it. To say that the soul
brings about a bodily motion is thus the same as to \opage{96}say that one does not
know how this motion comes about. Here, as so often, people dress up their ignorance in
fine-sounding words.---In the instinctive actions of animals, in the behavior of
sleepwalkers and somnambulists, we have cases where bodily nature acts alone, and where
it accomplishes things that everyone admires.---Further, experience shows that the
liveliness and activity of the soul stand in a definite relation to the bodily
condition. Things would certainly stand far better if men had the ability to govern
their bodies. But they do not even have their tongues in their power. And yet the
chatterer believes he acts entirely freely, however often he comes to regret his
openness of mouth. Since man does not know the causes that determine his acting, he
regards himself as absolutely free, just as the stone, if it could think, would suppose
that it fell to the earth of its own free will.---We cannot, after all, act according
to our decision unless we remember what it amounts to. But we are not always masters
of memory. Everyone must grant that we cannot at every moment call up before our
thought whatever we like. Thoughts and decisions arise in consciousness with the same
necessity as the representations of external objects. And the necessary law that holds
for consciousness corresponds to the one that holds for the corresponding bodily
process: ``The decision or desire of the mind and the determination of the body are by
their nature simultaneous, or rather, they are one and the same thing, which we call
decision when it is considered under the attribute of thought and explained through it,
and bodily determination when \opage{97}it is considered under the attribute of extension
and derived from the laws of motion and rest.''

This whole doctrine belongs to what secures Spinoza's writings a lasting significance.
He has here in a brilliant way anticipated an insight whose empirical working out our
own time has only now reached. He had some ground to build on, namely Descartes's
mechanical neurophysiology. But here, as on several points, Descartes stopped halfway.
Spinoza's eye for unity and connection led him to the conviction that the series of
bodily causes cannot be thought of as broken at any point. If it now nevertheless
remains firm that there is a (for us) insurmountable difference between consciousness
and matter (or in Spinoza's language: thought and extension), then the only possible
conception becomes that of thinking of these two forms of existence as simultaneous and
parallel expressions of one and the same fundamental being. The physiology and
psychology of our own time also seem to lead to a similar conception.\footnote{This I
hope in years to come to be able to demonstrate in a separate work.}

For Spinoza as a speculative philosopher the matter stood thus: that the fundamental
being, which lies behind mind and matter, is what is clearest and most certain for
thought. For us it makes the head swim at this height. We rejoice to see the threads
draw closer to one another the further we are able to follow them, to see even the
oppositions that are insurmountable for us within the circle of our experiences point
back to a deeper unity; but this unity itself we are not able \opage{98}to bring forth;
each time thought tries to hold it fast, it turns out that we are on the point of
leaving the proper domain of our knowledge. From the trackless regions of speculative
hypotheses we then seek our way back to what Kant has called ``the fruitful depth of
experience.'' We are therefore not surprised that Spinoza's striving to build philosophy
on the idea of the absolute being leads us to such difficulties as how the fundamental
being, its unity undiminished, can appear in a plurality of fundamental forms, and what
significance these have for it and for one another. The riddle of life and existence
meets us ever anew at the end of all the attempts the human mind has made, in its
religious and speculative systems, to solve it.

Spinoza is the first to have developed a theory that helps us beyond the one-sidedness
of spiritualism and of materialism, not by compromising, but by really thinking the
matter through to the end and courageously drawing all the consequences. But just as
his religious theory was not understood, and his pantheism was interpreted as atheism,
so his psychophysical theory was not understood, and his monism was interpreted as
materialism. And even if one does not go so far, it is still often pointed out as an
inconsistency in Spinoza that he sees the matter rather much from the bodily side, when,
for example, he teaches that the perfection of consciousness depends on the perfection
of the corresponding body. But it lies in the nature of the case that we are obliged to
begin with the bodily when we wish to compare the different \opage{99}stages of mental
life. With the exception of our own consciousness we know mental life only through the
bodily as intermediate link, and find it everywhere bound up with this. And what is the
main thing: our knowledge of the bodily can reach a perfection and exactness that the
knowledge of the mental will never reach. One has the surest starting point, therefore,
when one begins with the bodily. That on Spinoza's conception the matter could
nevertheless just as well be seen from the other side is shown by a proposition such as
this one (Et. V, 1): ``Just as the thoughts and representations are ordered and
connected in consciousness, exactly so are the states of the body ordered and connected
in the body.'' No connection of thoughts can take place without a connection of bodily
motions taking place at the same time (the modern expression would be: molecular motions
in the brain). Therefore it is no breach of the theory either when Spinoza shows how man
can attain mastery over his drives and spiritual freedom. To the higher and the lower
within the life of consciousness there naturally correspond different combinations of
bodily processes. Modern physiology shows us, in the relation between higher and lower
nerve centers, in the way in which the former can regulate the latter, a parallel to the
psychological-ethical opposition between thought and drive (``the spirit'' and ``the
flesh''). The objection mentioned would not so often have been raised against Spinoza if
people had remembered his doctrine of the fundamental forms as each by itself infinite,
so that there are bodily forms that correspond to the highest mental functions, as well
as \opage{100}mental forms that correspond to the simplest bodily motions, and if they
had at the same time considered the matter from a physiological point of view. Spinoza
refers expressly to involuntary actions (instinctive actions, somnambulism, etc.) and
finds in them, like modern psychology, a hint toward the understanding of voluntary
motions. It is the idea of connection, of continuity, that here too leads his thought
to a happy stroke.

Before we leave this point in Spinoza's doctrine, it must be brought out that he does
not always express himself clearly and consistently here. As an example of a certain
material form and the corresponding thought being one and the same, he cites in one
place (\textit{Et. II, 7. Schol.}) the circle existing in nature and the concept of the
circle. This example is also constantly cited in expositions of Spinoza's philosophy.
But I cannot believe otherwise than that Spinoza has here been guilty of a momentary
confusion and lack of clarity. The relation between the actual circle and the concept
of the circle is a relation between object and subject, not between matter and
consciousness. The form within the attribute of extension that corresponds to a certain
form within the attribute of thought is not that which is thought and represented, but
the material motion (brain process) that coincides with the thinking activity. When we
represent an object to ourselves, this happens, says Spinoza, through another body
affecting our body and producing a certain state in it: our representation of the
external object is then the form of thought that corresponds to the state into which
our \opage{101}body comes through being affected by the other body. Our consciousness
perceives no external body as actually existing except through the representations of
the states of our own body, and therefore the states of our body can be called images of
things, although they do not reproduce their figure. To my representation of the table
there corresponds within the attribute of extension---not the table, but---the state into
which my body (in particular the brain) is put by being affected by the table. To the
concept of the circle there corresponds in the same way the material motion (in the
brain) that takes place when I think the circle.---What occasioned the skew in the
example cited from Spinoza is no doubt his strange usage in calling the human body the
\emph{object} of human consciousness. The bodily process that corresponds to the process
of consciousness has, as the history of anatomy and physiology shows, only very late and
by many detours begun to become an object for consciousness, and immediate consciousness
perhaps has no inkling at all that such a material process is necessary. Although the two
problems---of the relation between subject and object and of the relation between
consciousness and matter---touch each other and pass over into each other, it is
nevertheless important not to confuse them. The speculative philosophy in Germany, which
at the beginning of this century renewed Spinozism, undertook precisely the essential
change of putting subject and object throughout in place of consciousness and matter,
and precisely thereby acquired an abstract, logical-idealistic character, \opage{102}which
opened the way for many fantasies. Schelling declared Spinoza's philosophy to be
``one-sided realism''; but what we miss in Schelling himself and in his whole school is
precisely Spinoza's brilliant eye for, and deep respect for, the real, lawful connection
of things. % the print: „Tingene reale, lovmæsssige“ (sense: Tingenes reale, lovmæssige)

\section*{III.}
\addcontentsline{toc}{section}{3. Theory of Knowledge}
\opage{103}To every thought there corresponds, according to Spinoza's doctrine, a bodily
motion. But the motions in our body do not arise of themselves; they are called forth by
motions proceeding from other bodies. Our body cannot exist at all unless it is
constantly regenerated by influences from other bodies. From this it further follows
that the state in which our body finds itself at any given moment is determined not
merely by its own nature, but by the nature of the other bodies that affect it. The
thoughts that correspond to the various bodily states and motions are therefore
expressions not merely of the nature of our own body, but also of that of several other
bodies. Thus arise what we call sensuous representations \textit{(imaginationes)}. We
suppose we have in them images of things; but they correspond only to the way in which
our body is affected by things. We therefore often imagine that we have sense
perceptions of things that are not present to us, since our body can be brought by
inner causes into the same state as by outer causes. The mistake we thereby make springs
from the fact that no other representation arises in us that excludes the presence of
the object supposedly perceived. Error is thus something negative and consists in the
lack of other representations \opage{104}that could convince us that the object is not
present. One who does not know the nature of the horse may very well believe that he
sees a winged horse, and have no doubt about its reality.\footnote{This theory agrees
entirely with modern psychology. Taine has given it epigrammatic expression in the
definition: la sensation est une hallucination vraie. % the print: „hallucition“ (Rettelser: l. hallucination)
Yet this way of seeing had already been put forward earlier by the able but little-known
inquirer \emph{Louis Peisse}.}

When two or more bodies have affected our body at the same time, consciousness, when it
later perceives one of them, will also come to think of the other. On this rests memory,
a connection or linking among representations, which forms itself in consciousness in
accordance with the order and connection of the bodily states. In this connection we
have no knowledge of the nature of things; it is of a purely subjective nature. That is
why it will also be different representations that are connected in this way in
different men. When a soldier sees hoofprints in the sand, he thinks at once of a rider,
then of war, etc.; in a peasant, on the other hand, the representation of the horse
leads to that of a plow, a field, etc. The connections depend on how each is accustomed
to see things follow one another.---This law, the law of the association of ideas,
Spinoza calls in one place a mental law of nature, a counterpart to the physical laws of
motion. There are thus laws in the mental as well as in the bodily domain. The doctrine
of the associations of ideas had shortly before Spinoza been developed in clear and
concise form by the Englishman \opage{105}Thomas Hobbes; but Spinoza is one of the older
authors who have treated this important psychological doctrine in the most acute way,
and who have seen its great and far-reaching significance, as will appear when we come
to his doctrine of the passions. If Spinoza had been better known in England, he would
surely have been named as one of the most important forerunners of modern English
psychology.

As long as the sensuous representations and the associations of ideas % the print: „Ideassosiationerne“
prevail, we behave passively and reach no real knowledge. No representation and no
association of ideas can be explained from our own nature alone. They always correspond,
after all, to the way in which our body is or has been affected by other bodies, and can
therefore be understood only when one takes into account all these bodies and their
interaction. If we keep to the human mind alone, sense perceptions and associations of
ideas are consequences without premises. The human being is only one link in an
infinite series and cannot be understood unless one takes into account the other links
of the series. Or, to see the matter from its highest point of view: a perfect
(adequate) knowledge of external bodies is not found in God insofar as he is considered
as expressed in the human mind alone, but only insofar as he is at the same time
considered as expressed in many other beings.

As long as we remain at this stage of knowledge, we reach only a confused and incomplete
knowledge, depending on the accidental way in which things affect us. This first kind
of knowledge, which Spinoza \opage{106}designates as opinion or sensuous representation
(imagination), has no other sources than what one hears from others or happens to
experience \textit{(experientia vaga)}. The essence of things constantly escapes us here;
it is only their outer forms and relations that we can perceive. But the outer forms
and relations of things depend on their inner essence, and so the task is to find a kind
of knowledge that can give us access to this.

It is characteristic of Spinoza that he leaves the method of experience in such great
haste and does not examine whether the defects from which it seems at first glance to
suffer might not be remedied. Its imperfection consists, according to Spinoza, in its
seeing things from the standpoint of the particular being, which is nevertheless only
one link in the series. But it has been precisely the whole endeavor of modern
experimental inquiry to get beyond this isolated and accidental element in the starting
point, and with what success it has striven for this, the history of natural science
bears brilliant witness. Bacon already strove to formulate a method that reached beyond
\textit{experientia vaga}. That one can reach general laws by the path of induction
Spinoza himself has furnished proof; for it is precisely by this path that he came to
set up his doctrine of the degrees of individuality in nature, his theory of knowledge,
his psychology, his politics. That in his theory of method he stands so far from modern
inquiry has its ground in his speculative starting point, which in turn is an expression
of his yearning to grasp the absolute, what is in itself, which cannot be given in any
experience.---On the other hand it might seem as if precisely Spinoza's fundamental
\opage{107}view should have led him to lay a weight of principle on the method of
experience and on the apprehension of the particular and finite. When dualism is
rejected, and everything in nature is seen as a form of expression of one fundamental
being, then knowledge of the reality given in experience also becomes a knowledge of
the infinite itself, of the fundamental being. Spinoza therefore also says in one
place: ``Knowledge of the effect is one with a more perfect knowledge of the cause. We
therefore cannot gain insight into nature without at the same time extending our
knowledge of the first cause.'' But Spinoza, like Descartes, conceives the right method
as the one that starts from the highest cause. Acquaintance with what is given in
experience acquires only the significance of directing our thought toward certain
definite things. Knowledge of the essence of things itself must be drawn from another
source.

The second kind of knowledge is raised above the isolated and subjective standpoint
that was the imperfection of sensuous representation (imagination). All bodies have
something in common: they all belong under the same attribute, and the general laws of
motion hold for them all. Of this common element there can be a perfect knowledge; for
it does not merely express the state or nature of the particular individual, but the
knowledge of it is universally valid. By grasping the general and lawful connection of
things man raises himself above the accidentalness and externality of the imagination.
This knowledge does not concern something that holds only for a certain definite point
in time, any more than for a particular thing, but something that holds without regard
to time. To see things in their lawful connection is \opage{108}thus to see them under
the aspect of eternity. The necessity of things is one with the necessity in God's own
nature: God, after all, produces everything in producing himself. Things here lose
their outer independence of one another; we know that the power by which each of them
persists in its existence follows from the natural necessity in God's essence. We can
thus consider things as actual in two different ways: either so that we consider them
as existing at a certain point in time and at a certain place in space---which is the
way of looking of the imagination, the first kind of knowledge---or so that we see them
as contained in God and following from the necessity of the divine nature---and this is
the second kind of knowledge, which Spinoza calls reason (ratio).

The more man considers things as accidental, the more he is himself subject to
accident. What we learn by sense perception and by tradition is not in our own control.
On the other hand, the clear and distinct concepts we form depend solely on our nature
and its fixed and determinate laws. With regard to these concepts we need fetch no
confirmation and corroboration from outside: truth and certainty coincide. With the
true idea there immediately comes the certainty of its truth. Just as light helps us to
distinguish between light and darkness, so it is truth by whose help we discover what
is untrue \textit{(verum est index sui et falsi)}. Doubt, untruth and fancy arise only
through our considering things in a confused way and not forming distinct and clear
concepts by going back to the simple elements and the general laws. The idea of
\opage{109}a square circle is confused; for as soon as one goes to its elements, the
concept of circle and the concept of square, one sees at once that they cannot be
combined in this way. And in any case a false and invented idea will soon show itself
in its untruth if we only carefully draw all its consequences.

To grasp rightly what Spinoza means by rational knowledge, it is necessary to
distinguish between the common and the abstract. Rational knowledge has as its object
the general laws of things and thereby the essence of things. The laws that spring
from the essence of things hold for the whole as well as for each particular part and
are thus expressions of the unity in everything. They correspond to the concept of
attributes or fundamental forms, which are, after all, each by itself one with substance
or the fundamental being; each time we consider a body according to the general law of
extension, or a thought according to the general law of thinking, we conceive them as
expressions and forms of the fundamental being.---Abstract concepts, on the other hand,
are formed, according to Spinoza, by the imagination or the sensuous mode of
representation and depend on a confused apprehension of a multitude of particulars.
Here belong all general concepts, which are really schematic images that arise in us
when we have often had representations of bodily objects (horses, men, etc.) and these
representations now, as it were, grind off their differences, so that there remains only
that in which they all agree. Each man forms these general concepts in his own way,
according to the features of the objects that have particularly made an impression on
him.---The difference between the abstract and the common thus \opage{110}consists in
this, that the former gives an indistinct schema, the latter a clear and distinct
knowledge of law and, in the knowledge of law, a knowledge of what works in everything.
In true knowledge our thoughts follow the eternal order of things; our mind is an
\textit{automaton spirituale}, whose working follows the same laws as the divine nature.

Spinoza's model in describing this form of knowledge is the mathematical-mechanical
conception of nature. He would have set up a closer relation between the first and the
second forms of knowledge if he had had greater sympathy for, and clearer understanding
of, the experimental method. For the path by which one arrives at the application of the
mathematical method in natural science is exact experimentation, which makes
quantitative determinations possible.---Rational knowledge concerns, after all, laws and
relations generally; but the elements, the actual existences, between which these
relations obtain and for which the laws hold, we can know only by experience. As shown
earlier, the very concepts of attribute and substance have sprung from analysis of what
is given in experience.---Here again we see that there is not, as Spinoza thought, any
complete opposition or any sharp boundary between empirical knowledge and rational
knowledge.

But rational knowledge still has for Spinoza % the print: „før“ (sense: for)
the unsatisfactory feature that it is discursive in character, draws conclusions from a
general law to the particular cases. To illustrate the relation between the kinds of
knowledge he uses the following comparison. When three numbers are given, and a fourth
is to be found that is related \opage{111}to the third as the second is related to the
first, the person who has merely learned the rule of three multiplies the second and
third numbers together and divides the product by the first. His knowledge is purely
empirical, rests on tradition and experience, and thus belongs under the first kind of
knowledge. The person who knows mathematics, on the other hand, performs the same
operation because it is an application of the doctrine of proportions. This corresponds
to the second kind of knowledge. But when the numbers are very simple, one need perform
no operation at all. One can, e.g., see immediately that 1 is related to 2 as 3 to 6.
The general rule here becomes superfluous; in place of calculation and reasoning there
appears immediate intuition of the relation; in place of the discursive, the intuitive.

The third and highest form of knowledge is intuitive knowledge. It beholds the essence
of things immediately, without the help of general propositions. As an example of such
intuitive knowledge Spinoza cites the fact that when we know something, we thereby also
have an immediate knowledge of what it is to know. In the first book of the ``Ethics''
he has, by his own account, employed the second form of knowledge. He there set up a
series of definitions and axioms and derived from them his doctrine that everything,
as regards both essence and existence, depends on the fundamental being. Only at the
end of the ``Ethics'' (in the fifth book) does this doctrine become the object of
intuitive knowledge, as the human mind immediately sees itself as an element in the
fundamental being. Here we no longer have general rules, but an \opage{112}immediate
beholding of the inner connection of the particular with the whole.

In the ``Short Treatise'' this highest form of knowledge is described in expressions
that have a certain mystical and religious stamp. We are thereby reminded of the
original practical motive of Spinoza's thinking. ``Rational inference is not what is
most excellent in us, but is only like a ladder by which we climb up to the desired
point, or like a good spirit that, without falsehood and deceit, brings us word of the
highest good, in order thereby to urge us to seek it and unite ourselves with it, which
union is our highest welfare and happiness.'' ``That this [third] kind of knowledge is
not a consequence of something else, but is immediate, is evident from the fact that God
is the cause of all knowledge and is known only through himself, not through anything
else, and from the fact that we are by nature so united with him that without him we
can neither exist nor be understood. From this close union between God and us it
follows that he can only be understood immediately.''---The way in which we thus
immediately apprehend the fundamental being or the Deity is designated partly as an
immediate revelation to the understanding, partly as a feeling and enjoyment of things
in themselves.---In his later writings too Spinoza designates the third form of
knowledge as the highest virtue and the highest blessedness.

There is this difficulty in Spinoza's description of the third form of knowledge, that
in some places he describes it as an altogether immediate intuition, but in other places
determines it as a movement \opage{113}of thought that goes from the knowledge of the
fundamental forms to the knowledge of the essence of things, or as a knowledge of the
thing through its proximate cause. With this latter determination it does not become
quite different from rational knowledge, above which it is nevertheless supposed to be
raised; it involuntarily shows its kinship with it, just as we saw above a natural
kinship between rational knowledge and empirical knowledge, which were not so foreign
to each other either as Spinoza supposed.

The motive for seeking beyond rational knowledge we shall easily find if we recall
Spinoza's fundamental concept and its definition: ``Substance is that which is in
itself and is conceived through itself, i.e., that whose concept does not need the
concept of any other thing from which it would have to be formed.'' If a concept is to
be apprehended in and through itself, without the use of other concepts, then it is
clear that no operation of thought or inference can take place. Where should the
movement of thought go when there is only one concept? What is understood ``in and
through itself'' is an object of beholding, not of thinking. Spinoza therefore could not
hold fast to his concept of substance without setting up a form of knowledge different
from ordinary experience and thinking and raised above them. The same holds of all the
philosophers who have maintained a direct and positive knowledge of the absolute. For
Plato the highest knowledge was ``a beholding of the best of what is,'' namely of ``the
first, presuppositionless ground of everything.'' Presuppositionless, after all,
Spinoza's substance too is supposed to be. But what is to be altogether without
presuppositions cannot be thought; it must be intuited. Schelling carried this to an
extreme in \opage{114}his doctrine of ``intellectual intuition,'' which is an immediate
apprehension of the absolute as the eternal unity of subject and object, of all
oppositions. This intuition was for Schelling a special organ, which one either had or
did not have. Those who lacked it simply lacked the condition for philosophizing, just
as one who could not construct and intuit figures in space would lack the condition for
being a mathematician.---Just as speculative philosophy must appeal to a special organ
when it wishes to maintain a positive knowledge of the absolute, so is it also the case
with positive religion and theology. These too point beyond experience and thinking to a
higher source of knowledge. Through rebirth and faith access is opened to a higher
nature than that which is accessible to the unbeliever.---We see here again how
speculative philosophy is reminiscent of the positive-religious consciousness. When
philosophy, on the other hand, acknowledges the law of relativity and carries it through
completely, it must give up the appeal to a special source of knowledge and place itself
on an equal footing with all other scientific inquiry. Science, like all culture,
originally unfolded under the protecting cover of religion, and philosophy still
displays many marks and remnants of this husk. It will no doubt also cast them off, and
it will be able to do so without injuring any real human interest.

This is not to say that the idea of an intuitive knowledge is to be rejected entirely.
The error consists only in regarding it as a special source of knowledge. The
development of knowledge does not reach its peak in \opage{115}mere reasoning. Everyone
knows that one has fully made a piece of knowledge one's own only when one can skip or
survey all, or at any rate many, of the intermediate determinations and see the whole
connection at a single glance. Intuition, or intellectual seeing, is thus a natural
result of effort and practice with scientific trains of thought. In this sense it holds
that intuitive knowledge alone is the crown of our labor of knowledge. But it then
stands in the closest connection with the other stages of knowledge and grows forth from
them in a natural way, just as good habits and instincts grow forth from the efforts of
the will. Spinoza's own example mentioned above points in this direction. For that 1 is
related to 2 as 3 to 6 is seen immediately only by one who is practiced in operating
with numbers; the child who does not yet know his multiplication table must first
laboriously work his way up to such an intuition.

---There is a class of concepts to which Spinoza denies absolute validity because they
suffer from the imperfections peculiar to the first kind of knowledge. They are such
concepts as good and evil, beautiful and ugly, order and disorder, purposiveness, etc.,
when they are applied to the essence of things. The imperfection of the imagination or
sensuous representation consisted, after all, in its considering everything from man's
individual and accidental standpoint and measuring everything by this. Man from the
outset judges everything practically, regards all things around him as means and tools,
and calls good what he can use, evil what harms and hinders him. He thus accustoms
himself first and foremost to ask, of \opage{116}every thing he sees, what it serves
for, instead of how it came into being. When he has then learned that the eyes are for
seeing with, plants and animals for eating, the sun for giving light, he has explanation
enough. From this stems the notion of the purposiveness of nature. Now on the other hand
there is much that is harmful and pernicious in nature; this, people then suppose, is
due to the wrath of the gods over men's sins. And when experience shows that misfortunes
strike good men just as well as bad, it is declared that the counsels of the gods
surpass human comprehension. Instead of remedying one's ignorance, one thus appeals to
it as an argument. If mathematics had not shown men another rule of truth, the human
race would have remained in its ignorance. But the old prejudices are still so strong
that one who seeks scientific insight into nature, instead of persisting in foolish
wonder, is regarded as impious.---Just as men ascribe ends to nature according to their
own wishes, so they also find order and disorder, beauty and ugliness in nature,
according as things are easy or difficult to intuit and apprehend, appeal to their
power of apprehension or not. But these concepts, and in general all concepts of
perfection and imperfection, spring from comparisons that we make, and assert nothing
about nature in itself. We form a general concept, a certain type, that fits all
individuals of a certain kind, and the more the particular individual approaches this
type, the more perfect we say it is. But this general concept, this type, belongs to our
\opage{117}understanding, and not to nature. The infinite being that we call God or
nature acts with the same necessity with which it exists. Our concepts of perfection and
imperfection therefore cannot be applied to it. Only when one understands by perfection
reality, being, can this concept be applied to nature; perfection then extends just as
far as the true being of things. But limitation belongs only to the finite, not to the
infinite being. When we find imperfection in nature, it is thus only because we
consider it from a finite point of view: in itself everything is perfect, when one does
not set it over against something else with which one compares it, i.e., when one does
not consider it as finite. Nature does not act according to the law of human reason and
does not direct itself according to human purposes and wishes, but according to
infinitely many other laws that spring from the eternal order of nature, in which man is
only a single link. When something therefore seems to us ridiculous, absurd or evil, it
comes from our limited knowledge. What our reason declares evil is evil only according
to the laws of \emph{our} nature, not according to the order and law of \emph{the whole
of} nature. One who strives for knowledge of the true essence of things will not blame
or mock anything in nature, because he knows that everything happens according to the
eternal and necessary laws of the divine nature.

If aesthetic and ethical concepts cannot be applied to nature in itself, because they
arise only through comparisons and valuations that we make from our finite standpoint,
then the same would seem to have to hold \opage{118}for still more concepts, and Spinoza
saw this too. As we have already seen, time, number and measure are determinations that
do not hold for true being. Time is not a determination of things themselves, but merely
a mode of representation in us by which we determine the duration of a thing's
existence. But in God there is nothing earlier or later; his existence---and thus all
existence---is a necessary or eternal consequence of his essence.

This way of looking at things ought surely to have been applied to all other concepts
as well. Spinoza would have done so had he not, as a good dogmatist, believed in the
possibility of our being able to come to consider things from another point of view than
our own. Such concepts as substance and cause, which Spinoza applies without hesitation
to the absolute, will on closer consideration also prove to be relative. Yet it is
remarkable to see to how high a degree Spinoza has succeeded in emancipating himself from
anthropomorphism. The naive view applies all determinations and concepts to things
without further ado. Modern philosophy has in its successive development been led to
refer all these concepts back to human consciousness and to see them as symbols that
express existence from a definite point of view, without our having in them the only
possible solution of the great riddle. Spinoza has often been sharply reproached for his
doctrine that concepts such as good and evil do not hold for the essence of things and
have only relative validity. In this, however, he is quite within his rights. His
treatment of aesthetic and ethical concepts is a necessary \opage{119}consequence of
Copernicus's astronomy and of Galilæi's and Descartes's physics. When the earth is no
longer the center of the universe, man cannot be so either, and his judgments hold only
from a point of view that is vanishingly small in relation to the great, infinite whole.
And when everything in nature is subject to the laws of mathematical mechanics, one
cannot explain its phenomena by inserting ideal forces. What must be objected to Spinoza
is, on the contrary, that he has not gone far enough, and further, that he in part
assumes that the relativity of concepts is equivalent to their invalidity. That this is
not so he in reality grants himself, as regards the ethical concepts, for in his ethical
doctrine he sets up ``an ideal pattern of human nature'' and calls good what brings us
nearer this pattern, evil what removes us from it. If such a type really expresses
human life as it would be, according to its natural laws, if it reached its highest
possible development, then it and the other ethical concepts attached to it have all the
validity that can be demanded, and this validity suffers no impairment from our not
being able to decide what significance this ideal and our striving for its realization
have for the universe as a whole.\footnote{Compare my work ``Om Grundlaget for den humane
Etik.'' Kbhvn. 1876, P.~84--86.}

\section*{IV.}
\addcontentsline{toc}{section}{4. Psychology and Ethics}
\opage{120}It is a quite common prejudice that an ethics is altogether impossible on the
basis of determinism. When a philosopher conceives human feelings and actions as subject
to laws of nature, a certain kind of criticism is at once ready with the remark that only
by a happy inconsistency can he set up an ethics, a doctrine of the ends of human action
and of the means of attaining them. Without absolutely free will, it is said, no ethics.
On closer consideration one will easily see on which side the ``happy inconsistency''
lies. Ethics would be impossible precisely if there were an absolutely free will in man;
then every coherent doctrine of human action would be impossible. The chain could be
broken at any moment. Of definite ends grounded in human nature, and of definite means
for their realization, there could then be no question. Spinoza was therefore, precisely
as a determinist, well fitted to set forth an ethical doctrine. In thorough psychological
knowledge of human nature and its laws there is hardly any ethicist who comes near him;
and if ethics is to be grounded in the laws of human nature, thorough psychology is the
first condition for an ethicist. \opage{121}It is another matter that ethics, like every
investigation of the development and laws of individual beings, acquires a significance
of principle in Spinoza's system only insofar as the experimental undercurrent is able to
assert itself against the speculative overcurrent that has its expression in Spinoza's
fundamental thought, the concept of substance. If one starts from this concept alone,
ethical development, like all concrete and finite development, becomes inexplicable and
illusory. But if we take both currents in Spinozism, as they already manifest themselves
in the first axiom in the first book of the Ethics (see above, p.~87), then Spinoza is
fully consistent in setting up an ethical doctrine.

Spinoza maintains repeatedly and with great energy that human feelings and passions are
to be investigated with the same calm and impartiality as other natural phenomena.
``Most of those,'' he says in the preface to the third book of the ``Ethics,'' ``who have
written about the passions and about men's way of life seem to speak not of natural
phenomena that follow the common laws of nature, but of things that stand outside
nature. They seem to regard man in nature as a state within the state. For they believe
that man disturbs the order of nature rather than follows it, and that he has absolute
power over his actions and cannot be determined from anywhere but himself. Nature is
everywhere the same and everywhere one; its power and efficacy, i.e., the laws and rules
of nature according to which everything happens and passes over from one form to
another, are everywhere and always the same, and there is therefore only one way in
which the nature of any thing whatever can \opage{122}be understood, namely through the
laws and universal rules of nature. The passions of hatred, anger and envy therefore
spring in themselves from the same necessity and power of nature as other finite
things; they have definite causes through which they are explained, and definite
properties that are just as worthy of our knowledge as the properties of any other thing
whose contemplation delights us. I shall therefore treat of the nature and powers of the
passions, and of the mind's power over them, by the same method I have used in what
precedes, and I shall consider human actions and drives in the same way as I consider
lines, planes and bodies.'' In the introduction to his Political Treatise Spinoza
likewise declares that he will treat these questions with the same freedom of mind that
we bring to mathematical investigations. ``I have earnestly endeavored not to mock,
lament or curse human actions, but to understand them, and I have therefore considered
the human passions, love, hatred, anger, envy, ambition, pity and the other emotions,
not as faults of human nature, but as properties that belong to it, just as heat, cold,
storm, thunder and the like belong to the nature of the air, which, to be sure, bring
inconveniences but are nevertheless necessary and have definite causes through which we
strive to understand their nature, an understanding that brings the mind just as much
joy as the knowledge of the things that bring pleasant consequences for us.''---These
words have sometimes been cited with a shudder and regarded as \opage{123}signs of a
cold and heartless way of thinking. To this it need only be said that if a psychology
at all---a knowledge of the laws of mental life---is to be possible, then one must
conceive mental phenomena as pure objects of knowledge, without any feeling other than
the wish to understand and explain them as well as possible. No inquiry is possible
without the inquirer's frame of mind, and this frame of mind is driven out as soon as
aesthetic and moral feelings assert themselves. Precisely because Spinoza saw this so
clearly, and because this inquirer's frame of mind was the prevailing one in him, he
has become one of the most important forerunners of modern scientific psychology.

Kuno Fischer has aptly shown how Spinoza here accomplished for philosophy what
Shakespeare before him had already accomplished in poetry: ``Every real character has
its passions, and every passion has its necessary course of development, which is
dramatic because it consists in definite actions; every true drama is nothing but the
natural history of human passions.---In philosophy Spinoza was the first who grasped
this truth and proved it with reasons; but the world-historical merit of having first of
all discovered dramatic human life in this sense belongs to an earlier mind. Beyond
philosophy, on the border between two ages, rises the gigantic Shakespeare, to whose
poetic intuition it first became clear that the events of the human world proceed from
actions, actions from passions, passions from characters, and characters from the
natural and \opage{124}historical connection of things.---Man acts in Shakespeare's
imagination and in Spinoza's philosophy as a nature that fulfills its law, and passion,
which moves human power and drives it forward, knows only one way, the way of the stream
to the fall.''---Kuno Fischer further points out that Shakespeare and Spinoza celebrated
their resurrection at the same time and among the same circle of men. It was Lessing in
particular who drew them both forth again and paved the way for their understanding.
(Geschichte der neuern Philosophie. I. Mannheim 1854. P.~416--420).\footnote{In the
preface to his first work \emph{Taine} refers to Spinoza, and this reference is
characteristic of the famous critic's whole standpoint and method.}

Spinoza's doctrine of the passions (the affects), which forms the groundwork of his
ethical doctrine, can be understood only if one starts from his doctrine of the
individual being. Each individual being is a finite form of the infinite fundamental
being; all individual beings stand in interaction with one another and together form
the great individual of nature. Insofar as the individual being is affected and
determined by the other individual beings, it is restricted, limited and passive. That
is why it is only in part the cause of the effects that proceed from it. A complete
(adequate) cause an individual being would be only if the effect could be explained
entirely from its nature alone. We are said to act when we are complete causes, but to
suffer when we are incomplete causes. By affect or passion we understand a state
(affection) by which our power of action \opage{125}is increased or diminished, and at
the same time the representation of this state. The passion is of an active character
when we ourselves are the complete cause of the state, but of a passive character when
we ourselves are only in part the cause.

---It is thus as a link in the series of natural beings that man is subject to passive
passions. He can no more than any other individual being withdraw himself from the
influence of external causes. The passions arise through the clash between the working
of external causes and the individual being's own fundamental drive. Every thing
strives, insofar as it has existence in itself, to go on existing. Only external causes
can abolish a thing's existence; % the print: „Existens“
insofar as it depends on the thing alone, it will go on existing. This drive of
self-preservation is in itself nothing other than human nature itself, from which there
follows with necessity \emph{that} which serves for its preservation, so that man is
naturally determined to seek it.---Spinoza seems here to have transferred the law of
inertia from the bodily to the mental domain. Galilæi, and after him Kepler and
Descartes, had shown that a body will remain in its state (of rest or motion) when the
influence of other bodies does not bring it out of it. This law Spinoza now regards as a
universal law, not restricted to physical bodies. In itself, too, it is only a form of
the universal principle of causality, % the print: „Kavsalitetsprincip“
which demands a definite cause for every change. For Spinoza this generalization lay
near at hand. All individual beings are finite forms of the fundamental being, and the
power of the fundamental being works in them. The individual being is, to be sure,
finite and limited, but as a form \opage{126}of the fundamental being it has a share in
its independence. ``The power of natural things to exist and to act can be none other
than the eternal power of God.'' ``Although each thing is determined by another
particular thing to exist in a certain way, yet the power by which each thing persists
in its existence follows from the eternal necessity of God's nature.''

The passion now arises through this fundamental drive being checked or furthered in its
working. When the mind, owing to external influences, passes over to a greater power of
thought, joy arises; when it passes over to a lesser power of thought, sorrow arises.
Desire, joy and sorrow are the three most original passions, from which all others
spring. The laws according to which they develop are the laws of the association of
ideas. When joy is accompanied by the representation of the cause that has called it
forth, love arises. When sorrow is accompanied by the representation of its cause,
hatred arises. The mind seeks as far as possible to hold fast what increases and
furthers its power of action, and turns away from what checks it. The one who loves
therefore necessarily strives to have the beloved object present and to preserve it,
whereas the one who hates seeks to remove and destroy the object of his
hatred.\footnote{Descartes had (in his Traité des passions) defined love as the will to
unite oneself with the beloved object in such a way that together with it one forms a
whole. Spinoza objects to this definition that it expresses, to be sure, an essential
property of love, but not the essence of love itself.---A similar remark is made by
Socrates in Plato's ``Symposium'' on the occasion of Aristophanes's mythical account of
how men in love seek their original other half.}

\opage{127}When we have had two passions at once, then, if the one of them arises again,
we shall also come to feel the other. In this way any thing can indirectly become a cause
of joy, sorrow or desire, namely when the feeling the thing itself calls forth awakens
another feeling, which we therefore often cannot explain ourselves, because we have
forgotten the original connection between the two feelings. Many inexplicable
sympathies and antipathies have this origin.---One and the same thing can call forth
mixed feelings, since it is directly the cause of one feeling, indirectly of another.

Because of the property of love that it binds the lover and the beloved object into a
whole, the feelings of the beloved object will also arise in the lover. Love thus
becomes like an extension of our being, through which new causes of feelings of various
kinds arise. What awakens joy or sorrow in the beloved object will call forth love or
hatred in the lover.---Sorrow brought about by the misfortune of others is pity, whereas
Spinoza missed a word in the language to express joy at the good fortune of others.
Love toward one who does good to another is called favor \textit{(favor)}; hatred toward
one who has inflicted evil on another, indignation \textit{(indignatio)}.---On the other
hand we feel love toward one who inflicts evil on someone we hate, and hatred toward one
who does him good.

When we imagine that someone loves, desires or hates something that we ourselves love,
desire or hate, we shall thereby come to love, desire or hate it \opage{128}still more,
whereas we shall feel a wavering of mind \textit{(animi fluctuatio)} when we imagine that
he abhors what we desire, or the reverse. Everyone will therefore by nature strive that
the others come to live according to his mind. And since all now strive for this, they
hinder one another in equal measure. Everyone wants to be praised or loved by all others,
and therefore all hate one another. Ambition therefore becomes the strongest of all
passions; it draws nourishment from many sides. It seeks the approval of men, but on the
other hand also wants all to approve what it itself approves.

The thought that a being resembling us, toward which we have no preconceived feeling,
has some feeling or other awakens the same feeling in us. For insofar as the other being
is like us, it comes to the same thing whether we imagine it in a certain state or we
imagine ourselves in the same state. We rejoice and grieve with others merely because
they are beings like us. Even the downfall of our enemy therefore awakens sorrow in us,
because he was a man like us. Here further belongs all involuntary imitation, as when one
draws back one's hand because one sees another burn himself. By the same law another's
joy in an object awakens our desire for it, and we seek to deprive him of the object. So
it goes with human nature, that we pity those who fare ill and envy those who fare well.
Compassion, envy and ambition spring from one and the same property of human nature.
This one will especially find confirmed when one \opage{129}observes man in the first
years of his life. The child cries and laughs according as it sees others cry and laugh,
and it feels desire for what it sees others rejoice in and desire. The role that envy
and ambition play in human life is nourished by upbringing, since parents especially
make use of these passions to spur children on to virtue.

In each particular case the passion acquires a different character according to the
particular object that awakens it, and according to the peculiarity of the individual in
whom it arises. One and the same individual is affected differently by different
objects; the same object affects different individuals differently, and one and the same
individual differently at different times. The passion is, after all, nothing but the
very nature of the individual being, and the essence of the particular, definite passion
depends on how it is affected by other individual beings, other links in the chain of
nature.

Through the combination and interaction of the various passions a multitude of derived
passions can arise. They form with the same necessity as the original passions. We are
set in motion in many ways by external causes and tossed about like the waves of the sea
driven by different winds, without knowing our goal and outcome.

---Spinoza's account of the origin and development of the passions is not only
characteristic of his intellectual direction and his standpoint, of his search for the
necessary connection in everything, but is also in itself a model of psychological
investigation and has great lasting value. In the psychological literature of the
Continent \opage{130}this account is still unsurpassed; England, on the other hand, has
in Hobbes, Hartley, Hume and in our time especially Bain a series of men who have
treated the same subject by the same method as Spinoza and can be placed beside him.

The basic determination of joy and sorrow (pleasure and pain) from which Spinoza starts
agrees entirely with that of modern psychology. Joy, he says, is\footnote{It would have
been more correct to say ``arises through'' instead of ``is.'' These simple psychical
phenomena can be defined only by what calls them forth.} transition to greater, sorrow to
lesser perfection; and he points out expressly that joy is not perfection itself, nor
sorrow imperfection itself. ``For if man were born with the perfection to which he
passes over, he would possess it without the feeling of joy. This shows itself still
more clearly in sorrow. That sorrow consists in the transition to lesser perfection, not
in the lesser perfection itself, no one can deny, since man cannot grieve insofar as he
partakes of any perfection.''---We have here the law of relativity in its application to
the domain of feelings. Pleasure and pain we feel only through a change in our state, an
advance or a decline in a certain respect. Yet they do not consist, as Spinoza likewise
expressly stresses, in a comparison between the two states that succeed each other. Such
a comparison would lead us beyond the domain of feeling to the purely intellectual.---We
shall see, however, that Spinoza's psychology \opage{131}suffers precisely from the
one-sidedness that this difference between the various expressions of mental life is not
held fast, since the intellectual side really becomes, in him, everything.

Pleasure and pain arise through advance and decline in our being, and our being Spinoza
sees as active, as striving for self-assertion---a finite form of the infinite
fundamental power working in everything, which is cause of itself. This is reminiscent
of the fundamental biological law that has been stated and carried through in our day in
so brilliant a way, and that sees life as a struggle for existence. Spinoza has
predecessors in this conception---not only in Descartes and the other physicists who set
up the law of inertia---but also in Thomas Hobbes, who built his whole ethical and
political doctrine on the drive of self-preservation.---

We have so far considered only the passive passions, which arise in man through
influences from outside and therefore cannot be explained entirely from his own nature
alone, since the nature of the external surroundings plays an essential part as well.
But passion is not always of a passive nature, although the Danish term
[\textit{Lidenskab}, from \textit{lide}, to suffer] takes account of that alone.\footnote{Spinoza's
general term is \textit{affectus}, and this can be either passio or actio.} Besides the
joy and desire that are of a passive kind, there are other feelings of joy and desire
that belong to us as acting beings.

Our power of action coincides for Spinoza with our power of thought. He sees the essence
of the mind in thought \opage{132}(just as he sees the essence of matter in extension);
the decision or act of will is for him only the assent to and affirmation of the truth
known. But thought is not a mute picture; it is activity and life. One can act against a
picture or a word that hovers before consciousness, but not against what has really
become the mind's own property and what it has assented to. Intellect and will are one.
We act insofar as we have adequate ideas, i.e., ideas that can be explained entirely from
our own nature and are therefore clear and distinct. We suffer insofar as we have obscure
and mutilated ideas, and this is the case when we are determined from outside. Our power
of action or of thought can be increased or diminished. Sorrow is a sign that it is
diminished, and therefore belongs to the passive feelings. Joy and desire, on the other
hand, \emph{can} be of an active nature.

All actions that belong to the mind as thinking spring from strength of mind
\textit{(fortitudo)}, which appears partly as courage \textit{(animositas)}, partly as
generosity \textit{(generositas)}. % the print: „gneerositas“ (Rettelser: l. generositas)
Courage is the desire by which each strives to preserve his being according to the
dictate of reason alone. By generosity Spinoza understands the desire by which each,
according to the dictate of reason, strives to help other men and to join them to
himself in friendship. Thus the actions that have the agent's own good in view belong
under courage; those that also have the good of others in view, under generosity. To the
first kind belong virtues such as temperance and presence of mind; to the latter kind
modesty, mildness, etc.

\opage{133}In the opposition between the suffering and the acting state, or between the
obscurity and the clarity of thought, we have the starting point for an evaluation of
the affects. Although the concepts good and evil do not designate anything positive in
things themselves, they nevertheless acquire a relative significance for us, since we
compare actual human nature with a pattern that we form of perfect human nature. Good we
then call what brings us nearer this pattern, evil what removes us from it. Virtue is
power---it is, namely, man's very nature, insofar as man is able to produce something
that can be conceived through the laws of his nature alone. Independence, spiritual
freedom, spiritual power is the essence of virtue, which, as said, appears partly as
courage, partly as generosity.

Since man's essence consists in thought, the dictate of reason cannot be different from
the fundamental drive itself. The source of passion was, after all, the individual
being's striving to assert itself. But precisely this striving is the first and only
foundation of virtue. No principle can be conceived that should lie deeper than this
striving, which is one with the very fundamental nature of the individual being. Each
strives for what is useful to him. This striving is checked only by external causes and
is in itself one with virtue. The dictate of reason is that each should seek his true
advantage. Reason is nothing but our mind itself, insofar as it has a clear and distinct
knowledge. Therefore the dictate of reason aims at nothing but knowledge, and striving
for true knowledge is the only virtue. Our power \opage{134}and our freedom consist in
true knowledge; we act, after all, only when we have adequate ideas; otherwise we
suffer. But the highest thing the mind grasps is God, the absolutely infinite being,
without which nothing can be or be conceived. To develop one's knowledge is to know God,
his attributes and his acting, which follow from his nature. Blessedness consists in the
peace of mind that springs from the intuitive knowledge of God. The highest desire of one
who is led by reason is the one that leads him to clear knowledge of himself and of all
things that present themselves to his understanding. In this striving human nature finds
its highest fulfillment. It develops here according to its own inner laws. No evil can
therefore befall man except from external causes, or insofar as man is a part of nature
and is affected by other parts. And this, as a finite being, he must necessarily be. But
when other individual beings agree with his nature, they will be able to be a help to
him in his striving. Nothing is so useful to us as what agrees entirely with our nature.
When two individuals of the same nature are united with each other, they form an
individual twice as great as the single individual. Man can wish for nothing better for
the preservation of his being than that all should so agree in everything that they
form, as it were, one mind and one body, and that all at once should seek the common
advantage of all. Those who are led by reason therefore wish nothing for themselves that
they do not also wish for other men.

Insofar as men are led by the passive passions, they differ from one another, and to
that extent each \opage{135}individual man is also changeable and inconstant. When, on
the other hand, they are led by reason, they agree. Each then does what is good for human
nature as a whole, and therefore also what agrees with the nature of each individual man.
By seeking his own true advantage each individual thus furthers the advantage of others.
The highest good, the knowledge of the eternal and the peace that follows from it, is
common to all; no one is excluded from access to it. And the more each for himself
becomes a partaker in this good, the more he will also wish to make others partakers in
it. Partly it is a psychological law that what we love gains still greater value for us
when we see that others love it too; partly we shall be furthered in our striving when
others strive with us. There is no way in which a man can better display what spiritual
powers he commands than by educating others to live according to the dictate of their
own reason. Moreover, only by united forces can men procure what they need and withstand
dangers. Admittedly very few live according to the dictate of reason, but it holds
nevertheless---despite all that satirists, melancholics and theologians object---that
far more goods than evils spring from the social life of men.

From the standpoint now reached Spinoza judges the various passions. Joy is in and of
itself good, since it is the expression of an advance in perfection. Sorrow, on the other
hand, is in and of itself evil, since it signifies a decline. Everything we do should aim
at advance and improvement. But as long as we grieve, we make ourselves unfit to act in
this way: sorrow \opage{136}is a decline in the power of action.---Spinoza therefore
evaluates many feelings in a way different from the usual one, and here in particular he
enters into a characteristic opposition to the theological conception. From his clear
insight into the necessary connection, and his faith that this necessary connection is
one with the essence of the Deity, there had to follow an evaluation of various moods
and feelings different from that of those who regard nature as corrupted by an original
fall. In a statement of a more personal character than one otherwise finds in his
writings, he says: ``Only a savage and gloomy superstition can forbid joy. Why should one
rather still hunger and thirst than drive away melancholy?---Thus I think, and thus the
matter stands for me: No deity, and in general no one who is not envious, rejoices in my
weakness and my misfortune, or regards tears, sobbing, fear and similar signs of
weakness of mind as signs of virtue. On the contrary, the greater the joy we feel, the
greater the perfection to which we pass over, and the more we share in the divine
nature. The wise man will therefore use things and take delight in them, thus
strengthening and refreshing himself with moderate and pleasant food and drink, with
fragrances, beautiful green plants, fine dress, music, bodily exercise, the theater and
other such things, which anyone can enjoy without harm to anyone else.''

That hatred is evil follows of itself. It falls away with the knowledge that the
endeavor to realize the human ideal excludes no one. Hatred is increased by reciprocal
hatred, but is extinguished \opage{137}by love and generosity and itself passes over into
a love that becomes the greater the stronger the hatred has been.---Hope and fear are not
good in themselves either. They are bound up with uncertainty and sorrow and are signs of
uncertain and defective knowledge. One who is sure of his ground neither hopes nor
fears.---Pity is in itself bad and useless in a man who lives according to the dictate
of reason. Sorrow is bound up with pity, and the good that can spring from it springs
also from the dictate of reason. One who has seen that everything follows from the
necessity of the divine nature and happens according to the eternal laws and rules of
nature will pity no one. Such a mood easily leads to actions that afterward prove to be
wrong, and one will easily be deceived by false tears.---Humility and repentance are not
virtues either, but signs of imperfection. Humility is a sorrow that arises from man's
contemplating his own weakness, and repentance likewise rests only on man's seeing
himself as powerless. None of these feelings can therefore be the motive of perfect
action, which is to spring precisely from man's spiritual power. We come to the right
way sooner by reason and love than by repentance and contrition. Only superstitious
teachers, who understand better how to rail at vice than to preach virtue, who break
men's minds instead of strengthening them, and who therefore would rather lead men by
fear than by reason---only they can praise these feelings. The desire or drive of
self-preservation that is one with man's rational nature is directed straight toward the
good. \opage{138}Evil is avoided indirectly. When we follow our true nature, we need
food, to be sure, but not medicine. It is only the sick man who, from fear of death,
takes medicine. But the free man thinks of nothing less than of death; his wisdom has
life, not death, as its content.

Yet it is only when we speak of the man who is led by reason, of the spiritually free
man, that this holds. One who is moved to help others neither by reason nor by pity is
rightly called inhuman. And since men seldom live according to the dictate of reason,
humility and repentance, hope and fear do more good than harm, and if one must err, one
had better repent and fear too much than too little. If the irrational, spiritually
unfree men were all equally arrogant, they would be ashamed of nothing and fear nothing,
and nothing would be able to govern them. The prophets, who cared not for the advantage
of a few men but for the common advantage, therefore zealously enjoined humility,
repentance and reverence. And in reality those who are governed by these feelings can
far more easily than others be led to live according to the dictate of reason, i.e., to
be free and lead a blessed life.

Spinoza gives a series of means by which passion can be brought to become active
instead of passive \textit{(remedia affectuum)}. First and foremost, the task is to know
the nature and cause of the passion; for as soon as it is known, it ceases to be a
passive passion. All passive passions arise, after all, from our being parts of nature
and therefore having only inadequate ideas; the passive therefore falls away with
adequate knowledge. It is a common experience that sorrow \opage{139}is softened when one
sees that the misfortune was necessary. Yet it is one and the same passion that first
appears as passive and afterward as active. The ambitious and proud man endeavors just as
much as the one who is led by reason that all should live according to his thoughts; but
what the former desires from inadequate insight, the latter strives for with full and true
knowledge.---Another important means is to separate the emotion or passion from the
thought of the external cause and connect it with other thoughts. Love and hatred arose,
after all, from joy and sorrow in connection with the thought of the things that caused
these feelings; when this connection is dissolved again, the passion must thus also take
on another character. We must use the intervals in which the passion is not raging to
bring about and practice such a connection among the representations as reason requires,
thus in particular to lead them all back to the highest thought, the idea of the Deity.
Since thought embraces the common laws that are always at work in everything, the
feelings that spring from the thought of the eternal will always have occasion to arise
and will gradually gain the upper hand over the other passions. Our finite nature
moreover entails that we must seek nourishment and strength in something outside us;
therefore we cannot be freed from love as from the other passions. It is altogether an
excellent thing that all passions that are good are of such a nature that without them
we cannot be or exist. The more causes and points of attachment a passion has, the
greater power it has. And this holds precisely of the passion that springs from
\opage{140}the thought of the eternal and divine, which lies at the ground of everything.
A clear and distinct knowledge brings joy with it, because our power of action reaches
its fulfillment in it. And since this joy, when knowledge goes back to the first cause,
has God as its cause, there arises the love of God, which is nourished by everything we
know, and must therefore be the mightiest passion of all. Even if love of the
unchangeable and eternal cannot entirely drive out the passions, it will nevertheless
restrict them to the smallest part of the soul.

Intellectual love---so Spinoza calls (after Descartes and older Jewish philosophers of
religion) the joy in the eternal and infinite that springs from true knowledge---which is
the transfigured passion, shows us man's essence from an entirely different side from the
one from which we have so far considered it: no longer as an individual being, but as a
being that, through its unity with the infinite fundamental being, acquires existence in
itself and is eternally free. But this intellectual love is not merely a goal that man
approaches and strives to reach. When the concept of God is thought according to the
strict consistency of the system, the matter stands otherwise.---As absolute being God
cannot pass over to greater or lesser perfection, and so cannot feel love or hatred
toward anyone. One who would wish that God should love him in return would wish that God
were not God. That God does not love men must not, however, be understood as though they
were indifferent to him. Man and everything that is, is in God in such a way, and God
consists of all beings in such a way, that he cannot \opage{141}feel love for anything
or anyone, since nothing is given outside his own essence. Man's love of God is a part of
the infinite love with which God embraces his own essence. In it man regards himself as a
member in, and form of, God's infinite essence. Since this cannot step forth in any
definite form, ``as an object among many others,'' intellectual love gets no object, does
not really become love \emph{for} anything. Therefore there can be no question of
requited love either; the subject and object of the love are not different. It is, one
may say, an immanent love; it has the same character as the acting of the fundamental
being in general: God can only be an immanent (reflexive) cause, not act on anything
outside himself.---This doctrine is not only consistent, but also contains the great
truth that every ideal striving has its goal not outside itself but in itself.
Blessedness, says Spinoza, is not the reward of virtue, but is virtue itself. For the
mythological imagination the matter stands thus, that the labor is one thing, the reward
another. But on Spinoza's view what holds is what stands in Heiberg's % the print: „Hejbergs“
poem:
\begin{quote}
In seeking him, you have already found him;\\
In asking, the answer is already won.\footnote{``This being so, we may rightly declare it
a great absurdity, what many who are otherwise regarded as great theologians say: that
if no eternal life followed from the love of God, one would do best to seek one's own
self, as though one could find in it something better than in God. This would be just as
foolish as if a fish, which after all cannot live except in the water, were to say: If
no eternal life follows for me after this life in the water, then I will out of the
water, up onto land.'' \emph{Spinoza}: ``Short Treatise'' II, 26. (Schaarschmidt's
German translation).}
\end{quote}

\opage{142}Intellectual love must, like true knowledge, be eternal, and has neither
beginning nor end. In it man is, after all, one with the fundamental being and thus
eternal as this is, and it is only owing to an inadequate knowledge that he is conceived
as being in time. It does no harm to this eternity that we cannot remember having existed
before our existence in time. Memory presupposes relations of time, but the eternal has
nothing whatever to do with time. True knowledge, which sees everything under the aspect
of eternity, is raised above sense representation and memory, which belong only to time.
Most people, to be sure, are not able to reach this knowledge. Just as by God they imagine
a single object, so by the eternity of the mind they understand a simple continuation and
persistence of their existence in time. But the eternity of the mind signifies nothing
other than its nature as a necessary element in the fundamental being (an eternal mode of
thought, \textit{æternus cogitandi modus}).

In this mystical doctrine of intellectual love Spinoza's philosophy bends back toward its
main principle. After the experimental undercurrent has been the stronger in the course
of the psychological and ethical investigations, religious speculation now again comes to
the fore. The opposition between the two currents makes itself sharply felt here too,
since Spinoza cuts off the possibility of reaching the ideal standpoint through finite
development. Just as \opage{143}he denied, in an earlier connection, that the individual
manifoldness of the individual beings could be derived from the fundamental being, so he
now denies that the individual being can reach knowledge of the eternal through a
development in time. The whole process of liberation we have seen man go through thus
rests on an illusion,\footnote{This is one of the points that \emph{Brøchner} has
illuminated with great clarity in his monograph ``Benedikt Spinoza'' (Kbhvn. 1857).
Chaps.~6 and 8.} the same illusion that is peculiar to the lowest form of knowledge:
namely, that the existence given in experience, unfolding itself in time, has reality.
In reality the fundamental being is the one thing existing, and man, as an element in its
nature, is himself eternal---and so can be subject to no passive passions and go through
no development.

Spinoza therefore abandons here too his earlier psychological determinations. He
conceived joy, after all, as transition to greater perfection, and we stressed what is
true in this determination. But it is no longer enough for him. In the absolute and
eternal there can be no question of transition: therefore he says that in place of joy
\textit{(lætitia)} there now appears blessedness \textit{(beatitudo)}, which consists in
the mind's being in possession of perfection itself.---It is no wonder that Spinoza here
has difficulty in making himself understood, and we in following him. A perfect
conclusion and fulfillment conflicts with psychology. Psychology, just as much as physics
and physiology, knows only relative forms of state. Spinoza's contemporary Thomas Hobbes
saw this clearly and expressed it thus (\textit{De homine} \opage{144}XI, 15): ``If there
is to be a final end, then, when it is reached, there is nothing to strive for and
desire; then there is nothing either that stands before man as a good; indeed he can no
longer even have any feeling, for every feeling is bound up with desire or aversion, and
not to feel is the same as not to live. The greatest good consists, on the contrary, in
going forward as unhindered as possible to ever higher ends. Life is a constant motion,
and when this motion cannot go in a straight line, it goes in a circle.''---One will
easily see that we here again stand at what has been mentioned several times, that all
deductive speculation really makes ethics impossible, since the eternal fundamental being
is the one thing that is. The law of relativity, on the other hand, shows us that our
life, like everything we know, is conditioned, is a link in a great development. Labor
and striving are thus not illusions, but the expression of a universally valid law. An
existence of pure being without becoming is incomprehensible to us.

Plato too, with whom Spinoza has so much in common altogether, stopped at an
insurmountable opposition between becoming and being. The ideal world in its eternity
was for him true being, and the existence given in experience, where becoming,
development, striving, and thus also ethical labor, belong, became in the end an
inexplicable appearance. That is why the highest thing for Plato, as for Spinoza, was
theory, the speculative intuition of the eternal. And yet Plato, like Spinoza, is
distinguished by his deep psychological understanding of the nature of drive.

Spinoza is reminiscent here also of the religious mysticism that flourished at the end
of the seventeenth century, \opage{145}especially within the Cartesian school (Fénélon,
Malebranche). He is in his way a champion of ``disinterested love,'' and here stands side
by side with the noble Archbishop of Cambray. The religious mystics too seek an immediate
union with the eternal and yearn for a standpoint where time and finitude vanish. What
separates Spinoza from them (apart from the opposition between his philosophy and their
theology) is that they move essentially in the domain of feeling, he in that of thought.
Their disinterested love of the eternal is sentimental,\footnote{This word used in its
proper, not in its pejorative, sense.} his is intellectual. He has, as Jakobi aptly said
of him, his heaven in the understanding. It belongs to the one-sidedness of his
psychology that he sees the essence of the mind exhausted in thought. Will is only assent
to the truth known, feeling only a confused thought. In his doctrine of the passions he
shows how these by a metamorphosis pass over into becoming thoughts. Pure intellectual
activity is for him the proper essence of the mind. It was only toward the end of the
eighteenth century that psychology became properly clear about the peculiarity and
independence of feeling as a mental phenomenon beside knowledge and will. It was surely
the eighteenth century's bringing forward of individuality, its significance and its
right, that opened people's eyes here. Feeling has a more individual character than the
other expressions of mental life. We can to a certain degree think and act like others,
but feel only as ourselves.
\begin{quote}
\opage{146}Allen gehört was du denkst, dein eigen ist nur,\\
\hspace*{4em}was du fühlest!\\
\hspace*{2em}(Schiller).
\end{quote}
Spinoza could, to be sure, immerse himself in the contemplation of the individual
emotions and their laws, but he nevertheless always sought beyond this individual
element to the deep unity and the great whole into which, for him, everything individual
flows out as it reaches its fulfillment. His personal character seems moreover to have
been decidedly contemplative; to be moved by hope and fear, by the many changing moods
under which most men's personal lives develop and must develop, was foreign to him,
because he had once and for all dedicated himself to knowledge and gathered up his
personal development into the development of thought. Feeling appears in him as the mild
glow, the quiet joy, that knowledge won spreads around it; the quiet shores where he
dwelt were not visited by life's more violent breakers.

\section*{V.}
\addcontentsline{toc}{section}{5. Political Theory. Philosophy of Religion}
\opage{147}However solitary and withdrawn a position Spinoza occupies in his thinking and in
his life, all his writings nevertheless bear the stamp of his having felt himself to be
``a man among men.'' The task he had set himself, and whose solution he regarded as the
highest good, the knowledge of our mind's connection with the whole of nature, could, by
his conviction, be reached not by the individual for himself alone, but only in community
with others. The solution of the task therefore presupposes partly a reformation of
knowledge, partly the formation of a society in which as many as possible are educated
and developed toward the stated goal. In his ethical doctrine we see how the strength of
mind that for him signifies the highest development of character appears not only as
courage (in the individual's own affairs), but also as generosity (in relation to
others).---The ethical character presupposes a development that is not present in all,
and when the question therefore becomes one of the foundation of human society, we must go
further back, back to the same point from which ethical development began: man as a mere
natural being, subject to the dominion of passive passions. The essence and origin of the
state cannot be understood unless one understands the nature of passion. Spinoza therefore
stresses here, \opage{148}as in the Ethics, the necessity of studying passion like any
other natural phenomenon. Statesmen, he says, have written far better about politics than
philosophers, because the former have followed experience, while the latter have
censured the passions and constructed a state for pure beings of reason, who do not
exist, or who in any case need no laws. Experience presents so many different forms of
state that we can hardly devise anything altogether new that the necessary demands of
practical life would not have brought to light in some place or other.---Since men are
led essentially by the passions, against which the dictates of reason and religion avail
but little, the state must be so ordered that nothing is left entirely to the fidelity
and morality of individuals, but that those who rule are prevented by their own interest
from acting dishonestly, whatever disposition might then lie at the ground of their
acting. Freedom and strength of mind are a private virtue; and in the state, where
security for all is what counts first and foremost, one dare not count on it.

Every natural being is a finite form of the infinite fundamental being, and the power by
which it exists, acts and asserts itself in its existence is a part of the infinite power
of the divine being. Therefore every natural being originally has---whether it is led by
reason or passion---just as much right as it has power. Natural right is one with natural
power, or with the law of nature that originally determines the individual's action. It
is thus by a natural right that the big fish devour the small. For each individual, what
\opage{149}he regards as good and evil is determined by his interest. He keeps a given
promise, for example, only so far as it agrees with his own advantage. If a robber has
forced me to promise to hand over my fortune to him, I do not hesitate to break this
promise when I can do so without danger. In the state of nature nothing is forbidden
except what no one thinks of doing. Now since the passions---especially anger, hatred and
envy---bring men into conflict with one another, they are by nature enemies. The one must
always fear the other. Under these conditions there exists no property---each possesses
only what he has the power to defend---and no security. But for that reason natural right
does not really exist at all as long as it is determined by each one's power. Only by
united forces can men sustain their lives and cultivate their minds. One can therefore
say that man is a social being: his natural conditions lead him to the founding of a
social life. The more there are who unite their forces, the greater the power, security
and right they possess.

When such a society is founded, individuals in part hand over their right to the power of
the community \textit{(imperium)}. What the power of the community declares right and
wrong, individuals must do or refrain from. Each individual now has only the right the
power of the community grants him. But how far does the right of society extend? Only so
far as its actual power extends. Society rests on the principle that a passion can be
checked only by another, stronger passion, and that man must therefore be kept from
inflicting harm on others by the fear of suffering greater harm himself. What society
\opage{150}cannot compel man to do, it has no right to command. The power of the
community must therefore, according to the law of self-preservation, which is one with
the dictate of reason, command nothing that would put an end to its own dominion. Just as
the individual is strongest when he follows reason, so only that society can last which
is determined by reason. To say that the power of the community may do what it pleases is
the same as to say that a man may take leave of his senses. It came into being to help
individuals to their natural right and must therefore seek what is most useful for all.
What conflicts with human nature and outrages human feeling the power of the state cannot
command without aiming at its own downfall. It then itself becomes an evil for man of the
same kind as those he suffered under in the state of nature; individuals will seek means
of securing themselves against it---and society will thus be dissolved. In particular the
state would overstep its right if it sought to intervene in the domain of free
conviction, where threats and promises avail nothing. No one can sign away his judgment
or prevent another from using his. In spite of all prohibitions and all compulsion men
will love and hate as agrees with their nature.

There is thus no absolute opposition between life in the state of nature and life in
society. On this point Spinoza differs from his contemporary Thomas Hobbes, who taught
that the drive of self-preservation, which leads men beyond the state of nature, this war
of all against all, commands complete submission \opage{151}to the will of the power of
the state as the only means of securing peace. Spinoza grants that every member of society
must, admittedly, give up the right to live as he himself sees fit; in place of his own
decision as to what is good and evil, he must now follow that of the state. Nor can he
have the right to interpret the laws in his own way, for then he would be his own judge,
and it was precisely in order to have a common measure of right and wrong that the state
was founded. But the resignation that must be shown here does not conflict with reason.
Even if in some cases I must do what I myself do not consider right, I nevertheless
choose the lesser of two evils: it would be a far greater evil if social life were
dissolved and the state of nature began again!

The origin of the state is thus, according to Spinoza, to be sought in human egoism. To
avoid the violent clash of interests one hands over to society as a whole the care of the
common welfare. But Spinoza does not see the task of the state exhausted in mere security.
Its purpose is not merely to rule and by fear to keep one from harming another. Of a state
whose citizens are kept only by fear from taking up arms against one another, one can
rather say that it is without war than that it has peace. Peace is not merely the absence
of war; true peace proceeds from strength of mind and from adherence to the common laws
of the state. It is precisely the task of the state to banish fear entirely. The state
does not aim at turning men into machines or irrational animals, but at making a
spiritual life possible. The end of the state is \opage{152}therefore freedom. That this
end is not reached must be ascribed not so much to the wickedness of the citizens as to
the badness of the constitution. Where discord, rebellion and contempt for the laws
prevail, it must be because the institutions are no good. For men are not born citizens,
but must develop into becoming such.

Spinoza regards democracy as both the most natural and the most perfect constitution.
It is most natural that men should not at once hand over the whole of power to one person
or a few, but to the whole assembled multitude that agrees to create a society. Each
individual then, since he himself belongs to the multitude, gets a share in future
deliberations on the use of the power of the state. Thereby hope rather than fear also
becomes the foundation of society. The most perfect constitution is democracy, because it
comes nearest to the freedom that belongs to each by nature.

Spinoza is the first modern political writer to have declared himself for a democratic
republic. The republican writers before him in the sixteenth and seventeenth centuries
were more or less aristocrats. He is in this respect Rousseau's predecessor. He grants,
however, that if the patricians let themselves be led not by their passions but by
regard for the public good, an aristocratic constitution would be best of all; but such a
presupposition he regards as contrary to experience. Unfortunately he died just as he had
begun, in his Political Treatise, to set forth the democratic constitution, so that we
cannot see what special institutions he considered \opage{153}necessary under
it.\footnote{The last paragraph in the unfinished treatise, and thus probably the last
thing Spinoza wrote before his death, deals with the question how far women in a
democratic republic ought to have equal rights with men. If women were by nature equal to
men, history---Spinoza thinks---would have offered examples of such equality of rights.
Moreover, women could not without great inconvenience have a share in political power:
the erotic drive and the passions it brings with it (jealousy, etc.) would then interfere
disruptively.---The first of these objections Stuart Mill has thoroughly refuted in his
well-known work; the second, on the other hand, he does not mention at all.} On the other
hand he sets forth in detail the institutions of aristocracy and monarchy. He seeks to
show how both rulers and ruled---with or against their will---can be brought to act
according to the dictate of reason. An absolute monarchy is impossible, since the king
needs counselors, on whom he will always be more or less dependent, and the constitution
will thus in reality become aristocratic. The kingship Spinoza describes corresponds to
what we call a constitutional monarchy \textit{(imperium monarchicum, qvod a libera
multitudine instituitur)}. The king is even bound to choose among the proposals that have
received the most votes in the great council, in which the whole people is represented.
The judicial power is exercised by a separate body. An absolute autocracy would, if it
were possible, perhaps lead to greater security, and state secrets could be better kept
under such a constitution; but these advantages would be bought by giving up far more
essential goods. When the autocrat has nothing to take into account but his own pleasure,
one has \opage{154}no guarantee that he will strive for the general good, whereas he will
naturally be an enemy of all eminent men who might seem to dispute his rank. After Spinoza
has discussed the forms and institutions by which the royal power ought to be restricted,
he remarks that his account will perhaps provoke laughter in those who are accustomed to
imagining all vices on the side of the multitude and all virtues among the great. But
human nature is one and the same in all, and precisely those who rule are by their
position exposed to many faults.

The starting point of Spinoza's political theory\footnote{Spinoza's political theory
stands---not merely for its own time---alone in clarity, consistency and
liberal-mindedness. But no doubt owing to his isolated position he had in this domain, as
well as in psychology, to stand in the shadow of Hobbes, although the latter's paradoxical
% the print: „parodoxe“
and shortsighted pessimism, at least at first glance, seemed to make him an opponent of
progress.} is the natural right of the individual, which is not to be given up but
precisely to be secured and developed in social life. He is here a predecessor of the
eighteenth century and its struggle for the rights of the individual. Kant's proposition
of freedom---in the sense of independence of the coercive arbitrariness of anyone
else---as the one original human right, and the development that Wilhelm von Humboldt and
Stuart Mill have given this proposition, is in principle stated in Spinoza. And what is
surprising in a speculative and abstract philosopher: the working out of his doctrine has
everywhere a practical and real character and is built on sharp psychological insight.

But although Spinoza recognizes no complete \opage{155}opposition between the state of
nature and social life, he nevertheless regards them as different and has the transition
from the one to the other take place by a compact. It was a notion that was often taken
as a basis in the seventeenth and eighteenth centuries, and that sprang essentially from
defective insight into the laws of human (and all) development. Just as in geology the
present state of the earth used to be thought of as the result of a series of great
revolutions and quite different from earlier states, so the transition between the
different states of human life was thought of as an arbitrary act. Such abrupt
transitions a more developed psychological and historical insight does not admit. The
law of continuity reigns here too. In spite of all drastic events and catastrophes, social
life has developed essentially under the same laws on which it now rests, by a quiet,
often imperceptible growth.

---There is a close connection between Spinoza's political theory and his doctrine of
religion. In his ``Theologico-Political Treatise'' he sought, on the basis of his doctrine
of the relation between faith and knowledge, theology and philosophy, to show the
possibility and necessity of religious freedom. And in the political theory we have seen
that there is a domain into which the power of the community cannot penetrate, namely
the individual's free conviction. His external actions the state can determine, but
religion belongs essentially to the inner life and consists in the knowledge of the
divine and in love of one's neighbor. The external forms of worship are indifferent, and
the religious man will \opage{156}show his love precisely by preserving peace in society
and not stirring up discord by sectarian conduct.

The religiousness that corresponds to Spinoza's standpoint he himself defines thus:
``Everything we desire and do, insofar as we are led by the idea of God, I refer to
religion. The desire to do good that arises from our living according to the dictate of
reason I call piety.'' But since this philosophical concept of religion conflicts with the
theological conception of it, he seeks to determine more closely the relation between
theology and philosophy. The theological concept of religion rests on the idea of a
revelation. Here, of course, there is no talk of revelation in the wider sense of the
word, in which all knowledge we gain of nature can be called revelation, even if we use
only ``the natural light.'' Revelation, in the strict sense of the word, is different from
rational knowledge: it does not come, as this does, through clear and distinct concepts,
but takes place through words and signs. Not all can receive such words and signs and
communicate them to others; for that, individuals with a special gift are required, who
are called prophets. But this special gift cannot consist in eminent understanding and
knowledge, since it is, after all, not pure thoughts but images and symbols that are to
be apprehended and communicated. What counts is rather, above all, a strong and lively
imagination. There is a natural opposition between individuals in whom imagination is
most developed and those in whom understanding is; the one of these faculties is
commonly developed at the expense of the other. That is why it is a great
misunderstanding to seek knowledge of \opage{157}the nature of the world and of the mind
in the writings of the prophets. Such knowledge must be gained in an entirely different
way. How, indeed, should a warrior like Joshua understand astronomy, or Noah geography?
Why, Adam did not even know that God was omniscient, since he hid among the trees of the
garden; nor did Abraham know it, since he asked God to count how many righteous men there
were in Sodom. The prophets had, moreover, to accommodate themselves to those to whom they
spoke. Only of Christ does it hold that he was not governed by imagination, but received
revelation by immediate and true knowledge. It is, as we have heard Spinoza say in one of
his letters (above, p.~49), not necessary to know the historical Christ (``Christ after
the flesh''), but ``the Son of God'' in the sense of the eternal wisdom that reveals itself
in all things, most in the human mind, and most of all in Christ, who taught men to know
the universal law of reason. In the Sermon on the Mount Spinoza sees, like modern
rationalism, the sum of true Christianity. The law of Moses was a national and political
institution that lapsed with the downfall of the Jewish kingdom. The Jews have nothing to
boast of over against other peoples. Nations other than the Jews too have had prophets;
it is not something peculiar to the Jewish people.

Revelation acquires a different character in the different prophets, since their
personality, bodily constitution and previous opinions necessarily come to influence
their preaching. The courtier Isaiah speaks differently from the shepherd Amos. One cannot
therefore without further ado use the writings of the one to \opage{158}interpret those
of the other. On the basis of sure knowledge of the language and acquaintance with the
author's life, the historical circumstances and the authenticity of the text, one must
find the actual meaning of each particular passage. Whether this meaning is true in
itself is an entirely different question. We cannot without further ado start from the
assumption that Scripture agrees with reason; that would lead to the consequence that the
final ecclesiastical authority would have to be sought among the philosophers, but this
authority would hardly meet with great respect. What Scripture teaches, and how far what
it teaches is true, are thus two entirely different questions.

Spinoza, who is the first to have spoken for a purely historical interpretation of
Scripture, is also the first to have applied a historical criticism to the biblical
books. He maintains that the authenticity of each particular writing must be proved for
itself, and comes to the conclusion that the Old Testament writings in their present form
cannot have the authors whose names they bear.

The difficulties at which the scientific interpretation of Scripture comes to a halt
need not, however, according to Spinoza, trouble those who seek edification in
Scripture. What is edifying in Scripture \textit{(pietatis documenta)} can always be
easily understood. And the real aim of the prophets was to proclaim a doctrine that could
lead to piety and to obedience to God. That is why they present God as a lawgiver and
prince whom men must obey. Theology differs from philosophy precisely in teaching that men
are saved by obedience, whereas philosophy maintains that true virtue springs
\opage{159}from insight into the connection of nature. Theology seeks piety, philosophy
seeks truth. Truth is in itself neither pious nor impious. Only insofar can one say that
there is piety or impiety in a man's opinions as he takes occasion from them for
obedience or disobedience. One who becomes arrogant by believing what is true has an
impious faith; one who becomes obedient by believing what is false, a pious one. The
faith that is to be common to all men can therefore include nothing about which honest
men can disagree. When it is demanded that one believe in God, then by God must be
understood an ideal pattern for human life \textit{(veræ vitæ exemplar)}, but it is
indifferent how each thinks of God---as fire or light, spirit, thought, in a pantheistic
or theological way, etc. Here the differences of intellectual gift and development come
to play a role. All can obey and lead a life in justice and love, but not all can acquire
a knowledge grounded in reason.

By referring positive religion to the practical domain Spinoza seeks to secure the
freedom of science. He leaves piety to theology, but reserves truth for philosophy, a
division with which the theologians could not be entirely satisfied. One is tempted to
regard such a division as ironically meant, and yet no irony can be traced in Spinoza's
account. It is clear, however, that it was not his real meaning thereby to set faith and
knowledge side by side as equals. Truth stands for him above piety and regulates it, not
the reverse. \opage{160}He not only remarks that, in order to bow in the obedience of
faith, one must not know that what one believes is untrue. He even says that it is
impossible for philosophy to see that men are saved by obedience, because obedience takes
account only of the will of the one who commands, not of the necessity and truth of the
matter. Further, in his ``Theologico-Political Treatise'' he criticizes miracle from a
purely philosophical point of view: since nature is one with God, one cannot, as in the
belief in miracles, appeal from the power of nature to the power of God, and it is no use
saying that miracle is above, not against, nature, since miracle, after all, happens in
the midst of nature and so interrupts its order. And finally he states outright that the
prophets, only because of their defective knowledge, presented God as lawgiver and prince
instead of as the eternal and infinite fundamental being. They did so in order to
accommodate themselves to the comprehension of the multitude. Most people do not do good
for its own sake and from love of it, but because they fear an almighty lawgiver and are
moved by the reward and punishment he has laid down. The multitude likewise needs
ceremonies and historical narratives and examples that can make an impression on its
imagination. Rational arguments make no impression on it. One who uses his reason, on the
other hand, does not need historical religion. The law is given only for the unfree, not
for those who have ``the fruits of the Spirit'' and possess true knowledge.

There thus remains a definite opposition between those who know and those who believe.
They are two different views of life that stand over against each other. % the print: „hinanden-“ (sense: hinanden.)
And \opage{161}the opposition is rather aptly stated by Spinoza's antithesis:
obedience---truth. What separates the two tendencies is, namely, essentially the different
conception of the significance of authority: positive faith holds fast that there are
absolute authorities; rational faith knows only relative authorities. Spinoza has not,
however, brought out the historical relation between the two tendencies, which, viewed
abstractly, are altogether irreconcilable opposites. For historical reality shows us
authorities as educating powers in the human world, and there takes place a constant
giving and receiving between the two poles; in the course of history they constantly
determine each other. If one compares the series of religions, and of authorities in
general, in history, one will find that the higher ones always have a more rational
character. On the other hand authority is---as Spinoza himself has shown in his political
theory---a necessary support for developing freedom. Just as the universe, in the opinion
of some astronomers, shows us groups of stars at very different stages of development
from chaos to system, so we see at all times in the human world many different stages of
development from authority to freedom. And Spinoza has shown by his own example that the
transition from the one standpoint to the other does not entail the abandonment of true
religiousness. The religious need can find its satisfaction in more forms than the
dogmatism of positive religion admits. Spinoza's innermost line of thought was of a
religious character, if by religiousness one understands \opage{162}devotion to the
higher connection of which the individual is a member. A feeling of this connection is the
lifeblood of his system.---

Systematic a mind as Spinoza was, men must nevertheless, by his conviction, seek
satisfaction for their religious need each in his own way. The mental nature of men
\textit{(ingenium)} is so different that they cannot find rest in the same opinions. The
power of the community will therefore only do harm and work against itself if it seeks to
violate freedom of thought and speech. Each is master of his conviction and must have the
right to express it, provided he does not do so in a way that can stir up hatred and
rebellion in the state. Piety and impiety concern, after all, as we have heard, not
opinions in themselves, but the actions that spring from opinions. If the state were to
intervene in the domain of free conviction, it would meet resistance precisely in the
noblest and best men. Those who think only of amassing money and tending their bellies
would naturally not care about such an intervention. But what greater misfortune can be
conceived for a state than that honorable men be sent into exile as criminals because
they hold peculiar views and do not know how to dissemble? What is more terrible than that
men be led to death, not on account of any crime or misdeed, but on account of their free
cast of mind? Freedom may, to be sure, entail inconveniences; but there are so many
inconveniences one must put up with in order to obtain the greater goods.---Who, in
\opage{163}any case, should have the right to intervene in the individual's conviction?
One cannot appeal here to any divine authority. God can only falsely be represented as
prince and lawgiver; the divine law is one with the necessity of eternal nature. History,
particularly the history of the Jews, shows plainly enough that it leads to dubious
consequences when the clergy obtains political power. The state must therefore have the
right to order external worship as agrees with its welfare. But the best arrangement is
nevertheless that the various religious parties obtain freedom to exist side by side.
Spinoza cites Amsterdam as an example of such an arrangement being carried through
without unfortunate consequences. Nothing is safer for the state than that piety and
religiousness be placed only in love and justice, and that the right of the authorities,
in sacred as in secular matters, be restricted to actions alone, while each is otherwise
allowed to think what he will and to say what he thinks.

Spinoza could no doubt have gone still further and restricted the right of the power of
the community to intervene even in the domain of actions still more. In this direction
Kant and Stuart Mill have continued his line of thought. But it is already a great thing
that, in the midst of the dogmatism and rule of authority of the seventeenth century, he
stated the principle of freedom of thought and speech. He shows here how little right
people have so often had to deny him an eye for the significance of personality. As his
religious and philosophical \opage{164}conviction was in harmony with his clear, quiet
and yet energetic personality, so he also shows a deep understanding of the significance
of personal conviction for the full and rich development of human life.

% ===== END MATTER (pp. 165--166) =====
\chapter*{List of the Most Important Philosophical Terms and the Places in the Book Where Their Meaning Can Be Seen}
\addcontentsline{toc}{chapter}{List of Philosophical Terms}
\markboth{List of Terms}{List of Terms}
\opage{165}\relax[The page numbers are those of the original edition, given in the margin of
this translation. The Danish equivalents in parentheses are the author's.]

\begin{description}\setlength{\itemsep}{0pt}\setlength{\parskip}{0pt}
\item[Absolute] (cf.\ relative). 69--71. 97 f.
\item[Anthropomorphism] (conception of the divine under human forms). 118. % the print: „mennelige“
\item[Atheism.] 54--58.
\item[Attribute] (\textit{Grundform}, fundamental form). 70 f. 73 ff. 80 f. 89 f.
\item[Deduction, deductive] (cf.\ induction). 6. 14 f. 44--46. 68. 95.
\item[Deism, deist.] 54--59.
\item[Determinism] (\textit{Nødvendighedslære}, doctrine of necessity), deterministic. 26. 49. 51--54. 120 ff.
\item[Discursive] (cf.\ intuitive). 110 ff.
\item[Dogmatism, dogmatic.] 4--6. 118.
\item[Dualism] (\textit{Tohedslære}, doctrine of twoness), dualistic (cf.\ monism). 7--8. 59. 94. 107.
\item[Empiricism, empirical.] See induction.
\item[Immanence] (\textit{Iboen}, indwelling), immanent. 48. 83. 141.
\item[Induction] (\textit{Erfaringsmetode}, method of experience), inductive. 6. 12--15. 44--46. 76--79. 106 f. 110.
\item[Intuition] (\textit{anskuende Tænken}, intuiting thought), intuitive. 110--115.
\item[Critical philosophy] (cf.\ dogmatism). 4--6.
\item[Mode] (\textit{Enkeltvæsen}, individual being). 73. 84 ff.
\item[Monad.] 67. 71. 93 f.
\item[Monism] (\textit{Enhedslære}, doctrine of unity), monistic (cf.\ dualism). 8. 59.
\item[Optimism.] 31. 134 ff.
\item[Pantheism, pantheist.] 8--10. 53. 57--59. 82. 140--141
\item[Relation] (\textit{Forhold}), relativity (\textit{Betingethed}, conditionedness), relative. 59. 70 f. 81. 88. 118 f. 130. 143 f.
\item[Speculation, speculative.] 2. 13. 16. 86 f. 97 f.
\item[Substance] (\textit{Grundvæsen}, fundamental being). 6--9. 17. 26. 67. 70 f. 72 ff. 76 ff.
\item[Theism.] 53, cf.\ 57.
\end{description}

\chapter*{Translation of the French Verse on Page 56; Note to Page 127; Errata}
\addcontentsline{toc}{chapter}{Translation of the French Verse; Note to Page 127; Errata}
\markboth{Note to Page 127}{Note to Page 127}
\opage{166}\section*{Translation of the French Verse on Page 56}
% the Danish verse translation, signed „H. T.“, here put into English
\begin{quote}
Then came a little, sallow and long-nosed Jew,\\
Poor but content, thoughtful, turned inward,\\
A fine, subtle mind, less read than acknowledged.\\
Wrapped in the garments of Descartes, his master,\\
Slowly he steps before the Most High\\
And whispers softly to him: ``Pardon me, I would\\
Confide to you, unter uns, I do not believe you exist.''\\
\hspace*{\fill}H. T.
\end{quote}

\section*{Note to Page 127}
A note ought to have been added here to the effect that Spinoza is surely not right in
applying the laws of the association of ideas directly to the feelings. It deserves in any
case closer psychological investigation whether the one feeling can directly call forth
the other, or whether the way from the one feeling to the other does not rather go through
the thoughts and representations attached to the feeling and the associations of ideas
that these call forth. Any thing can thus indirectly become a cause of joy, sorrow and
desire, not because these feelings pass directly over into one another, but because
various representations are attached to the thing, of which the one awakens joy, the
other sorrow, etc. A memory will, for example, awaken joy when the representation of the
earlier happiness is predominant, but sorrow when that representation is succeeded by the
representation that this happiness is now over.

\section*{Errata}
% corrected in this translation at the sites (pp. 8, 104, 132); the second is printed as "S. 10"
\begin{center}
P.~8, l.~5 from bottom: for \textit{for} read \textit{fra}.\\
P.~10 [i.e.\ 104], l.~3 from bottom: for \textit{hallucition} read \textit{hallucination}.\\
P.~132, l.~11 from bottom: for \textit{gneerositas} read \textit{generositas}.
\end{center}

\end{document}
